% Leçon 19 — Expression écrite : offre de services aux usagers des administrations (collectivités locales) — Chinois Tle A — Cours signé OnBuch+
% Programme officiel MINESEC. Compiler avec : tectonic ZH19-ecrit-offre-de-services-aux-usagers-des-administra.tex
\newcommand{\DOCMATIERE}{Chinois}
\newcommand{\DOCNIVEAU}{Tle A}
\newcommand{\DOCMODULE}{Module 3 : Citoyenneté et ouverture au monde}
\newcommand{\DOCLECON}{ZH19}
\newcommand{\DOCTITRE}{Leçon 19 — Expression écrite : offre de services aux usagers des administrations (collectivités locales)}
% OnBuch+ — préambule « cours signés » ENRICHI (compilé avec tectonic / XeLaTeX)
% Système visuel « Pop » d'OnBuch+ : fond crème (jamais blanc), bordures encre
% épaisses, ombres « dures » décalées, titres Archivo Black en capitales.
% Version MATHÉMATIQUES : graphes (pgfplots/tikz), boîtes pédagogiques
% (définition, propriété, méthode, attention, le savais-tu, exercice résolu,
% exercice, TP, bilan), page de garde et signature OnBuch+ ancrée en bas.
%
% Macros attendues dans main.tex AVANT \input{preamble} :
%   \DOCMATIERE  \DOCNIVEAU  \DOCMODULE  \DOCLECON (numéro)  \DOCTITRE
\documentclass[11pt]{article}

\usepackage{fontspec}
\usepackage[french]{babel} % français = langue principale ; le chinois via \zh{..} / textebox (xeCJK)
\usepackage{xeCJK}
\usepackage{pifont} % \ding{51} \ding{55} utilisés par les modèles
\usepackage[a4paper,margin=2cm,top=3.4cm,bottom=2.8cm,headheight=1.4cm,footskip=1.3cm]{geometry}
\usepackage{amsmath,amssymb,amsfonts}
\usepackage{mathtools} % \xmapsto, \coloneqq... souvent utilisés par les modèles
\usepackage{cancel} % simplifications barrées (bilans de forces)
% \abs{z} (module, valeur absolue) et \norm{u} (norme), utilisés par les modèles
\providecommand{\abs}{}\let\abs\relax\DeclarePairedDelimiter{\abs}{\lvert}{\rvert}
\providecommand{\norm}{}\let\norm\relax\DeclarePairedDelimiter{\norm}{\lVert}{\rVert}
\usepackage[table]{xcolor} % [table] : fournit aussi \rowcolors (alternance de lignes)
\usepackage{pagecolor}
\usepackage{tikz}
\usetikzlibrary{arrows.meta,positioning,calc,shapes.geometric,shapes.misc,shapes.arrows,decorations.pathmorphing,decorations.pathreplacing,decorations.markings,patterns,shadows,mindmap,trees,fit,backgrounds,matrix,intersections,angles,quotes}
\usepackage{pgfplots}
\usepackage[version=4]{mhchem} % équations nucléaires \ce{^{235}_{92}U}
\usepackage[european,siunitx]{circuitikz} % schémas électriques
\pgfplotsset{compat=1.18}
\usepgfplotslibrary{fillbetween,groupplots} % aires entre courbes (calcul intégral)
\usepackage{enumitem}
\usepackage{fancyhdr}
% [most] charge toutes les bibliothèques tcolorbox (listings, minted, raster,
% documentation...) et gonfle inutilement la mémoire TeX sur un document long
% (~300 boîtes) : on ne charge que ce qui est réellement utilisé.
\usepackage[skins,breakable]{tcolorbox}
\usepackage{graphicx}
\usepackage{colortbl}
\usepackage{booktabs}
\usepackage{tabularx}
\usepackage{diagbox} % case d'en-tête diagonale des tableaux à double entrée
% Colonnes extensibles centrée / gauche / droite (C, L, R), souvent utilisées par les modèles
\newcolumntype{C}{>{\centering\arraybackslash}X}
\newcolumntype{L}{>{\raggedright\arraybackslash}X}
\newcolumntype{R}{>{\raggedleft\arraybackslash}X}
\usepackage{array}
\usepackage{multicol}
\usepackage{multirow}
\usepackage{siunitx}
% parse-numbers=false : accepte aussi \SI{2,5 \times 10^{-3}}{...}
\sisetup{locale=FR,output-decimal-marker={,},group-minimum-digits=4,parse-numbers=false}
\DeclareSIUnit\ohm{\ensuremath{\symup{\Omega}}} % U+2126 absent de la police
\DeclareSIUnit\division{div} % réglages d'oscilloscope (ms/div, V/div)
\DeclareSIUnit\tr{tr} % tours (vitesse de rotation, ex. \SI{1500}{\tr\per\minute})
\DeclareSIUnit\molar{M} % concentration molaire (ex. \SI{3}{\molar}, \SI{1}{\micro\molar})
\DeclareSIUnit\an{an} % année (le modèle confond parfois avec \year, un compteur TeX) — ex. \SI{5}{\kilo\meter\per\an}
\DeclareSIUnit\FCFA{FCFA} % franc CFA (ex. \SI{500}{\FCFA\per\kilo\gram})
\DeclareSIUnit\degreeBrix{\textdegree Brix} % teneur en sucre d'un sirop/jus (réfractométrie)
\DeclareSIUnit\month{mois} % le modèle utilise parfois \month, un compteur TeX, comme unité de durée
\usepackage{qrcode}
% \degree : commande fréquemment utilisée par les modèles pour un angle ou une
% température en dehors de \SI{}{} (non fournie par les paquets chargés).
\providecommand{\degree}{^{\circ}}
% \qed (fin de démonstration), utilisé par les modèles sans amsthm
\providecommand{\qed}{\hfill$\square$}
% \barre : double barre des valeurs interdites dans un tableau de variations (tkz-tab)
\providecommand{\barre}{\ensuremath{\|}}
% environnement proof (amsthm non chargé) utilisé par les modèles
\ifdefined\proof\else\newenvironment{proof}[1][Démonstration]{\par\noindent\textit{#1.}\ }{\qed\par}\fi
\usepackage{lastpage}
\usepackage{hyperref}
\hypersetup{hidelinks}

% ============================================================
% Palette « Pop » OnBuch+ — mêmes hex que la classe P (plus_ui.dart)
% ============================================================
\definecolor{poporange}{HTML}{F59321}
\definecolor{popdark}{HTML}{E07A0C}
\definecolor{popcream}{HTML}{FFF6E8}
\definecolor{popink}{HTML}{1C1714}
\definecolor{poppurple}{HTML}{7A5AE0}
\definecolor{popgreen}{HTML}{2E9E6A}
\definecolor{poppink}{HTML}{E0457B}
\definecolor{popblue}{HTML}{2F7DE1}
\definecolor{popgold}{HTML}{C98A12}
\definecolor{popmuted}{HTML}{8A7F76}
\definecolor{popline}{HTML}{EADFD2}
\definecolor{popgray}{HTML}{8A7F76}
\colorlet{popgrey}{popgray}
% Alias tolérants (le modèle nomme parfois les couleurs différemment)
\colorlet{popwhite}{white}
\colorlet{popblack}{popink}
\colorlet{popred}{poppink}
\colorlet{popyellow}{popgold}
% Teintes claires (fonds de tableaux/encadrés) — le modèle nomme parfois
% les couleurs avec un suffixe L (ex. popblueL) sans les avoir définies.
\colorlet{poporangeL}{poporange!20}
\colorlet{popdarkL}{popdark!20}
\colorlet{poppurpleL}{poppurple!20}
\colorlet{popgreenL}{popgreen!20}
\colorlet{poppinkL}{poppink!20}
\colorlet{popblueL}{popblue!20}
\colorlet{popgoldL}{popgold!20}
\colorlet{popredL}{popred!20}
\colorlet{popgrayL}{popgray!20}
% Teintes claires pour fonds de tableaux / zones
\colorlet{popblueL}{popblue!10!white}
\colorlet{poporangeL}{poporange!14!white}
\colorlet{popgreenL}{popgreen!12!white}
\colorlet{poppurpleL}{poppurple!12!white}
\colorlet{poppinkL}{poppink!12!white}
\colorlet{popgoldL}{popgold!14!white}

\pagecolor{popcream}

% ============================================================
% Typographie
% ============================================================
\defaultfontfeatures{Ligatures=TeX}
\setmainfont{PlusJakartaSans-Medium.ttf}[Path=../../../fonts/,BoldFont=PlusJakartaSans-Bold.ttf]
% Symboles, API et emojis absents de la police principale : repli sur DejaVu Sans (caractères actifs)
\newfontface\symfont{DejaVuSans.ttf}[Path=../../../fonts/]
\makeatletter
\newcommand\symactive[1]{\catcode#1=13 \begingroup\lccode`\~=#1 \lowercase{\endgroup\def~}{{\symfont\char#1}}}
\newcommand\symblank[1]{\catcode#1=13 \begingroup\lccode`\~=#1 \lowercase{\endgroup\def~}{}}
\newcommand\symhyph[1]{\catcode#1=13 \begingroup\lccode`\~=#1 \lowercase{\endgroup\def~}{\mbox{-}}}
\newcount\symc
\newcommand\symrange[2]{\symc=#1 \@whilenum\symc<#2 \do{\expandafter\symactive\expandafter{\the\symc}\advance\symc by 1 }}
\newcommand\symblankrange[2]{\symc=#1 \@whilenum\symc<#2 \do{\expandafter\symblank\expandafter{\the\symc}\advance\symc by 1 }}
\makeatother
\symrange{"0250}{"02FF}\symactive{"2713}\symactive{"2714}\symactive{"2717}\symactive{"2718}\symactive{"2610}\symactive{"2611}\symactive{"2612}
\symrange{"2600}{"27BF}
\catcode"2705=13 \begingroup\lccode`\~="2705 \lowercase{\endgroup\def~}{{\symfont\char"2713}}\symblank{"26BD}\symblank{"26A1}
\symrange{"2460}{"257F}\symactive{"223C}\symactive{"2248}\symactive{"2260}
\symactive{"01F8}\symactive{"01F9}\symactive{"1E3E}\symactive{"1E3F}
\symhyph{"2011}\symhyph{"2E1A}\symblankrange{"1F300}{"1FAFF}\symblank{"FE0F}
% Caractères chinois : WenQuanYi Zen Hei (sous-ensemble GB2312 + pinyin), gras/italique simulés
\setCJKmainfont{WenQuanYiZenHei-subset.ttf}[Path=../../../fonts/,AutoFakeBold=3,AutoFakeSlant=0.2]
\xeCJKsetup{CJKspace=false}
% Pinyin : voyelles à caron (ǎ ǐ ǒ ǔ ǖ ǘ ǚ ǜ) absentes de la police principale -> caractères actifs
\newfontface\pyfont{WenQuanYiZenHei-subset.ttf}[Path=../../../fonts/]
\newcommand\pyactive[1]{\catcode#1=13 \begingroup\lccode`\~=#1 \lowercase{\endgroup\def~}{{\pyfont\char#1}}}
\pyactive{"01CD}\pyactive{"01CE}\pyactive{"01CF}\pyactive{"01D0}\pyactive{"01D1}\pyactive{"01D2}\pyactive{"01D3}\pyactive{"01D4}\pyactive{"01D5}\pyactive{"01D6}\pyactive{"01D7}\pyactive{"01D8}\pyactive{"01D9}\pyactive{"01DA}\pyactive{"01DB}\pyactive{"01DC}
% Maths en sans-serif (Fira Math) avec chiffres et lettres droites de Plus
% Jakarta, pour que les formules se fondent dans le texte.
\usepackage[math-style=ISO,bold-style=ISO]{unicode-math}
\setmathfont{FiraMath-Regular.otf}[Path=../../../fonts/]
\setmathfont{PlusJakartaSans-Medium.ttf}[Path=../../../fonts/,range={up/num,up/{Latin,latin},\mathup}]
\usepackage{icomma} % virgule décimale sans espace en mode math
% Caractères mathématiques Unicode absents des polices de texte (Plus Jakarta,
% Archivo Black) : rendus par leur équivalent mathématique, en texte comme en
% formule — sinon ils s'affichent en carré vide (ex. « Arithmétique dans ℤ »).
\usepackage{newunicodechar}
\newunicodechar{ℝ}{\ensuremath{\mathbb{R}}}
\newunicodechar{ℤ}{\ensuremath{\mathbb{Z}}}
\newunicodechar{ℂ}{\ensuremath{\mathbb{C}}}
\newunicodechar{ℕ}{\ensuremath{\mathbb{N}}}
\newunicodechar{ℚ}{\ensuremath{\mathbb{Q}}}
\newunicodechar{𝔻}{\ensuremath{\mathbb{D}}}
\newunicodechar{α}{\ensuremath{\alpha}}
\newunicodechar{β}{\ensuremath{\beta}}
\newunicodechar{γ}{\ensuremath{\gamma}}
\newunicodechar{δ}{\ensuremath{\delta}}
\newunicodechar{ε}{\ensuremath{\varepsilon}}
\newunicodechar{θ}{\ensuremath{\theta}}
\newunicodechar{λ}{\ensuremath{\lambda}}
\newunicodechar{μ}{\ensuremath{\mu}}
\newunicodechar{σ}{\ensuremath{\sigma}}
\newunicodechar{τ}{\ensuremath{\tau}}
\newunicodechar{φ}{\ensuremath{\varphi}}
\newunicodechar{ψ}{\ensuremath{\psi}}
\newunicodechar{ω}{\ensuremath{\omega}}
\newunicodechar{Σ}{\ensuremath{\Sigma}}
% majuscules produites par \MakeUppercase dans les titres (α-aminés -> Α)
\newunicodechar{Α}{\ensuremath{\alpha}}
\newunicodechar{Β}{\ensuremath{\beta}}
\newunicodechar{Γ}{\ensuremath{\Gamma}}
\newunicodechar{Δ}{\ensuremath{\Delta}}
\newunicodechar{Ε}{\ensuremath{\varepsilon}}
\newunicodechar{Θ}{\ensuremath{\Theta}}
\newunicodechar{Λ}{\ensuremath{\Lambda}}
\newunicodechar{Μ}{\ensuremath{\mu}}
\newunicodechar{Τ}{\ensuremath{\tau}}
\newunicodechar{Φ}{\ensuremath{\Phi}}
\newunicodechar{Ψ}{\ensuremath{\Psi}}
\newunicodechar{Ω}{\ensuremath{\Omega}}
\newunicodechar{↦}{\ensuremath{\mapsto}}
\newunicodechar{∈}{\ensuremath{\in}}
\newunicodechar{∉}{\ensuremath{\notin}}
\newunicodechar{∧}{\ensuremath{\wedge}}
\newunicodechar{⇔}{\ensuremath{\Leftrightarrow}}
\newunicodechar{⟺}{\ensuremath{\Longleftrightarrow}}
\newunicodechar{⇒}{\ensuremath{\Rightarrow}}
\newunicodechar{⟹}{\ensuremath{\Longrightarrow}}
\newunicodechar{∘}{\ensuremath{\circ}}
\newunicodechar{∪}{\ensuremath{\cup}}
\newunicodechar{∩}{\ensuremath{\cap}}
\newunicodechar{≡}{\ensuremath{\equiv}}
\newunicodechar{⊂}{\ensuremath{\subset}}
\newunicodechar{⊥}{\ensuremath{\perp}}
\newunicodechar{∃}{\ensuremath{\exists}}
\newunicodechar{∀}{\ensuremath{\forall}}
\newunicodechar{∛}{\ensuremath{\sqrt[3]{\,}}}
\newunicodechar{∜}{\ensuremath{\sqrt[4]{\,}}}
\newunicodechar{⃗}{\ensuremath{{}^{\to}}}
\newunicodechar{ⁿ}{\ifmmode^{n}\else\textsuperscript{\ensuremath{n}}\fi}
\newunicodechar{ˣ}{\ifmmode^{x}\else\textsuperscript{\ensuremath{x}}\fi}
\newunicodechar{ᵇ}{\ifmmode^{b}\else\textsuperscript{\ensuremath{b}}\fi}
\newunicodechar{ʳ}{\ifmmode^{r}\else\textsuperscript{\ensuremath{r}}\fi}
\newunicodechar{ᵘ}{\ifmmode^{u}\else\textsuperscript{\ensuremath{u}}\fi}
\newunicodechar{ʸ}{\ifmmode^{y}\else\textsuperscript{\ensuremath{y}}\fi}
\newunicodechar{ᵃ}{\ifmmode^{a}\else\textsuperscript{\ensuremath{a}}\fi}
\newunicodechar{ᵉ}{\ifmmode^{e}\else\textsuperscript{\ensuremath{e}}\fi}
\newunicodechar{ᵏ}{\ifmmode^{k}\else\textsuperscript{\ensuremath{k}}\fi}
\newunicodechar{ᵖ}{\ifmmode^{p}\else\textsuperscript{\ensuremath{p}}\fi}
\newunicodechar{ᴺ}{\ifmmode^{N}\else\textsuperscript{\ensuremath{N}}\fi}
\newunicodechar{ᶜ}{\ifmmode^{c}\else\textsuperscript{\ensuremath{c}}\fi}
\newunicodechar{ʰ}{\ifmmode^{h}\else\textsuperscript{\ensuremath{h}}\fi}
\newunicodechar{ᵠ}{\ifmmode^{q}\else\textsuperscript{\ensuremath{q}}\fi}
\newunicodechar{ₙ}{\ifmmode_{n}\else\textsubscript{\ensuremath{n}}\fi}
\newunicodechar{ₐ}{\ifmmode_{a}\else\textsubscript{\ensuremath{a}}\fi}
\newunicodechar{ᵢ}{\ifmmode_{i}\else\textsubscript{\ensuremath{i}}\fi}
\newunicodechar{ᵣ}{\ifmmode_{r}\else\textsubscript{\ensuremath{r}}\fi}
\newunicodechar{ₖ}{\ifmmode_{k}\else\textsubscript{\ensuremath{k}}\fi}
\newunicodechar{ᵦ}{\ifmmode_{\beta}\else\textsubscript{\ensuremath{\beta}}\fi}
\newunicodechar{ₓ}{\ifmmode_{x}\else\textsubscript{\ensuremath{x}}\fi}
\newfontfamily\popdisplay{ArchivoBlack-Regular.ttf}[Path=../../../fonts/]
\newfontfamily\poplogo{SpaceGrotesk-Bold.ttf}[Path=../../../fonts/]
\newfontfamily\popsemi{PlusJakartaSans-SemiBold.ttf}[Path=../../../fonts/]
\newfontfamily\popxbold{PlusJakartaSans-ExtraBold.ttf}[Path=../../../fonts/]
\newfontfamily\popxboldls{PlusJakartaSans-ExtraBold.ttf}[Path=../../../fonts/,LetterSpace=3.0]

\setlist[enumerate]{leftmargin=1.7em,itemsep=5pt,topsep=4pt}
\setlist[itemize]{leftmargin=1.7em,itemsep=4pt,topsep=4pt,label={\color{poporange}\textbullet}}
\setlength{\parindent}{0pt}
\setlength{\parskip}{5pt}
\renewcommand{\arraystretch}{1.25}

% pgfplots : style Pop par défaut
\pgfplotsset{
  every axis/.append style={axis line style={popink,line width=1pt},tick style={popink},
    grid style={popline,line width=0.5pt},label style={font=\small},tick label style={font=\footnotesize}},
  popcurve/.style={poporange,line width=1.8pt,no markers,smooth},
}
% Même style disponible dans un \draw TikZ simple (hors pgfplots)
\tikzset{popcurve/.style={poporange,line width=1.8pt}}
\tikzset{popgrid/.style={popline,very thin}}
\colorlet{popcurve}{poporange} % le nom du style est parfois utilisé comme couleur

% ============================================================
% Pill / squiggle
% ============================================================
\newcommand{\poppill}[3]{%
  \tikz[baseline=-0.55ex]\node[draw=popink,line width=1.2pt,rounded corners=2.6pt,%
    fill=#2,inner xsep=6.5pt,inner ysep=3pt]{%
    \popxboldls\fontsize{7}{7}\selectfont\color{#3}\MakeUppercase{#1}};%
}
\newcommand{\popsquiggle}[1]{%
  \tikz[baseline=-0.4ex]{
    \draw[#1,line width=2.6pt,line cap=round]
      plot [smooth] coordinates {(0,0.09) (0.34,0.01) (0.68,0.21) (1.02,0.01) (1.36,0.10)};
  }%
}
% Mot-clé surligné (marqueur orange sous le texte)
\newcommand{\cle}[1]{{\popxbold\color{popdark}#1}}

% ============================================================
% En-tête / pied
% ============================================================
\pagestyle{fancy}
\fancyhf{}
\renewcommand{\headrulewidth}{0pt}
\renewcommand{\footrulewidth}{0pt}
\lhead{\raisebox{-0.15ex}{\poplogo\fontsize{13}{13}\selectfont\color{popink}OnBuch\color{poporange}+}}
\rhead{\poppill{\DOCMATIERE\ \textbullet\ \DOCNIVEAU\ \textbullet\ Leçon \DOCLECON}{white}{popink}}
\lfoot{\popxbold\fontsize{7.6}{7.6}\selectfont\color{popmuted}\MakeUppercase{Cours signé OnBuch+}\ \textbullet\ {\popsemi\DOCTITRE}}
\rfoot{\tikz[baseline=-0.6ex]\node[rounded corners=8.5pt,fill=popink,minimum height=17pt,minimum width=17pt,inner xsep=5pt,inner sep=0]{\popxbold\fontsize{8}{8}\selectfont\color{popcream}\thepage\,/\,\pageref*{LastPage}};}

% ============================================================
% Titres
% ============================================================
\newcommand{\chaptitre}[2]{%
  \begin{tcolorbox}[enhanced,colback=poporange,colframe=popink,boxrule=2.6pt,arc=11pt,
      drop shadow=popink,left=15pt,right=15pt,top=12pt,bottom=11pt,before skip=3pt,after skip=18pt]
    \poppill{#2}{popink}{popcream}\par\vspace{9pt}
    {\raggedright\hyphenpenalty=10000\exhyphenpenalty=10000\popdisplay\fontsize{22}{24}\selectfont\color{popcream}\MakeUppercase{#1}\par}
  \end{tcolorbox}
}
% Section numérotée : pastille orange avec le numéro + titre Archivo
\newcounter{popsec}
\newcommand{\coursec}[1]{%
  \stepcounter{popsec}\setcounter{popsub}{0}%
  \par\vspace{16pt}\noindent
  \begin{minipage}{\linewidth}
  \tikz[baseline=-0.7ex]\node[draw=popink,line width=1.6pt,fill=poporange,rounded corners=4pt,minimum size=22pt,inner sep=0]{\popdisplay\fontsize{12}{12}\selectfont\color{popcream}\thepopsec};%
  \hspace{8pt}{\popdisplay\fontsize{15}{17}\selectfont\color{popink}\MakeUppercase{#1}}\par
  \vspace{4pt}\noindent{\color{poporange}\rule{36pt}{3.6pt}}
  \end{minipage}\par\nopagebreak\vspace{8pt}
}
\newcounter{popsub}
\newcommand{\courssub}[1]{%
  \stepcounter{popsub}%
  \par\vspace{9pt}\noindent{\popxbold\fontsize{11.5}{13}\selectfont\color{popink}{\color{poporange}\thepopsec.\thepopsub}\hspace{6pt}#1}\par\nopagebreak\vspace{4pt}
}

% ============================================================
% Boîtes pédagogiques — même recette : fond blanc, bordure épaisse colorée,
% ombre dure encre, badge plein. Titre optionnel [#1].
% ============================================================
\tcbset{popbase/.style={breakable,enhanced,colback=white,boxrule=2.3pt,arc=9pt,
  shadow={3pt}{-3pt}{0pt}{popink},left=14pt,right=14pt,top=11pt,bottom=11pt,before skip=11pt,after skip=15pt,
  fontupper=\popsemi\fontsize{10.6}{14.5}\selectfont\color{popink},
  fontlower=\popsemi\fontsize{10.6}{14.5}\selectfont\color{popink}}}
% Coche dessinée (le glyphe ✓ n'existe pas dans les polices du document)
\newcommand{\popcheck}[1]{\tikz[baseline=-0.5ex]\draw[#1,line width=1.6pt,line cap=round,line join=round](-0.13,0.0)--(-0.04,-0.09)--(0.13,0.1);}
\providecommand{\coche}{\popcheck{popgreen}}
\providecommand{\resultat}[1]{\par\smallskip\noindent\fcolorbox{popgreen}{popgreenL}{\parbox{\dimexpr\linewidth-2\fboxsep-2\fboxrule\relax}{\textbf{Résultat.} #1}}\par\smallskip}
\newcommand{\popboxhead}[3]{\poppill{#1}{#2}{white}\ifx\relax#3\relax\else\hspace{6pt}{\popxbold\fontsize{10.6}{12}\selectfont\color{popink}#3}\fi\par\vspace{7pt}}

\newtcolorbox{aretenir}[1][]{popbase,colframe=popgreen,before upper={\popboxhead{À retenir}{popgreen}{#1}}}
% Texte en chinois (document de lecture, dialogue, extrait) : cadre
% d'étude. Un titre optionnel (source, auteur) s'affiche dans l'en-tête.
\newtcolorbox{textebox}[1][]{popbase,colframe=popblue,before upper={\popboxhead{文本 · Texte}{popblue}{#1}},after upper={}}
% Marques ✓ / ✗ (forme correcte / fautive) souvent utilisées par les modèles
\providecommand{\cross}{\textcolor{poppink}{\ensuremath{\times}}}
\providecommand{\xmark}{\cross}\providecommand{\crossmark}{\cross}\providecommand{\cmark}{\ensuremath{\checkmark}}
% Ligne de réponse / justification à compléter (inventée par les modèles)
\providecommand{\justification}[1]{\par\noindent\textit{Justification~:} \dotfill\par}
% \textipa (package tipa non chargé) : la police ne couvre pas l'API -> simple passage
\providecommand{\textipa}[1]{#1}
% Soulignement (ulem non chargé) : \uline, \ul
\providecommand{\uline}[1]{\underline{#1}}\providecommand{\ul}[1]{\underline{#1}}
% \glossaire{\item ...} inventée par les modèles : liste de glossaire
\providecommand{\glossaire}[1]{\begin{itemize}#1\end{itemize}}
% \alert (beamer) inventée par les modèles : texte mis en évidence
\providecommand{\alert}[1]{\textbf{\textcolor{poppink}{#1}}}
% variantes ASCII d'ordinaux écrites en macro
\providecommand{\iere}{\textsuperscript{re}}\providecommand{\ieme}{\textsuperscript{e}}\providecommand{\ere}{\textsuperscript{re}}\providecommand{\eme}{\textsuperscript{e}}\providecommand{\er}{\textsuperscript{er}}\providecommand{\ie}{\textsuperscript{e}}
% Ordinaux français écrits en macro par les modèles : 1\ière, 3\ième
\newcommand{\ière}{\textsuperscript{re}}\newcommand{\ième}{\textsuperscript{e}}
% \exemple (inventée par les modèles) : étiquette « Exemple : »
\providecommand{\exemple}{\textbf{Exemple~:}\ }
% \square utilisable aussi hors mode math (cases à cocher)
\AtBeginDocument{\let\squareMath\square\renewcommand{\square}{\ensuremath{\squareMath}}}
% Mot ou phrase en chinois à l'intérieur d'un paragraphe français
\newcommand{\zh}[1]{\ifmmode\text{#1}\else{#1}\fi}
\let\eng\zh \let\esp\zh \let\all\zh % alias tolérés
% flèches d'intonation inventées par les modèles
\providecommand{\textsearrow}{}\renewcommand{\textsearrow}{\ensuremath{\searrow}}\providecommand{\textnearrow}{}\renewcommand{\textnearrow}{\ensuremath{\nearrow}}
% \fr{..} : mot français (inventé par les modèles)
\providecommand{\fr}[1]{\textit{#1}}
% zhbox : encadré de texte chinois inventé par les modèles
\newenvironment{zhbox}{\begin{quote}}{\end{quote}}
% \courssubsub : titre de niveau 3 inventé par les modèles
\providecommand{\courssubsub}[1]{\par\medskip\noindent\textbf{#1}\par\smallskip}
% \zhbf : texte chinois en gras inventé par les modèles
\providecommand{\zhbf}[1]{\textbf{#1}}
% \vs : « versus » inventé par les modèles
\providecommand{\vs}{\textit{vs}\ }
% \blank : case à compléter inventée par les modèles
\providecommand{\blank}[1][]{\underline{\hspace{1.6em}}}
\newtcolorbox{exemplebox}[1][]{popbase,colframe=poporange,before upper={\popboxhead{Exemple}{poporange}{#1}}}
\newtcolorbox{definition}[1][]{popbase,colframe=popblue,before upper={\popboxhead{Définition}{popblue}{#1}}}
\newtcolorbox{propriete}[1][]{popbase,colframe=poppurple,before upper={\popboxhead{Propriété}{poppurple}{#1}}}
\newtcolorbox{methode}[1][]{popbase,colframe=popgold,before upper={\popboxhead{Méthode}{popgold}{#1}}}
\newtcolorbox{attention}[1][]{popbase,colframe=poppink,before upper={\popboxhead{Attention}{poppink}{#1}}}
\newtcolorbox{experience}[1][]{popbase,colframe=popblue,colback=popblueL,before upper={\popboxhead{Expérience}{popblue}{#1}}}
% « Le savais-tu ? » : carte violette pleine (contexte, culture, Cameroun)
\newtcolorbox{savaistu}[1][]{popbase,colback=poppurple,colframe=popink,
  fontupper=\popsemi\fontsize{10.4}{14.2}\selectfont\color{white},
  before upper={\poppill{Le savais-tu\,?}{white}{poppurple}\ifx\relax#1\relax\else\hspace{6pt}{\popxbold\color{white}#1}\fi\par\vspace{7pt}}}
% Exercice résolu : énoncé en haut, \tcblower puis la solution sur fond crème
\newcounter{popexo}\newcounter{popexr}
\newtcolorbox{exoresolu}[1][]{popbase,colframe=poporange,colbacklower=popcream,
  segmentation style={popink,line width=1.2pt,dashed},
  before upper={\stepcounter{popexr}\popboxhead{Exercice résolu \thepopexr}{poporange}{#1}},
  before lower={\poppill{Solution}{popink}{popcream}\par\vspace{6pt}}}
% Exercice (non résolu en place) — la correction est dans la partie Corrigés
\newtcolorbox{exercice}[1][]{popbase,colframe=popink,
  before upper={\stepcounter{popexo}\popboxhead{Exercice \thepopexo}{popink}{#1}}}
% Corrigé
\newtcolorbox{corrige}[1][]{popbase,colframe=popgreen,colback=popgreenL,
  before upper={\popboxhead{Corrigé}{popgreen}{#1}}}
% Objectifs / prérequis / bilan
\newtcolorbox{objectifs}{popbase,colframe=popink,colback=poporangeL,
  before upper={\popboxhead{Objectifs de la leçon}{popink}{}}}
\newtcolorbox{prerequis}{popbase,colframe=popink,colback=popblueL,
  before upper={\popboxhead{Prérequis}{popblue}{}}}
\newtcolorbox{bilan}{popbase,colframe=popink,colback=white,boxrule=2.8pt,
  before upper={\popboxhead{Fiche bilan}{poporange}{L'essentiel à connaître pour l'examen}}}
% Figure encadrée avec légende
\newtcolorbox{popfigure}[1][]{enhanced,breakable=false,colback=white,colframe=popline,boxrule=1.4pt,arc=8pt,
  left=10pt,right=10pt,top=10pt,bottom=8pt,before skip=10pt,after skip=12pt,halign=center,
  lower separated=false,halign lower=center,
  fontlower=\popsemi\fontsize{9}{11.5}\selectfont\color{popmuted},
  #1}
\newcommand{\legende}[1]{\par\vspace{6pt}{\popsemi\fontsize{9}{11.5}\selectfont\color{popmuted}#1}}

% ============================================================
% Signature OnBuch+
% ============================================================
\newcommand{\poplogoinline}[1]{{\poplogo\fontsize{#1}{#1}\selectfont\color{popink}OnBuch\color{poporange}+}}
\newcommand{\signatureonbuch}{%
  \begin{tcolorbox}[enhanced,colback=popink,colframe=popink,boxrule=2.2pt,arc=10pt,
      drop shadow=poporange,left=14pt,right=14pt,top=10pt,bottom=10pt,before skip=0pt,after skip=0pt]
    \tikz[baseline=-0.7ex]\node[circle,fill=popgreen,draw=popcream,line width=1.3pt,minimum size=22pt,inner sep=0]{\popcheck{white}};%
    \hspace{9pt}\begin{minipage}[c]{0.62\linewidth}
      {\popxbold\fontsize{12}{13}\selectfont\color{popcream}Signé\ \textbullet\ {\poplogo OnBuch\color{poporange}+}}\par\vspace{2pt}
      {\raggedright\popsemi\fontsize{8}{10.5}\selectfont\color{popline}\MakeUppercase{Équipe pédagogique OnBuch+}\\\MakeUppercase{Conforme au programme officiel MINESEC}\par}
    \end{minipage}\hfill
    {\poplogo\fontsize{22}{22}\selectfont\color{popcream}OnBuch\color{poporange}+}
  \end{tcolorbox}%
}

% ============================================================
% Page de garde
% #1 titre · #2 matière · #3 niveau · #4 module · #5 n° leçon · #6 durée ·
% #7 description / objectifs (texte ou itemize)
% ============================================================
\newcommand{\pagegarde}[7]{%
  \thispagestyle{empty}
  \newgeometry{a4paper,margin=1.7cm,top=1.6cm,bottom=1.4cm}%
  \begin{tikzpicture}[remember picture,overlay]
    \fill[poporange!18] ([xshift=-3.2cm,yshift=-3.6cm]current page.north east) circle (4.6cm);
    \fill[poppurple!14] ([xshift=2.4cm,yshift=7.6cm]current page.south west) circle (3.8cm);
  \end{tikzpicture}%
  {\poplogo\fontsize{30}{30}\selectfont\color{popink}OnBuch\color{poporange}+}\hfill
  \raisebox{0.6ex}{\poppill{Cours signé \textbullet\ Premium}{popink}{popcream}}\par
  \vspace{1pt}\popsquiggle{poppurple}\par
  \vspace{8pt}
  \begin{tcolorbox}[enhanced,colback=poporange,colframe=popink,boxrule=2.8pt,arc=13pt,
      drop shadow=popink,left=18pt,right=18pt,top=12pt,bottom=13pt,after skip=10pt]
    \poppill{#2 \textbullet\ #3}{popink}{popcream}\hspace{5pt}\poppill{#4}{white}{popink}\par\vspace{8pt}
    {\popdisplay\fontsize{14}{14}\selectfont\color{popink}LEÇON #5}\par\vspace{4pt}
    {\raggedright\hyphenpenalty=10000\exhyphenpenalty=10000\popdisplay\fontsize{25}{27}\selectfont\color{popcream}\MakeUppercase{#1}\par}
    \vspace{8pt}
    {\popxbold\fontsize{9.5}{9.5}\selectfont\color{popink}\MakeUppercase{Durée conseillée : #6}}
  \end{tcolorbox}
  \begin{tcolorbox}[enhanced,colback=white,colframe=popink,boxrule=2.2pt,arc=10pt,
      drop shadow=popink,left=16pt,right=16pt,top=10pt,bottom=10pt,after skip=8pt]
    \poppill{Dans cette leçon}{poppurple}{white}\par\vspace{5pt}
    \popsemi\fontsize{9.3}{12.6}\selectfont\color{popink}#7
  \end{tcolorbox}
  \noindent\begin{minipage}[t]{0.487\linewidth}
    \begin{tcolorbox}[enhanced,colback=popgreen,colframe=popink,boxrule=2.2pt,arc=10pt,
        drop shadow=popink,left=11pt,right=11pt,top=10pt,bottom=10pt,height=3.1cm,valign=center]
      \tcbox[colback=white,colframe=popink,boxrule=1.3pt,arc=5pt,boxsep=0pt,nobeforeafter,
        left=3pt,right=3pt,top=3pt,bottom=3pt,valign=center]{\qrcode[height=1.9cm]{https://play.google.com/store/apps/details?id=com.luvvix.onbuch&hl=fr}}%
      \hspace{8pt}\begin{minipage}[c]{0.5\linewidth}
        {\popdisplay\fontsize{11}{12.5}\selectfont\color{white}\MakeUppercase{Télécharge OnBuch}}\par\vspace{4pt}
        {\raggedright\popsemi\fontsize{8.4}{11}\selectfont\color{white}Gratuit \textbullet\ Google Play\\Tuteur IA, annales, concours, résultats\par}
      \end{minipage}
    \end{tcolorbox}
  \end{minipage}\hfill
  \begin{minipage}[t]{0.487\linewidth}
    \begin{tcolorbox}[enhanced,colback=poppurple,colframe=popink,boxrule=2.2pt,arc=10pt,
        drop shadow=popink,left=11pt,right=11pt,top=10pt,bottom=10pt,height=3.1cm,valign=center]
      \tikz[baseline=-0.7ex]\node[circle,fill=white,draw=popink,line width=1.3pt,minimum size=1.6cm,inner sep=0]{\popdisplay\fontsize{24}{24}\selectfont\color{poppurple}+};%
      \hspace{8pt}\begin{minipage}[c]{0.6\linewidth}
        {\popdisplay\fontsize{11}{12.5}\selectfont\color{white}\MakeUppercase{Passe à OnBuch+}}\par\vspace{4pt}
        {\raggedright\popsemi\fontsize{8.4}{11}\selectfont\color{white}Tous les cours signés, Tuteur IA illimité, fiches PDF — onglet OnBuch+ de l'app\par}
      \end{minipage}
    \end{tcolorbox}
  \end{minipage}
  \vspace{8pt}
  \popcontenu
  \vfill
  \signatureonbuch
  \restoregeometry
  \newpage
}

% Rangée « Ce cours contient » (page de garde)
\newcommand{\poptile}[3]{% #1 couleur #2 gros texte #3 légende
  \begin{tcolorbox}[enhanced,colback=#1,colframe=popink,boxrule=2pt,arc=9pt,shadow={2.5pt}{-2.5pt}{0pt}{popink},
      left=6pt,right=6pt,top=9pt,bottom=9pt,halign=center,valign=center,height=1.75cm,width=\linewidth,nobeforeafter]
    {\popdisplay\fontsize{12.5}{13}\selectfont\color{white}\MakeUppercase{#2}}\par\vspace{3pt}
    {\centering\popsemi\fontsize{7.8}{9.5}\selectfont\color{white}#3\par}
  \end{tcolorbox}}
\newcommand{\popcontenu}{%
  \par\noindent{\popxboldls\fontsize{7.5}{7.5}\selectfont\color{popmuted}CE COURS SIGNÉ CONTIENT}\par\vspace{7pt}
  \noindent\begin{minipage}[t]{0.235\linewidth}\poptile{popblue}{Cours}{complet, illustré, pas à pas}\end{minipage}\hfill
  \begin{minipage}[t]{0.235\linewidth}\poptile{poporange}{Méthodes}{et exercices résolus}\end{minipage}\hfill
  \begin{minipage}[t]{0.235\linewidth}\poptile{poppink}{Exercices}{type examen + corrigés}\end{minipage}\hfill
  \begin{minipage}[t]{0.235\linewidth}\poptile{popgreen}{Fiche bilan}{l'essentiel à retenir}\end{minipage}\par}

% Sommaire
\newtcolorbox{sommaire}{enhanced,colback=white,colframe=popink,boxrule=2.2pt,arc=10pt,
  drop shadow=popink,left=18pt,right=18pt,top=13pt,bottom=13pt,before skip=6pt,after skip=20pt,
  before upper={\poppill{Sommaire}{poppurple}{white}\par\vspace{9pt}\popsemi\fontsize{10.6}{14.5}\selectfont\color{popink}}}

% Signature de fin de cours — poussée tout en bas de la dernière page
\newcommand{\findecours}{%
  \par\vfill\nopagebreak
  \begin{center}\popsquiggle{poporange}\end{center}\vspace{4pt}
  \signatureonbuch
}
\AtBeginDocument{\def\llbracket{\mathopen{[\mkern-3.5mu[}}\def\rrbracket{\mathclose{]\mkern-3.5mu]}}} % intervalles d'entiers (glyphe ⟦ absent de la police)
% Glyphes absents de Fira Math : redéfinis à partir de glyphes disponibles
\AtBeginDocument{%
  \def\setminus{\mathbin{\backslash}}\let\smallsetminus\setminus
  \def\vdots{\mathord{\vbox{\baselineskip3.2pt\lineskiplimit0pt\kern4pt\hbox{.}\hbox{.}\hbox{.}}}}%
  \def\ddots{\mathinner{\mkern1mu\raise6.5pt\hbox{.}\mkern2mu\raise3.6pt\hbox{.}\mkern2mu\raise0.7pt\hbox{.}\mkern1mu}}%
  \def\triangle{\mathord{\scalebox{0.9}{$\symup{\Delta}$}}}%
  \def\nabla{\mathord{\raisebox{\depth}{\scalebox{1}[-1]{$\symup{\Delta}$}}}}%
  \def\checkmark{\popcheck{popdark}}%
  \def\star{\mathbin{\ast}}\def\bigstar{\mathord{\ast}}%
  \def\diamond{\mathbin{\scalebox{0.6}{$\Diamond$}}}%
  \def\bigcup{\mathop{\mathchoice{\raisebox{-0.3em}{\scalebox{1.7}{$\cup$}}}{\scalebox{1.25}{$\cup$}}{\cup}{\cup}}\displaylimits}%
  \def\bigcap{\mathop{\mathchoice{\raisebox{-0.3em}{\scalebox{1.7}{$\cap$}}}{\scalebox{1.25}{$\cap$}}{\cap}{\cap}}\displaylimits}%
  \def\lesseqqgtr{\mathrel{\lessgtr}}\def\bigoplus{\mathop{\oplus}\displaylimits}%
}
\providecommand{\ssi}{si et seulement si} % abréviation fréquente
\AtBeginDocument{\providecommand{\wideparen}{\overparen}} % arc de cercle

%% FIGFIX-BEGIN v5 — correctifs globaux des figures (docs/onbuch-generation/tools/figfix)
\usepackage{adjustbox}\usepackage{etoolbox}\tcbuselibrary{raster}
% toute figure TikZ/pgfplots plus large que la ligne est réduite à la largeur disponible
\BeforeBeginEnvironment{tikzpicture}{\begin{adjustbox}{max width=\linewidth}}
\AfterEndEnvironment{tikzpicture}{\end{adjustbox}}
% légendes pgfplots lisibles par-dessus les courbes
\pgfplotsset{every axis/.append style={legend style={fill=white,fill opacity=0.92,text opacity=1,draw=black!15,inner sep=3pt}}}
% étiquettes d'arêtes (syntaxe edge["..."]) avec fond blanc
\tikzset{every edge quotes/.append style={fill=white,inner sep=1.2pt}}
%% FIGFIX-END
\begin{document}

\pagegarde{Leçon 19 — Expression écrite : offre de services aux usagers des administrations (collectivités locales)}{Chinois}{Tle A}{Module 3 : Citoyenneté et ouverture au monde}{ZH19}{2 h}{Cette leçon d'expression écrite apprend à rédiger en chinois une offre de services destinée aux usagers des administrations locales (mairies, communautés urbaines, conseils). Elle mobilise l'écriture correcte des caractères du thème, les connecteurs de texte et la traduction (thème et version) de phrases simples.
\vspace{4pt}
\begin{multicols}{2}\small
\begin{itemize}
\item À la fin de cette leçon, je sais identifier la structure d'une offre de services administrative en chinois (titre, objet, services, horaires, contact, formule de clôture)
\item je sais utiliser le lexique des services aux usagers (办理, 服务, 窗口, 申请, 材料, 时间...)
\item je sais écrire correctement les caractères clés du thème en respectant radicaux et ordre des traits
\item je sais distinguer les caractères proches (办/为, 请/清, 务/各, 间/问)
\item je sais employer les connecteurs utiles (首先, 其次, 另外, 最后, 如有...请...) dans un texte administratif
\item je sais traduire du français vers le chinois des phrases simples sur les services publics (thème)
\end{itemize}\end{multicols}}

\begin{sommaire}
\begin{multicols}{2}
\begin{enumerate}
\item Découverte : un avis de services à la mairie
\item Lexique des services aux usagers
\item Écriture des caractères du thème
\item Structure et connecteurs de l'offre de services
\item Thème et version guidés
\item Production guidée et évaluation
\item Activité d'intégration
\item Méthodes et exercices résolus
\item Exercices
\item Corrigés des exercices
\item Fiche bilan
\end{enumerate}
\end{multicols}
\end{sommaire}

\begin{objectifs}
\begin{itemize}
\item À la fin de cette leçon, je sais identifier la structure d'une offre de services administrative en chinois (titre, objet, services, horaires, contact, formule de clôture)
\item je sais utiliser le lexique des services aux usagers (办理, 服务, 窗口, 申请, 材料, 时间...)
\item je sais écrire correctement les caractères clés du thème en respectant radicaux et ordre des traits
\item je sais distinguer les caractères proches (办/为, 请/清, 务/各, 间/问)
\item je sais employer les connecteurs utiles (首先, 其次, 另外, 最后, 如有...请...) dans un texte administratif
\item je sais traduire du français vers le chinois des phrases simples sur les services publics (thème)
\item je sais traduire du chinois vers le français un avis ou une offre de services simple (version)
\item je sais produire un avis d'offre de services complet et correct d'environ 80 à 120 caractères
\end{itemize}
\end{objectifs}

\begin{prerequis}
\begin{itemize}
\item Lexique de la ville et des lieux publics (Seconde/Première : 市政府, 办公室, 地址)
\item Expression de l'heure et des jours (星期, 点, 上午/下午)
\item Structures de politesse et verbes modaux (请, 可以, 需要, 应该)
\item Nombres et classificateurs courants
\item Connecteurs simples vus en Première (因为...所以..., 如果...就...)
\end{itemize}
\end{prerequis}

\newpage

\coursec{Découverte : un avis de services à la mairie}

À la mairie de Yaoundé ou au council de Buea, un usager vient chercher un acte de naissance, payer une taxe ou déposer une plainte. Comment l'administration locale l'informe-t-elle de ses services ? Au Cameroun, on affiche des avis au tableau d'affichage, on annonce à la radio communautaire, on publie sur les réseaux sociaux. En Chine, les collectivités locales affichent leurs offres de services — horaires, guichets, démarches en ligne — dans des avis rédigés selon des codes précis, avec un titre officiel, une énumération numérotée et une formule de politesse finale.

Savoir lire, puis rédiger en chinois une offre de services claire et polie, c'est une compétence professionnelle concrète — et un thème classique d'expression écrite au Bac. C'est toute l'ambition de cette leçon.

\textbf{Problématique :} quelles sont les pièces obligatoires d'un avis de services en chinois, et dans quel ordre se présentent-elles ?

\courssub{Lecture du texte support : un avis du Centre de services aux citoyens}

Avant toute explication, lisons le texte. Lis-le d'abord à voix haute en respectant les tons, puis essaie de le traduire sans dictionnaire : tu connais déjà la plupart des mots.

\begin{textebox}[市民服务中心公告 — Avis du Centre de services aux citoyens]
市民服务中心公告。\\
Shìmín fúwù zhōngxīn gōnggào.\\
\emph{Avis du Centre de services aux citoyens.}\\[4pt]
为了方便居民，本中心提供以下服务：一、办理身份证和户口本；二、办理出生证明；三、接受居民的意见和建议。\\
Wèile fāngbiàn jūmín, běn zhōngxīn tígōng yǐxià fúwù : yī, bànlǐ shēnfènzhèng hé hùkǒuběn ; èr, bànlǐ chūshēng zhèngmíng ; sān, jiēshòu jūmín de yìjian hé jiànyì.\\
\emph{Pour faciliter la vie des habitants, notre centre offre les services suivants : 1. délivrance des cartes d'identité et des livrets de famille ; 2. délivrance des actes de naissance ; 3. réception des avis et suggestions des habitants.}\\[4pt]
办公时间：星期一至星期五，上午八点到下午五点。\\
Bàngōng shíjiān : xīngqīyī zhì xīngqīwǔ, shàngwǔ bā diǎn dào xiàwǔ wǔ diǎn.\\
\emph{Horaires : du lundi au vendredi, de 8 h à 17 h.}\\[4pt]
地址：和平路十二号。电话：12345。\\
Dìzhǐ : Hépíng lù shí'èr hào. Diànhuà : yāo èr sān sì wǔ.\\
\emph{Adresse : 12, rue de la Paix. Téléphone : 12345.}\\[4pt]
如有问题，请随时联系我们。谢谢合作！\\
Rú yǒu wèntí, qǐng suíshí liánxì wǒmen. Xièxie hézuò !\\
\emph{Pour toute question, contactez-nous à tout moment. Merci de votre coopération !}
\end{textebox}

\textbf{Glossaire du texte}

\begin{center}
\begin{tabularx}{\linewidth}{|X|X|X|X|}
\hline
\rowcolor{popblueL} Caractères & Pinyin & Nature & Traduction \\
\hline
市民 & shìmín & nom & citoyen, habitant de la ville \\
\hline
服务 & fúwù & nom / verbe & service ; servir \\
\hline
中心 & zhōngxīn & nom & centre \\
\hline
公告 & gōnggào & nom & avis officiel, annonce publique \\
\hline
为了 & wèile & préposition & pour, afin de \\
\hline
方便 & fāngbiàn & adjectif / verbe & pratique ; faciliter \\
\hline
居民 & jūmín & nom & habitant, résident \\
\hline
本 & běn & déterminant / pronom & notre (officiel : « le présent »), ce \\
\hline
提供 & tígōng & verbe & fournir, offrir \\
\hline
以下 & yǐxià & locution & suivant, ci-dessous \\
\hline
办理 & bànlǐ & verbe & effectuer (une démarche), délivrer \\
\hline
身份证 & shēnfènzhèng & nom & carte d'identité \\
\hline
户口本 & hùkǒuběn & nom & livret de famille (registre de résidence) \\
\hline
出生证明 & chūshēng zhèngmíng & nom & acte de naissance \\
\hline
接受 & jiēshòu & verbe & recevoir, accepter \\
\hline
意见 & yìjian & nom & avis, opinion \\
\hline
建议 & jiànyì & nom / verbe & suggestion ; suggérer \\
\hline
办公时间 & bàngōng shíjiān & nom & horaires de bureau \\
\hline
至 & zhì & préposition / conjonction & jusqu'à (soutenu, = 到) \\
\hline
地址 & dìzhǐ & nom & adresse \\
\hline
如有 & rú yǒu & locution conjonctive & si (jamais) il y a, en cas de \\
\hline
随时 & suíshí & adverbe & à tout moment \\
\hline
联系 & liánxì & verbe & contacter \\
\hline
合作 & hézuò & verbe / nom & coopérer ; coopération \\
\hline
\end{tabularx}
\end{center}

\begin{definition}[Trois mots-clés de l'administration]
\begin{itemize}
\item \cle{公告} \zh{公告} (gōnggào) : avis officiel affiché ou publié par une administration. Composé de \zh{公} (gōng, « public ») et \zh{告} (gào, « annoncer ») : littéralement « annonce publique ». C'est le mot qui figure en titre de tout avis administratif chinois.
\item \cle{服务} \zh{服务} (fúwù) : service. On le retrouve partout : \zh{服务中心} (centre de services), \zh{服务员} (serveur, employé de service), \zh{为人民服务} (servir le peuple).
\item \cle{居民} \zh{居民} (jūmín) : habitant, résident. \zh{居} signifie « habiter », \zh{民} signifie « peuple ». À ne pas confondre avec \zh{市民} (shìmín, habitant de la ville, citoyen).
\end{itemize}
\end{definition}

\courssub{Repérage global : qui écrit ? à qui ? pour quoi ?}

Toute lecture d'un texte administratif commence par trois questions simples. Répondons-y pour notre avis.

\begin{center}
\begin{tabularx}{\linewidth}{|X|X|X|}
\hline
\rowcolor{popblueL} Question & Réponse & Indice dans le texte \\
\hline
Qui écrit ? (émetteur) & Le Centre de services aux citoyens & \zh{市民服务中心} et \zh{本中心} (« notre centre ») \\
\hline
À qui ? (destinataire) & Les habitants, les usagers & \zh{居民} (« les habitants »), \zh{我们} (« nous ») \\
\hline
Pour quoi ? (fonction) & Informer des services offerts et des modalités pratiques & \zh{提供以下服务} (« offre les services suivants ») \\
\hline
\end{tabularx}
\end{center}

La fonction du texte est donc \cle{informative} et \cle{incitative} : il informe (quels services ? quand ? où ?) et il invite à l'action (\zh{请随时联系我们}, « contactez-nous à tout moment »). C'est exactement la double fonction de l'avis d'offre de services que tu devras rédiger au Bac.

\courssub{Repérage des marqueurs : l'architecture de l'avis}

Un avis administratif chinois n'est pas un texte continu : c'est une \cle{architecture} en cinq blocs, toujours dans le même ordre. Observons-les dans notre texte.

\begin{enumerate}
\item \textbf{Le titre} : \zh{公告} (gōnggào). Un seul mot suffit, souvent précédé du nom de l'administration : \zh{市民服务中心公告}. En français, on écrirait « AVIS » ou « COMMUNIQUÉ » en haut de l'affiche.
\item \textbf{L'objet} : introduit par \zh{为了...} (wèile..., « pour, afin de ») : \zh{为了方便居民} (« pour faciliter la vie des habitants »). Cette formule justifie l'avis et annonce sa finalité.
\item \textbf{La liste des services} : énumération numérotée avec les chiffres chinois suivis du point de suspension énumératif : \zh{一、二、三} (yī, èr, sān). Chaque point commence par un verbe d'action : \zh{办理} (délivrer), \zh{接受} (recevoir).
\item \textbf{Les informations pratiques} : présentées sous forme d'étiquettes suivies de deux-points : \zh{办公时间：} (horaires), \zh{地址：} (adresse), \zh{电话：} (téléphone).
\item \textbf{La clôture} : formule de politesse stéréotypée : \zh{谢谢合作！} (Xièxie hézuò !, « Merci de votre coopération ! »), souvent précédée d'une invitation à contacter l'administration : \zh{如有问题，请随时联系我们。}
\end{enumerate}

\begin{popfigure}
\begin{tikzpicture}[
  node distance=0.55cm,
  bloc/.style={rectangle, rounded corners=4pt, draw=popblue, fill=popblueL, align=center, text width=8.6cm, inner sep=5pt, font=\small},
  fleche/.style={-{Stealth[length=2.5mm]}, thick, popdark}
]
\node[bloc, fill=popgoldL, draw=popgold] (titre) {\textbf{1. Titre} : 公告 (gōnggào) — « Avis »};
\node[bloc, below=of titre] (objet) {\textbf{2. Objet} : 为了方便居民… (wèile…) — « Pour faciliter la vie des habitants… »};
\node[bloc, below=of objet] (liste) {\textbf{3. Liste des services} : 一、办理…；二、办理…；三、接受…};
\node[bloc, below=of liste] (infos) {\textbf{4. Infos pratiques} : 办公时间：… 地址：… 电话：…};
\node[bloc, fill=popgreenL, draw=popgreen, below=of infos] (clot) {\textbf{5. Clôture} : 谢谢合作！— « Merci de votre coopération ! »};
\draw[fleche] (titre) -- (objet);
\draw[fleche] (objet) -- (liste);
\draw[fleche] (liste) -- (infos);
\draw[fleche] (infos) -- (clot);
\end{tikzpicture}
\legende{Organigramme : les cinq blocs d'un avis de services en chinois}
\end{popfigure}

\begin{methode}[Lire un texte administratif en trois passes]
Face à un avis chinois inconnu, ne lis pas mot à mot du début à la fin. Procède en trois passes :
\begin{enumerate}
\item \textbf{Première passe — titre et émetteur.} Cherche \zh{公告} ou \zh{通知} (avis) et le nom de l'administration qui l'accompagne. Tu sais immédiatement \emph{qui parle} et de quel type de texte il s'agit.
\item \textbf{Deuxième passe — la liste.} Repère les chiffres énumératifs \zh{一、二、三} et lis le verbe qui ouvre chaque point (\zh{办理}, \zh{接受}, \zh{提供}…). Tu obtiens le \emph{cœur du message} : les services offerts.
\item \textbf{Troisième passe — les informations pratiques.} Repère les étiquettes suivies de deux-points : \zh{时间：}, \zh{地址：}, \zh{电话：}. Tu recueilles \emph{quand, où, comment}.
\end{enumerate}
La formule finale (\zh{谢谢合作！}) ne demande pas de traduction fine : c'est une clôture de politesse.
\end{methode}

\begin{popfigure}
\begin{center}\tcbox[colback=popink,colframe=popink,coltext=white,arc=9pt,boxsep=4pt,left=6pt,right=6pt]{\bfseries\small 服务 fúwù (services)}\end{center}
\vspace{-2pt}
\begin{tcbraster}[raster columns=2,raster equal height=rows,raster column skip=6pt,raster row skip=6pt,enhanced,blankest]
\begin{tcolorbox}[enhanced,colback=popblue,colframe=popblue,coltext=white,boxrule=0.9pt,arc=3pt,left=4pt,right=4pt,top=3pt,bottom=3pt,boxsep=1pt,halign=left]\bfseries\scriptsize 办理证件 bànlǐ zhèngjiàn délivrer les documents\end{tcolorbox}
\begin{tcolorbox}[enhanced,colback=poporange,colframe=poporange,coltext=white,boxrule=0.9pt,arc=3pt,left=4pt,right=4pt,top=3pt,bottom=3pt,boxsep=1pt,halign=left]\bfseries\scriptsize 接受意见 jiēshòu yìjian recevoir les avis\end{tcolorbox}
\begin{tcolorbox}[enhanced,colback=popgreen,colframe=popgreen,coltext=white,boxrule=0.9pt,arc=3pt,left=4pt,right=4pt,top=3pt,bottom=3pt,boxsep=1pt,halign=left]\bfseries\scriptsize 办公时间 bàngōng shíjiān horaires\end{tcolorbox}
\begin{tcolorbox}[enhanced,colback=poppurple,colframe=poppurple,coltext=white,boxrule=0.9pt,arc=3pt,left=4pt,right=4pt,top=3pt,bottom=3pt,boxsep=1pt,halign=left]\bfseries\scriptsize 联系方式 liánxì fāngshì contacts\end{tcolorbox}
\end{tcbraster}

\legende{Carte mentale lexicale autour de 服务 (fúwù, « service »)}
\end{popfigure}

\courssub{Premier point de langue : 为了 + but}

L'avis s'ouvre sur la structure \zh{为了方便居民} : c'est la marque typique de l'objet dans un texte officiel.

\begin{propriete}[La structure 为了 + but + , + action]
\textbf{Structure :} 为了 + nom ou verbe (le but) + ，+ proposition principale (l'action).\par
\textbf{Emploi :} exprime le but poursuivi, « pour, afin de ». Registre plutôt soutenu : très fréquent dans les avis, discours et textes officiels.\par
\textbf{Comparaison avec le français :} le français place souvent le but après l'action (« notre centre offre ces services \emph{pour} faciliter… ») ; le chinois place presque toujours le but \emph{avant} l'action dans ce type de texte.
\end{propriete}

\begin{exemplebox}[为了 en action]
\zh{为了方便居民，本中心提供以下服务。} \emph{Wèile fāngbiàn jūmín, běn zhōngxīn tígōng yǐxià fúwù.} « Pour faciliter la vie des habitants, notre centre offre les services suivants. »\\[3pt]
\zh{为了帮助学生，学校开设了中文课。} \emph{Wèile bāngzhù xuésheng, xuéxiào kāishè le Zhōngwén kè.} « Pour aider les élèves, l'école a ouvert un cours de chinois. »\\[3pt]
\zh{为了大家的安全，请排队。} \emph{Wèile dàjiā de ānquán, qǐng páiduì.} « Pour la sécurité de tous, veuillez faire la queue. »
\end{exemplebox}

\courssub{Écriture des caractères du thème}

Maîtriser l'écriture des caractères clés du domaine administratif permet de rédiger un avis lisible et correct. Nous travaillons sur six caractères à haute fréquence : \zh{公}, \zh{告}, \zh{办}, \zh{证}, \zh{时}, \zh{电}.

\begin{definition}[Analyse graphique : traits, radicaux, pièges]
\begin{center}
\begin{tabularx}{\linewidth}{|X|X|X|X|X|}
\hline
\rowcolor{popblueL} Caractère & Pinyin & Radical (clé) & Nb traits & Ordre des traits principal / Pièges \\
\hline
\zh{公} & gōng & \zh{八} (huit, n°12) & 4 & Horizontal → verticale descendante → deux traits « huit » (gauche puis droite). \textbf{Piège} : ne pas confondre avec \zh{翁} (wēng, vieillard) ni \zh{瓜} (guā, pastèque). \\
\hline
\zh{告} & gào & \zh{口} (bouche, n°30) & 7 & Haut : \zh{牛} (vache : horizontal, verticale, horizontal, horizontal) → Bas : \zh{口} (bouche : verticale, angle+horizontal, horizontal). \textbf{Piège} : le trait vertical du \zh{牛} ne traverse pas le \zh{口}. \\
\hline
\zh{办} & bàn & \zh{力} (force, n°19) & 8 & Partie gauche \zh{卜} (divination : point, horizontal, crochet) + \zh{力} (force : crochet horizontal, point). \textbf{Sens} : « force appliquée à une tâche » → faire, traiter. \\
\hline
\zh{证} & zhèng & \zh{讠} (parole, n°149) & 6 & Radical \zh{讠} (deux traits : point, courbe montante) + phonétique \zh{正} (correct : horizontal, verticale, horizontal, horizontal). \textbf{Simplifié} : traditionnel \zh{證} (phonétique \zh{登} + \zh{言}). \\
\hline
\zh{时} & shí & \zh{日} (soleil/jour, n°72) & 7 & \zh{日} (6 traits : verticale, angle+horizontal, horizontal, horizontal, horizontal) + \zh{寸} (pouce : horizontal, crochet, point). \textbf{Sens} : le temps qui passe (soleil + mesure). \\
\hline
\zh{电} & diàn & \zh{田} (champ, n°102) & 5 & \zh{田} (5 traits : verticale, angle+horizontal, horizontal, horizontal, verticale) + \zh{乙} (deuxième trait : courbe descendante) en bas au centre. \textbf{Simplifié} : traditionnel \zh{電} (pluie + éclair). \\
\hline
\end{tabularx}
\end{center}
\end{definition}

\begin{experience}[Atelier écriture : dictée muette et contrôle croisé]
\begin{enumerate}
\item L'enseignant dicte en français : « avis », « service », « traiter », « carte d'identité », « heure », « téléphone ». Les élèves écrivent les caractères au tableau ou sur ardoise.
\item Échange d'ardoises : chaque élève vérifie l'ordre des traits (numérote les traits au stylo rouge) et le radical de son camarade.
\item Collectif : on relève les erreurs fréquentes (traits dans le désordre sur \zh{告}, radical \zh{力} oublié sur \zh{办}, \zh{乙} mal placé sur \zh{电}).
\end{enumerate}
\end{experience}

\begin{attention}[Caractères proches : ne les confondez pas]
\begin{itemize}
\item \zh{公} (gōng, public) vs \zh{翁} (wēng, vieillard : \zh{公} + \zh{羽} plume) vs \zh{瓜} (guā, pastèque : trait vertical central traversant).
\item \zh{证} (zhèng, preuve) vs \zh{正} (zhèng, correct) — \zh{证} a le radical \zh{讠} à gauche.
\item \zh{时} (shí, temps) vs \zh{待} (dài, attendre : \zh{彳} + \zh{寺}) vs \zh{特} (tè, spécial : \zh{牛} + \zh{寸}).
\item \zh{电} (diàn, électricité) vs \zh{由} (yóu, par/à cause de : \zh{田} + \zh{人} au lieu de \zh{乙}).
\end{itemize}
\end{attention}

\courssub{Structure et connecteurs de l'offre de services}

Au-delà de \zh{为了}, un avis administratif utilise des connecteurs et des structures figées pour organiser l'information avec clarté et autorité.

\begin{propriete}[L'énumération formelle : 一、二、三 + 、]
En chinois administratif, l'énumération utilise les chiffres chinois suivis du signe \zh{、} (point de suspension énumératif, Unicode U+3001), \textbf{jamais} la virgule \zh{，} ni le point \zh{。}.
\par \textbf{Schéma :} \zh{一、} [Verbe + Objet] \zh{；} \zh{二、} [Verbe + Objet] \zh{；} \zh{三、} [Verbe + Objet] \zh{。}
\par Les items sont séparés par des points-virgules \zh{；} (ou virgules \zh{，} si courts), le dernier se termine par un point \zh{。}.
\end{propriete}

\begin{propriete}[L'étiquette d'information : Nom + ： + Valeur]
Les informations pratiques (horaires, adresse, téléphone, site web) suivent le modèle : \textbf{Étiquette + deux-points chinois \zh{：} + Contenu}.
\par \textbf{Exemples :} \zh{办公时间：星期一至星期五} ; \zh{地址：和平路十二号} ; \zh{电话：12345} ; \zh{网址：www.example.cn}。
\end{propriete}

\begin{propriete}[Formules de clôture standardisées]
Deux formules ferment quasi systématiquement l'avis :
\begin{enumerate}
\item \textbf{Invitation au contact :} \zh{如有问题，请随时联系我们。} (Rú yǒu wèntí, qǐng suíshí liánxì wǒmen.) — « En cas de question, merci de nous contacter à tout moment. »
\item \textbf{Remerciement de coopération :} \zh{谢谢合作！} (Xièxie hézuò !) — « Merci de votre coopération ! » (formule figée, ne pas traduire littéralement « merci de coopérer » mais bien comme une marque de politesse administrative).
\end{enumerate}
\end{propriete}

\begin{exemplebox}[Connecteurs utiles pour l'expression écrite]
\begin{center}
\begin{tabularx}{\linewidth}{|X|X|X|}
\hline
\rowcolor{popblueL} Connecteur chinois & Pinyin & Emploi / Traduction \\
\hline
为了 & wèile & But (en tête de phrase) : « Pour, afin de » \\
\hline
本 & běn & Déterminant officiel : « le présent, notre » (\zh{本中心}, \zh{本单位}) \\
\hline
以下 & yǐxià & « Ci-dessous, suivants » (\zh{以下服务}) \\
\hline
及 & jí & « Et » (formel, lie deux noms) : \zh{身份证及户口本} \\
\hline
或 & huò & « Ou » (formel) : \zh{电话或网络预约} \\
\hline
请 & qǐng & Impératif poli : « Veuillez, merci de » (\zh{请排队}, \zh{请携带证件}) \\
\hline
须知 & xūzhī & « Doit savoir / Note importante » (titre de section) \\
\hline
特此公告 & tè cǐ gōnggào & Formule de clôture solennelle : « Spécialement par la présente, avis est donné » \\
\hline
\end{tabularx}
\end{center}
\end{exemplebox}

\courssub{Thème et version guidés}

Applique les structures et le vocabulaire vus ci-dessus. Travaille d'abord seul, puis compare en binôme.

\begin{exercice}[Thème : Français → Chinois]
Traduis les phrases suivantes en chinois (caractères + pinyin). Utilise le vocabulaire de la leçon.
\begin{enumerate}
\item Pour faciliter les démarches des habitants, la mairie offre les services suivants.
\item 1. Délivrer les cartes d'identité ; 2. Délivrer les actes de naissance ; 3. Recevoir les plaintes.
\item Horaires : du lundi au vendredi, de 8h à 17h.
\item Adresse : 10, avenue Kennedy. Téléphone : 119.
\item En cas de question, veuillez nous contacter. Merci de votre coopération !
\end{enumerate}
\end{exercice}

\begin{corrige}[Exercice Thème]
\begin{enumerate}
\item \zh{为了方便居民办事，市政府提供以下服务。} (Wèile fāngbiàn jūmín bànshì, shìzhèngfǔ tígōng yǐxià fúwù.) \emph{Note : « mairie » = \zh{市政府} (gouvernement municipal) ou \zh{区政府} (arrondissement). « Démarches » = \zh{办事}.}
\item \zh{一、办理身份证；二、办理出生证明；三、接受投诉。} (Yī, bànlǐ shēnfènzhèng ; èr, bànlǐ chūshēng zhèngmíng ; sān, jiēshòu tóusù.) \emph{Note : « plainte » = \zh{投诉} (tóusù).}
\item \zh{办公时间：星期一至星期五，上午八点到下午五点。} (Bàngōng shíjiān : xīngqīyī zhì xīngqīwǔ, shàngwǔ bā diǎn dào xiàwǔ wǔ diǎn.)
\item \zh{地址：肯尼迪大道十号。电话：一一九。} (Dìzhǐ : Kěnnídí dàdào shí hào. Diànhuà : yāo yāo jiǔ.) \emph{Note : 119 se lit \zh{幺幺九} (yāo yāo jiǔ).}
\item \zh{如有问题，请随时联系我们。谢谢合作！} (Rú yǒu wèntí, qǐng suíshí liánxì wǒmen. Xièxie hézuò !)
\end{enumerate}
\end{corrige}

\begin{exercice}[Version : Chinois → Français]
Traduis en français les phrases suivantes, extraites d'autres avis administratifs.
\begin{enumerate}
\item \zh{本单位为了提高服务质量，决定延长办公时间。}
\item \zh{办理业务，请携带身份证及相关材料。}
\item \zh{服务热线：12345，二十四小时受理。}
\item \zh{如有疑问，请登录官网查询或拨打热线。}
\item \zh{特此公告。}
\end{enumerate}
\end{exercice}

\begin{corrige}[Exercice Version]
\begin{enumerate}
\item « Notre unité / Le présent service, afin d'améliorer la qualité du service, a décidé de prolonger les horaires de bureau. » (\zh{本单位} = notre service / l'administration émettrice ; \zh{提高} = améliorer ; \zh{服务质量} = qualité du service ; \zh{决定} = décider ; \zh{延长} = prolonger).
\item « Pour effectuer les démarches, veuillez apporter votre carte d'identité et les documents concernés. » (\zh{办理业务} = traiter les affaires/démarches ; \zh{携带} = apporter/emporter ; \zh{及} = et (formel) ; \zh{相关材料} = documents/materials concernés).
\item « Ligne d'assistance : 12345, traitement 24h/24. » (\zh{服务热线} = ligne d'assistance/service ; \zh{二十四小时} = 24 heures ; \zh{受理} = traiter/recevoir une demande).
\item « En cas de doute/question, veuillez consulter le site officiel ou appeler la hotline. » (\zh{疑问} = doute/question ; \zh{登录} = se connecter/accéder ; \zh{官网} = site officiel ; \zh{查询} = consulter/rechercher ; \zh{拨打} = composer/appeler ; \zh{热线} = ligne directe/hotline).
\item « Spécialement par la présente, avis est donné. » (Formule finale très formelle, équivalent de « Fait pour servir et valoir ce que de droit. »)
\end{enumerate}
\end{corrige}

\courssub{Production guidée et évaluation}

\begin{methode}[Rédiger un avis de services (5 étapes)]
\begin{enumerate}
\item \textbf{Analyser la commande :} Quelle administration ? Quels services (3-4) ? Quels horaires/adresse/téléphone ? Contexte : Cameroun ou Chine ?
\item \textbf{Structurer le squelette :} Écrire les 5 blocs en français : Titre — Objet (\zh{为了...}) — Liste (\zh{一、二、三}) — Infos pratiques (\zh{时间/地址/电话}) — Clôture.
\item \textbf{Passer au chinois :} Traduire bloc par bloc en réutilisant les modèles (\zh{本中心提供...}, \zh{办理...}, \zh{接受...}, \zh{办公时间：...}).
\item \textbf{Vérifier les points de vigilance :} Chiffres chinois + \zh{、}, deux-points \zh{：}, \zh{幺} pour le 1 au téléphone, \zh{至} pour les horaires, \zh{谢谢合作！}.
\item \textbf{Relire à voix haute :} Vérifier les tons, la fluidité, l'ordre des mots (Sujet + 为了... + Verbe principal).
\end{enumerate}
\end{methode}

\begin{experience}[Production écrite : Avis pour une collectivité camerounaise]
\textbf{Consigne :} Tu es chargé de la communication à la \textbf{Mairie de Yaoundé VI} (ou \textbf{Conseil de Buea}). Rédige un avis \zh{公告} en chinois (caractères + pinyin) pour informer les usagers des services suivants :
\begin{itemize}
\item Délivrance des actes de naissance (\zh{出生证明}).
\item Établissement des cartes d'identité nationale (\zh{身份证}).
\item Réception des réclamations sur l'éclairage public et la voirie (\zh{路灯和道路维修投诉}).
\end{itemize}
Horaires : Lundi–Vendredi 7h30–15h30. Adresse : Carrefour Nlongkak (Yaoundé) / Molyko (Buea). Téléphone : 119 (urgence) ou 2220 0000 (standard).
\textbf{Contraintes :} Respecter les 5 blocs. Utiliser \zh{为了方便市民}... \zh{本单位}... \zh{一、二、三}. Clôture standard.
\end{experience}

\begin{experience}[Évaluation par les pairs : Grille de correction]
Échange ta production avec un camarade. Coche chaque critère réussi :
\begin{center}
\begin{tabularx}{\linewidth}{|X|c|}
\hline
\rowcolor{popblueL} Critère & Validé ? \\
\hline
Titre : \zh{公告} + Nom de l'administration & \square \\
\hline
Objet : \zh{为了方便市民/居民，本单位提供以下服务：} & \square \\
\hline
Liste : 3 items numérotés \zh{一、二、三} + \zh{、} + verbes \zh{办理/接受} + \zh{；} & \square \\
\hline
Infos pratiques : \zh{办公时间：} + \zh{至} + \zh{地址：} + \zh{电话：} (lecture \zh{幺}) & \square \\
\hline
Clôture : \zh{如有问题，请随时联系我们。谢谢合作！} & \square \\
\hline
Écriture : Caractères lisibles, traits dans l'ordre, pas de confusion \zh{公/翁}, \zh{证/正}, \zh{电/由} & \square \\
\hline
Pinyin : Tons corrects sur tout le texte & \square \\
\hline
\end{tabularx}
\end{center}
\textbf{Bilan :} 7/7 = Maîtrise parfaite ; 5-6/7 = Bon, revois les points manquants ; <5 = Refaire l'exercice avec le modèle.
\end{experience}

\begin{aretenir}[L'essentiel de la leçon 19 — Offre de services aux usagers]
\begin{itemize}
\item \textbf{Architecture de l'avis (\zh{公告}) :} 1. Titre (\zh{公告}) → 2. Objet (\zh{为了...}) → 3. Liste numérotée (\zh{一、二、三} + \zh{、}) → 4. Infos pratiques (\zh{时间/地址/电话} + \zh{：}) → 5. Clôture (\zh{谢谢合作！}).
\item \textbf{Structures clés :} \zh{为了} + but (antéposé) ; \zh{本} + nom (émetteur) ; \zh{办理} (délivrer) / \zh{接受} (recevoir) + objet ; \zh{如有} + nom, \zh{请} + verbe (invitation polie).
\item \textbf{Vocabulaire administratif :} \zh{市民/居民}, \zh{服务中心}, \zh{身份证}, \zh{出生证明}, \zh{户口本} (spécifique Chine), \zh{意见/建议/投诉}, \zh{办公时间}, \zh{联系方式}.
\item \textbf{Écriture :} Radicaux \zh{八} (\zh{公}), \zh{口} (\zh{告}), \zh{力} (\zh{办}), \zh{讠} (\zh{证}), \zh{日} (\zh{时}), \zh{田} (\zh{电}). Ordre des traits rigoureux.
\item \textbf{Conventions typographiques :} Chiffres chinois, point énumératif \zh{、}, deux-points \zh{：}, lecture téléphone \zh{幺} (1), \zh{至} pour « à/jusqu'à ».
\item \textbf{Compétence Bac :} Savoir produire un avis complet, lisible, politement formulé, en caractères + pinyin, sur un contexte camerounais ou chinois.
\end{itemize}
\end{aretenir}

\coursec{Lexique des services aux usagers}

Dans la section précédente, nous avons découvert un avis de services affiché à l'entrée d'une mairie : horaires d'ouverture, guichets, documents à présenter. Pour comprendre un tel avis — et surtout pour en rédiger un — il faut d'abord maîtriser un vocabulaire précis et fréquent. C'est l'objet de cette section : apprendre les mots du service aux usagers, non pas isolément, mais dans leurs combinaisons naturelles, celles que l'on retrouve dans tous les textes administratifs chinois et qui sont évaluées au Baccalauréat.

\courssub{Observation : les mots dans leur contexte}

Avant tout tableau, lisons ces phrases extraites d'avis authentiques de type municipal :

\begin{exemplebox}[Observer avant d'apprendre]
\zh{本窗口办理身份证。}\par
\emph{Běn chuāngkǒu bànlǐ shēnfènzhèng.}\par
« Ce guichet délivre les cartes d'identité. »

\zh{我们提供以下服务。}\par
\emph{Wǒmen tígōng yǐxià fúwù.}\par
« Nous offrons les services suivants. »

\zh{请带齐所有材料。}\par
\emph{Qǐng dài qí suǒyǒu cáiliào.}\par
« Apportez toutes les pièces (au complet). »

\zh{您需要提交申请表。}\par
\emph{Nín xūyào tíjiāo shēnqǐngbiǎo.}\par
« Vous devez déposer un formulaire de demande. »
\end{exemplebox}

Observons : chaque phrase contient un \cle{verbe de démarche} (\zh{办理}, \zh{提供}, \zh{提交}) suivi d'un \cle{nom d'objet administratif} (\zh{身份证}, \zh{服务}, \zh{材料}, \zh{申请表}). C'est exactement cette architecture « verbe + nom » qu'il faut mémoriser.

\courssub{Glossaire principal : les verbes et noms de la démarche}

\begin{center}
\begin{tabularx}{\linewidth}{|X|X|X|X|}
\hline
\rowcolor{popblueL} Caractères & Pinyin & Nature & Traduction \\
\hline
服务 & fúwù & n. / v. & service ; servir \\
\hline
办理 & bànlǐ & v. & traiter, effectuer (une démarche) \\
\hline
提供 & tígōng & v. & fournir, offrir \\
\hline
申请 & shēnqǐng & v. / n. & demander ; demande \\
\hline
材料 & cáiliào & n. & documents, pièces à fournir \\
\hline
窗口 & chuāngkǒu & n. & guichet (litt. « fenêtre ») \\
\hline
证明 & zhèngmíng & n. / v. & attestation ; prouver \\
\hline
居民 & jūmín & n. & habitant, résident \\
\hline
意见 & yìjian & n. & avis, opinion, remarque \\
\hline
建议 & jiànyì & n. / v. & suggestion ; suggérer \\
\hline
\end{tabularx}
\end{center}

\begin{center}
\begin{tabularx}{\linewidth}{|X|X|X|X|}
\hline
\rowcolor{popblueL} Caractères & Pinyin & Nature & Traduction \\
\hline
办公 & bàngōng & v. & travailler (au bureau), assurer le service \\
\hline
时间 & shíjiān & n. & temps, horaires \\
\hline
地址 & dìzhǐ & n. & adresse \\
\hline
联系 & liánxì & v. & contacter, se mettre en relation avec \\
\hline
方便 & fāngbiàn & adj. & pratique, commode \\
\hline
随时 & suíshí & adv. & à tout moment \\
\hline
合作 & hézuò & v. / n. & coopérer ; coopération \\
\hline
\end{tabularx}
\end{center}

\begin{definition}[\zh{办理} bànlǐ : le verbe des démarches]
\zh{办理} signifie « traiter, effectuer, accomplir » et s'emploie \textbf{uniquement pour les démarches administratives ou officielles} : sa structure est \cle{办理 + document / démarche}.
\begin{itemize}
\item \zh{办理身份证} \emph{bànlǐ shēnfènzhèng} : délivrer / faire établir une carte d'identité
\item \zh{办理护照} \emph{bànlǐ hùzhào} : faire une demande de passeport
\item \zh{办理手续} \emph{bànlǐ shǒuxù} : accomplir les formalités
\end{itemize}
On ne dit \textbf{jamais} \zh{办理饭} ou \zh{办理作业} : pour « faire » en général, le chinois utilise \zh{做} (\emph{zuò}). \zh{办理} appartient au registre administratif ; \zh{做} appartient à la vie quotidienne.
\end{definition}

\begin{attention}[\zh{意见} ne veut pas dire « opinion favorable »]
Les francophones traduisent souvent \zh{意见} (\emph{yìjian}) par « opinion » au sens large. En réalité, dans le contexte administratif, \zh{意见} désigne un \textbf{avis, une remarque, souvent une critique} :
\begin{itemize}
\item \zh{提意见} \emph{tí yìjian} : formuler une remarque (souvent pour critiquer ou améliorer)
\item \zh{接受意见} \emph{jiēshòu yìjian} : recevoir / prendre en compte les avis
\item \zh{征求意见} \emph{zhēngqiú yìjian} : recueillir les avis (des usagers)
\end{itemize}
Pour « suggestion » constructive, on préfère \zh{建议} (\emph{jiànyì}), qui est plus positif : \zh{提建议} « faire une suggestion ». Dans un avis de mairie, \zh{欢迎居民提出意见和建议} signifie « les habitants sont invités à formuler avis et suggestions ».
\end{attention}

\courssub{Le champ lexical des documents}

Tout avis de services mentionne les pièces à fournir. Voici les quatre documents essentiels à connaître :

\begin{center}
\begin{tabularx}{\linewidth}{|X|X|X|X|}
\hline
\rowcolor{popblueL} Caractères & Pinyin & Nature & Traduction \\
\hline
身份证 & shēnfènzhèng & n. & carte d'identité \\
\hline
出生证明 & chūshēng zhèngmíng & n. & acte de naissance \\
\hline
护照 & hùzhào & n. & passeport \\
\hline
申请表 & shēnqǐngbiǎo & n. & formulaire de demande \\
\hline
\end{tabularx}
\end{center}

Remarquons la logique de construction des mots chinois : \zh{身份} « identité » + \zh{证} « certificat » donnent \zh{身份证} ; \zh{申请} « demande » + \zh{表} « formulaire » donnent \zh{申请表} ; \zh{出生} « naître » + \zh{证明} « attestation » donnent \zh{出生证明}. Comprendre les morceaux permet de deviner les mots nouveaux : c'est une stratégie précieuse en compréhension de l'écrit.

\courssub{Les collocations clés : verbe + nom}

\begin{methode}[Mémoriser par collocations, pas mot par mot]
Au Baccalauréat, ce n'est pas le mot isolé qui est évalué, mais la \cle{combinaison correcte verbe + nom}. Procédez ainsi :
\begin{enumerate}
\item Pour chaque verbe (\zh{办理}, \zh{提供}, \zh{提交}, \zh{接受}, \zh{联系}), apprenez son ou ses noms partenaires par cœur.
\item Fabriquez une phrase complète avec chaque collocation et traduisez-la.
\item En thème (français $\rightarrow$ chinois), demandez-vous toujours : « Quel verbe ce nom exige-t-il ? » On « délivre » un document (\zh{办理}), on « dépose » des pièces (\zh{提交}), on « offre » un service (\zh{提供}).
\end{enumerate}
\end{methode}

\begin{center}
\begin{tabularx}{\linewidth}{|X|X|X|}
\hline
\rowcolor{popgreenL} Collocation & Pinyin & Traduction \\
\hline
办理证件 & bànlǐ zhèngjiàn & délivrer / établir un document officiel \\
\hline
提供服务 & tígōng fúwù & offrir un service \\
\hline
提交材料 & tíjiāo cáiliào & déposer les pièces (le dossier) \\
\hline
接受意见 & jiēshòu yìjian & recevoir les avis (des usagers) \\
\hline
联系工作人员 & liánxì gōngzuò rényuán & contacter le personnel \\
\hline
\end{tabularx}
\end{center}

\begin{exemplebox}[Les collocations en phrases]
\zh{市政府为居民提供多种服务。}\par
\emph{Shì zhèngfǔ wèi jūmín tígōng duō zhǒng fúwù.}\par
« La mairie offre de nombreux services aux habitants. »

\zh{请先提交材料，再办理证件。}\par
\emph{Qǐng xiān tíjiāo cáiliào, zài bànlǐ zhèngjiàn.}\par
« Veuillez d'abord déposer les pièces, puis faire établir le document. »

\zh{如有问题，请随时联系工作人员。}\par
\emph{Rú yǒu wèntí, qǐng suíshí liánxì gōngzuò rényuán.}\par
« Pour toute question, contactez le personnel à tout moment. »

\zh{网上办公更方便。}\par
\emph{Wǎngshang bàngōng gèng fāngbiàn.}\par
« Les démarches en ligne sont plus pratiques. »
\end{exemplebox}

\begin{popfigure}
\begin{tikzpicture}[
  box/.style={rectangle, rounded corners=3pt, draw=popblue, fill=popblueL, align=center, minimum height=0.9cm, inner sep=4pt, font=\small},
  arr/.style={-{Stealth[length=2.5mm]}, thick, popdark}
]
\node[box, fill=poporangeL, draw=poporange, align=center] (u) at (0,0){居民\\ \emph{jūmín}\\ usager};
\node[box, align=center] (g) at (3.4,0){窗口\\ \emph{chuāngkǒu}\\ guichet};
\node[box, align=center] (m) at (6.8,0){提交材料\\ \emph{tíjiāo cáiliào}\\ déposer les pièces};
\node[box, align=center] (b) at (10.2,0){办理\\ \emph{bànlǐ}\\ traitement};
\node[box, fill=popgreenL, draw=popgreen, align=center] (z) at (13.4,0){领取证明\\ \emph{lǐngqǔ zhèngmíng}\\ retirer l'attestation};
\draw[arr] (u) -- (g);
\draw[arr] (g) -- (m);
\draw[arr] (m) -- (b);
\draw[arr] (b) -- (z);
\end{tikzpicture}
\legende{Le parcours de l'usager : du guichet au retrait de l'attestation.}
\end{popfigure}

Ce schéma résume la logique de tout texte d'offre de services : on y retrouve, dans l'ordre, les cinq collocations du tableau. Rédiger un avis, c'est souvent décrire ce parcours avec des connecteurs (\zh{先} « d'abord », \zh{然后} « ensuite », \zh{最后} « enfin »), que nous étudierons dans la section 4.

\begin{popfigure}
\begin{center}\tcbox[colback=popink,colframe=popink,coltext=white,arc=9pt,boxsep=4pt,left=6pt,right=6pt]{\bfseries\small 服务 \emph{fúwù} services}\end{center}
\vspace{-2pt}
\begin{tcbraster}[raster columns=2,raster equal height=rows,raster column skip=6pt,raster row skip=6pt,enhanced,blankest]
\begin{tcolorbox}[enhanced,colback=white,colframe=poporange,boxrule=0.9pt,arc=3pt,title={verbes},fonttitle=\bfseries\scriptsize,coltitle=white,colbacktitle=poporange,left=4pt,right=4pt,top=2pt,bottom=2pt,boxsep=1pt,toptitle=1.5pt,bottomtitle=1.5pt,halign=left]
\begin{itemize}[nosep,leftmargin=*,itemsep=1pt,font=\scriptsize]
\item 办理
\item 提供
\item 提交
\end{itemize}
\end{tcolorbox}
\begin{tcolorbox}[enhanced,colback=white,colframe=popgreen,boxrule=0.9pt,arc=3pt,title={documents},fonttitle=\bfseries\scriptsize,coltitle=white,colbacktitle=popgreen,left=4pt,right=4pt,top=2pt,bottom=2pt,boxsep=1pt,toptitle=1.5pt,bottomtitle=1.5pt,halign=left]
\begin{itemize}[nosep,leftmargin=*,itemsep=1pt,font=\scriptsize]
\item 身份证
\item 护照
\item 申请表
\end{itemize}
\end{tcolorbox}
\begin{tcolorbox}[enhanced,colback=white,colframe=poppurple,boxrule=0.9pt,arc=3pt,title={lieu / contact},fonttitle=\bfseries\scriptsize,coltitle=white,colbacktitle=poppurple,left=4pt,right=4pt,top=2pt,bottom=2pt,boxsep=1pt,toptitle=1.5pt,bottomtitle=1.5pt,halign=left]
\begin{itemize}[nosep,leftmargin=*,itemsep=1pt,font=\scriptsize]
\item 窗口
\item 地址
\item 联系
\end{itemize}
\end{tcolorbox}
\begin{tcolorbox}[enhanced,colback=white,colframe=popgold,boxrule=0.9pt,arc=3pt,title={avis},fonttitle=\bfseries\scriptsize,coltitle=white,colbacktitle=popgold,left=4pt,right=4pt,top=2pt,bottom=2pt,boxsep=1pt,toptitle=1.5pt,bottomtitle=1.5pt,halign=left]
\begin{itemize}[nosep,leftmargin=*,itemsep=1pt,font=\scriptsize]
\item 意见
\item 建议
\end{itemize}
\end{tcolorbox}
\end{tcbraster}

\legende{Carte mentale du champ lexical : quatre familles autour du mot 服务.}
\end{popfigure}

\courssub{Texte d'application : un avis de la mairie de Yaoundé}

\begin{textebox}[Avis aux usagers — Mairie de Yaoundé]
\zh{雅温得市政府通知}\\[4pt]
\zh{为了方便居民，本市政府提供以下服务：办理身份证、出生证明和护照。}\\[4pt]
\zh{办公时间：星期一至星期五，上午八点到下午三点。}\\[4pt]
\zh{请带齐所有材料，先到三号窗口提交申请表，然后办理证件。}\\[4pt]
\zh{如有问题，请随时联系工作人员。我们也接受居民的意见和建议。}\\[4pt]
\zh{地址：雅温得市中心区政府大楼。谢谢合作！}\\[8pt]
\emph{Yǎwēndé shì zhèngfǔ tōngzhī}\\
\emph{Wèile fāngbiàn jūmín, běn shì zhèngfǔ tígōng yǐxià fúwù : bànlǐ shēnfènzhèng, chūshēng zhèngmíng hé hùzhào.}\\
\emph{Bàngōng shíjiān : xīngqīyī zhì xīngqīwǔ, shàngwǔ bā diǎn dào xiàwǔ sān diǎn.}\\
\emph{Qǐng dài qí suǒyǒu cáiliào, xiān dào sān hào chuāngkǒu tíjiāo shēnqǐngbiǎo, ránhòu bànlǐ zhèngjiàn.}\\
\emph{Rú yǒu wèntí, qǐng suíshí liánxì gōngzuò rényuán. Wǒmen yě jiēshòu jūmín de yìjian hé jiànyì.}\\
\emph{Dìzhǐ : Yǎwēndé shì zhōngxīn qū zhèngfǔ dàlóu. Xièxie hézuò !}\\[8pt]
« Avis de la mairie de Yaoundé. Pour faciliter la vie des habitants, la mairie offre les services suivants : établissement des cartes d'identité, des actes de naissance et des passeports. Horaires : du lundi au vendredi, de 8 h à 15 h. Apportez toutes les pièces, déposez d'abord le formulaire au guichet n° 3, puis faites établir votre document. Pour toute question, contactez le personnel à tout moment. Nous recevons aussi les avis et suggestions des habitants. Adresse : bâtiment administratif du centre-ville de Yaoundé. Merci de votre coopération ! »
\end{textebox}

Questions de compréhension : (1) Quels services la mairie offre-t-elle ? (2) Quels sont les horaires d'ouverture ? (3) Que doit faire l'usager en arrivant ? (4) Où se trouve le service ?

\courssub{Exercice résolu et entraînement}

\begin{exoresolu}[Thème : traduire en chinois]
Traduisez : « Ce guichet délivre les cartes d'identité. Apportez toutes les pièces. »
\tcblower
\textbf{Méthode.} On identifie d'abord les collocations : « délivrer une carte d'identité » $\rightarrow$ \zh{办理身份证} ; « apporter toutes les pièces » $\rightarrow$ \zh{带齐所有材料}. On place ensuite le sujet et le démonstratif : « ce guichet » $\rightarrow$ \zh{这个窗口}.

\textbf{Traduction.} \zh{这个窗口办理身份证。请带齐所有材料。}

\emph{Zhège chuāngkǒu bànlǐ shēnfènzhèng. Qǐng dài qí suǒyǒu cáiliào.}

\textbf{Piège évité.} Ne pas traduire « délivrer » par \zh{做} : on dirait \zh{做身份证}, ce qui signifierait « fabriquer une fausse carte » ! Le verbe administratif correct est \zh{办理}.
\end{exoresolu}

\begin{exercice}[À vous de jouer]
\begin{enumerate}
\item Traduisez en français : \zh{我们随时接受居民的意见。}
\item Traduisez en chinois : « Vous devez d'abord déposer le formulaire de demande. »
\item Complétez avec le bon verbe (\zh{办理} / \zh{提供} / \zh{提交}) : \zh{市政府为居民\_\_\_\_\_\_多种服务。}
\end{enumerate}
\end{exercice}

\begin{corrige}[Exercice « À vous de jouer »]
\begin{enumerate}
\item « Nous recevons à tout moment les avis des habitants. » (\zh{随时} = à tout moment ; \zh{接受意见} = recevoir les avis.)
\item \zh{您需要先提交申请表。} \emph{Nín xūyào xiān tíjiāo shēnqǐngbiǎo.} (« déposer » $\rightarrow$ \zh{提交} ; « d'abord » $\rightarrow$ \zh{先}, placé avant le verbe.)
\item \zh{提供} : la collocation correcte est \zh{提供服务} « offrir des services ». \zh{办理} exige un document, \zh{提交} exige des pièces ou un formulaire.
\end{enumerate}
\end{corrige}

\begin{savaistu}[Du guichet à l'écran : les services en ligne en Chine]
\zh{窗口} signifie littéralement « fenêtre » : c'est la petite ouverture du guichet derrière laquelle l'employé reçoit les dossiers. Mais la Chine est devenue l'un des pays les plus numérisés au monde : de nombreuses démarches se font désormais sur téléphone, et les administrations parlent de \zh{网上服务} (\emph{wǎngshang fúwù}), les « services en ligne », et de \zh{网上办公} « démarches en ligne ». Le mot \zh{窗口} a même pris un sens figuré : une « fenêtre » sur le monde, une vitrine. Au Cameroun aussi, certains services publics proposent désormais des démarches en ligne : un point commun de modernisation entre les deux pays.
\end{savaistu}

\begin{aretenir}[L'essentiel du lexique]
\begin{itemize}
\item \zh{办理} + document : le verbe des démarches officielles (jamais \zh{做}).
\item \zh{提供服务} « offrir un service », \zh{提交材料} « déposer les pièces », \zh{接受意见} « recevoir les avis », \zh{联系工作人员} « contacter le personnel ».
\item \zh{意见} = avis, remarque (souvent critique) ; \zh{建议} = suggestion positive.
\item Documents : \zh{身份证}, \zh{出生证明}, \zh{护照}, \zh{申请表}.
\item Le parcours type : \zh{窗口} $\rightarrow$ \zh{提交材料} $\rightarrow$ \zh{办理} $\rightarrow$ \zh{领取证明}.
\end{itemize}
\end{aretenir}

Ce lexique maîtrisé, nous pouvons passer à la section suivante : l'écriture correcte des caractères du thème, en particulier ceux de \zh{办理}, \zh{窗} et \zh{证}, qui cachent quelques pièges graphiques.

\coursec{Écriture des caractères du thème}

Dans la section précédente, nous avons découvert le lexique des services aux usagers : \zh{办理} (bànlǐ, « effectuer une démarche »), \zh{服务} (fúwù, « service »), \zh{证明} (zhèngmíng, « certificat »), \zh{窗口} (chuāngkǒu, « guichet »), \zh{材料} (cáiliào, « documents »). Savoir prononcer et comprendre ces mots ne suffit pas : à l'écrit, l'administration chinoise exige des caractères tracés correctement. Un trait mal placé peut transformer \zh{办} en \zh{为} et changer complètement le sens d'une phrase. Cette section vous apprend à écrire les sept caractères clés du thème, à reconnaître leurs radicaux et à ne pas les confondre avec leurs « sosies ».

\courssub{Observer d'abord : les mots dans l'avis de la mairie}

Relisons quelques phrases de l'avis de services étudié en début de leçon :

\begin{exemplebox}[Les caractères du thème en contexte]
\zh{为了方便居民，本市政府开设了新的服务窗口。}\\
\emph{Wèile fāngbiàn jūmín, běn shì zhèngfǔ kāishè le xīn de fúwù chuāngkǒu.}\\
« Pour faciliter la vie des habitants, la mairie a ouvert de nouveaux guichets de service. »

\zh{办理各种证明，请带好所有材料。}\\
\emph{Bànlǐ gè zhǒng zhèngmíng, qǐng dài hǎo suǒyǒu cáiliào.}\\
« Pour effectuer vos demandes de certificats, veuillez apporter tous vos documents. »
\end{exemplebox}

Observez : chaque caractère est un petit dessin organisé. Pour l'écrire correctement, il faut connaître trois choses : son \cle{radical} (la clé qui donne une indication de sens), son \cle{nombre de traits} et l'\cle{ordre des traits}.

\courssub{La règle générale d'ordre des traits}

\begin{propriete}[Les quatre règles d'or de l'ordre des traits]
\begin{enumerate}
\item \textbf{Haut avant bas} : on écrit la partie supérieure d'abord. Ex. \zh{务} : d'abord \zh{夂} en haut, puis \zh{力} en bas.
\item \textbf{Gauche avant droite} : on écrit la composante de gauche d'abord. Ex. \zh{请} : d'abord \zh{讠} à gauche, puis \zh{青} à droite ; \zh{料} : d'abord \zh{米}, puis \zh{斗}.
\item \textbf{Extérieur avant intérieur} : pour les caractères à enclos, on trace l'extérieur, puis l'intérieur. Ex. \zh{间} : d'abord la porte \zh{门}, puis le soleil \zh{日} à l'intérieur.
\item \textbf{Fermeture en dernier} : le trait qui ferme l'enclos se trace en tout dernier. Ex. \zh{国} : le trait horizontal du bas ferme le cadre à la fin ; dans \zh{门}, le cadre est tracé avant le contenu.
\end{enumerate}
\end{propriete}

\begin{attention}[Pourquoi l'ordre des traits est-il si important ?]
Contrairement au français, où l'ordre des lettres suffit, le chinois exige un ordre de tracé précis. Un caractère écrit « dans le désordre » est souvent illisible, déséquilibré, et les applications de reconnaissance d'écriture (dictionnaires électroniques, claviers manuscrits) ne le reconnaissent pas. À l'examen, une écriture fautive est pénalisée. Entraînez-vous toujours en comptant les traits à voix haute.
\end{attention}

\courssub{Étude caractère par caractère}

\textbf{1. 办 (bàn) — « faire, gérer »}\par
Radical : \zh{力} (lì, la force). 5 traits. Ordre : 1. le grand crochet horizontal-vertical \zh{(trait brisé)} (横折钩) en haut à gauche ; 2. le point de gauche ; 3. le point de droite ; 4. l'horizontal du radical \zh{力} ; 5. le crochet vertical du radical \zh{力}. L'idée : la force \zh{力} entourée de deux gouttes de sueur — travailler avec effort, donc « s'occuper de ». On le trouve dans \zh{办理} (bànlǐ, effectuer une démarche), \zh{办公室} (bàngōngshì, bureau).

\textbf{2. 务 (wù) — « affaire, tâche »}\par
Radical : \zh{力} placé en bas. 5 traits. On écrit d'abord la partie haute \zh{夂} (zhǐ, 3 traits : un trait oblique \emph{piě}, un crochet horizontal-vertical, un point final), puis \zh{力} en bas (2 traits : horizontal, crochet vertical). Attention à ne pas confondre la partie haute avec \zh{各} (gè, « chaque »), qui possède un \zh{口} supplémentaire en bas. On le trouve dans \zh{服务} (fúwù, service), \zh{业务} (yèwù, activité professionnelle).

\textbf{3. 请 (qǐng) — « prier, inviter, s'il vous plaît »}\par
Radical : \zh{讠} (la parole, forme simplifiée de \zh{言}). 10 traits. D'abord \zh{讠} à gauche (2 traits : le point, puis le pli horizontal-vertical), puis \zh{青} à droite (8 traits : le trait oblique, les trois horizontaux et le vertical du haut, puis \zh{月} en bas : oblique, pli, horizontal, horizontal). C'est le caractère de la politesse administrative : \zh{请带好材料} (qǐng dài hǎo cáiliào, « veuillez apporter vos documents »).

\textbf{4. 证 (zhèng) — « certificat, prouver »}\par
Radical : \zh{讠} (parole). 7 traits. D'abord \zh{讠} (2 traits), puis \zh{正} (5 traits : l'horizontal du haut, le vertical, le petit horizontal, le vertical gauche, l'horizontal de fermeture). Sens : une parole \zh{讠} droite et correcte \zh{正} — donc un document qui atteste la vérité. On le trouve dans \zh{证明} (zhèngmíng, certificat), \zh{身份证} (shēnfènzhèng, carte d'identité). Ne pas confondre avec \zh{政} (zhèng, politique), qui prend \zh{攵} à droite au lieu de \zh{讠} à gauche.

\textbf{5. 窗 (chuāng) — « fenêtre »}\par
Radical : \zh{穴} (xué, la grotte, le toit percé). 12 traits. On écrit \textbf{toujours le \zh{穴} d'abord} (5 traits : le point, le point vertical, le crochet, puis les deux petits traits intérieurs), puis la partie basse \zh{囱} (7 traits : vertical, pli, horizontal, horizontal, vertical, pli, horizontal de fermeture). On le trouve dans \zh{窗口} (chuāngkǒu, guichet — littéralement « fenêtre-bouche »), \zh{窗户} (chuānghu, fenêtre).

\textbf{6. 料 (liào) — « matériau, document »}\par
Radical : \zh{斗} (dǒu, la mesure à grain). 10 traits. D'abord \zh{米} (mǐ, riz, 6 traits) à gauche, puis \zh{斗} (4 traits : les deux points, l'horizontal, le vertical) à droite. Image : on mesure \zh{斗} le riz \zh{米} — donc la matière première. Ne pas confondre avec \zh{科} (kē, science, branche), qui a \zh{禾} (céréale) à gauche au lieu de \zh{米}. On le trouve dans \zh{材料} (cáiliào, documents, matériaux).

\textbf{7. 间 (jiān) — « entre, pièce, intervalle »}\par
Radical : \zh{门} (mén, la porte). 7 traits. Règle « extérieur avant intérieur » : on trace d'abord la porte \zh{门} (3 traits : le point, le vertical, le crochet horizontal-vertical), puis le soleil \zh{日} (4 traits : vertical, pli, horizontal, horizontal) à l'intérieur. Image : le soleil \zh{日} vu à travers la porte \zh{门} — l'espace entre les deux battants. Ne pas confondre avec \zh{问} (wèn, demander), qui contient une bouche \zh{口} au lieu du soleil \zh{日}. On le trouve dans \zh{时间} (shíjiān, le temps), \zh{房间} (fángjiān, la pièce).

\courssub{Tableau récapitulatif des caractères du thème}

\begin{center}
\begin{tabularx}{\linewidth}{|X|X|X|X|}\hline
\rowcolor{popblueL} Caractère & Pinyin & Radical / traits & Caractère proche \\ \hline
\zh{办} & bàn & \zh{力} / 5 & \zh{为} (wèi, pour) \\ \hline
\zh{务} & wù & \zh{力} (bas) / 5 & \zh{各} (gè, chaque) \\ \hline
\zh{请} & qǐng & \zh{讠} / 10 & \zh{清} (qīng), \zh{情} (qíng) \\ \hline
\zh{证} & zhèng & \zh{讠} / 7 & \zh{政} (zhèng, politique) \\ \hline
\zh{窗} & chuāng & \zh{穴} / 12 & --- \\ \hline
\zh{料} & liào & \zh{斗} / 10 & \zh{科} (kē, science) \\ \hline
\zh{间} & jiān & \zh{门} / 7 & \zh{问} (wèn, demander) \\ \hline
\end{tabularx}
\end{center}

\courssub{Le radical 讠, indice de sens}

\begin{propriete}[Le radical de la parole 讠]
Le radical \zh{讠}, forme simplifiée de \zh{言} (yán, « parole »), se place à gauche du caractère et indique un lien avec le langage, la parole ou la communication. Exemples : \zh{请} (qǐng, prier — une parole polie), \zh{证} (zhèng, certifier — une parole qui atteste), \zh{议} (yì, discuter), \zh{话} (huà, parole), \zh{说} (shuō, parler), \zh{谢} (xiè, remercier). \textbf{Repérer le radical aide à deviner le sens d'un caractère inconnu} : si vous voyez \zh{讠} à gauche, pensez « parole ».
\end{propriete}

\begin{methode}[Distinguer 请 / 清 / 情 grâce au radical]
Ces trois caractères partagent la même partie phonétique \zh{青} (qīng), qui donne le son. Seul le radical change et donne le sens :
\begin{enumerate}
\item \zh{请} = \zh{讠} (parole) + \zh{青} : on \emph{parle} pour prier quelqu'un → « prier, inviter » (qǐng).
\item \zh{清} = \zh{氵} (eau) + \zh{青} : de l'\emph{eau} pure → « clair, limpide » (qīng).
\item \zh{情} = \zh{忄} (cœur) + \zh{青} : ce qui vient du \emph{cœur} → « sentiment » (qíng).
\end{enumerate}
Astuce de mémorisation : « On \textbf{parle} pour prier, l'\textbf{eau} est claire, le \textbf{cœur} sent. »
\end{methode}

\begin{popfigure}
\begin{tikzpicture}[font=\small]
\node[draw=popblue, fill=popblueL, rounded corners, minimum width=1.4cm, minimum height=1.2cm, align=center] (qing) at (0,0) {\zh{请}\\ qǐng};
\node[draw=popgreen, fill=popgreenL, rounded corners, minimum width=1.4cm, minimum height=1.2cm, align=center] (yan) at (-3,0) {\zh{讠}\\ parole};
\node[draw=poporange, fill=poporangeL, rounded corners, minimum width=1.4cm, minimum height=1.2cm, align=center] (qing2) at (3,0) {\zh{青}\\ qīng (son)};
\draw[-{Stealth}, thick, popgreen] (yan) -- node[above, align=center]{1. sens} (qing);
\draw[-{Stealth}, thick, poporange] (qing2) -- node[above, align=center]{2. son} (qing);
\node[align=center, popmuted] at (0,-1.4) {radical à gauche + phonétique à droite};
\end{tikzpicture}
\legende{Décomposition du caractère 请 : le radical 讠 donne le sens (parole), la partie 青 donne la prononciation.}
\end{popfigure}

\courssub{Les caractères « sosies » : ne pas se tromper}

\begin{popfigure}
\begin{tikzpicture}[font=\small]
\node[draw=popdark, fill=popcream, rounded corners, align=center, minimum width=2.2cm, minimum height=1.1cm] at (-4,2) {\zh{\textcolor{popink}{办}} bàn\\ gérer};
\node[draw=popdark, fill=popcream, rounded corners, align=center, minimum width=2.2cm, minimum height=1.1cm] at (-1,2) {\zh{\textcolor{popink}{为}} wèi\\ pour};
\node[draw=popdark, fill=popcream, rounded corners, align=center, minimum width=2.2cm, minimum height=1.1cm] at (2,2) {\zh{间} jiān\\ \zh{门}+\textcolor{popink}{\zh{日}}};
\node[draw=popdark, fill=popcream, rounded corners, align=center, minimum width=2.2cm, minimum height=1.1cm] at (5,2) {\zh{问} wèn\\ \zh{门}+\textcolor{popink}{\zh{口}}};
\node[draw=popdark, fill=popcream, rounded corners, align=center, minimum width=2.2cm, minimum height=1.1cm] at (-4,0) {\zh{请} qǐng\\ \textcolor{popink}{\zh{讠}}+\zh{青}};
\node[draw=popdark, fill=popcream, rounded corners, align=center, minimum width=2.2cm, minimum height=1.1cm] at (-1,0) {\zh{清} qīng\\ \textcolor{popink}{\zh{氵}}+\zh{青}};
\node[draw=popdark, fill=popcream, rounded corners, align=center, minimum width=2.2cm, minimum height=1.1cm] at (2,0) {\zh{情} qíng\\ \textcolor{popink}{\zh{忄}}+\zh{青}};
\node[align=center, popmuted] at (0.5,-1.4) {L'élément surligné est celui qui change le sens.};
\end{tikzpicture}
\legende{Paires et familles de caractères proches : le radical différenciant est surligné.}
\end{popfigure}

\begin{attention}[Erreur fréquente des francophones : 为 au lieu de 办]
Les caractères \zh{办} (bàn, gérer) et \zh{为} (wèi, pour) se ressemblent beaucoup : mêmes points latéraux, mais un corps central différent. Écrire \zh{为理} au lieu de \zh{办理} est une faute grave, car le mot n'existe pas. Retenez : \zh{为} signifie « pour, pour le compte de » et se combine avec \zh{了} dans la structure \zh{为} + \zh{了} + verbe, comme dans \zh{为了方便居民} (wèile fāngbiàn jūmín, « pour faciliter la vie des habitants »). \zh{办}, lui, est un verbe d'action : \zh{办理证明} (bànlǐ zhèngmíng, « effectuer une demande de certificat »).
\end{attention}

\begin{exemplebox}[Comparer pour mémoriser]
\zh{为了方便居民，请到这里办理手续。}\\
\emph{Wèile fāngbiàn jūmín, qǐng dào zhèlǐ bànlǐ shǒuxù.}\\
« Pour faciliter la vie des habitants, veuillez venir ici effectuer vos démarches. »

\zh{请问，办理身份证需要什么材料？}\\
\emph{Qǐngwèn, bànlǐ shēnfènzhèng xūyào shénme cáiliào ?}\\
« S'il vous plaît, quels documents faut-il pour demander une carte d'identité ? »

\zh{三号窗口办理出生证明。}\\
\emph{Sān hào chuāngkǒu bànlǐ chūshēng zhèngmíng.}\\
« Le guichet numéro 3 délivre les actes de naissance. »
\end{exemplebox}

\courssub{Exercice résolu et dictée d'entraînement}

\begin{exoresolu}[Identifier le radical et corriger la faute]
Énoncé : dans chaque phrase, un caractère est mal écrit. Retrouvez la faute, écrivez le caractère correct et justifiez avec le radical.\\
a) \zh{请到三号窗日办理。} \quad b) \zh{为理证明需要材料。} \quad c) \zh{这个问间很大。}
\tcblower
Solution détaillée :\\
a) \zh{窗日} doit s'écrire \zh{窗口} (chuāngkǒu, guichet). La faute vient de la confusion entre \zh{日} (soleil) et \zh{口} (bouche, ouverture) : un guichet est une « fenêtre-ouverture », donc \zh{口}.\\
b) \zh{为理} doit s'écrire \zh{办理} (bànlǐ). \zh{为} (wèi, « pour ») ne se combine pas avec \zh{理} ; le verbe « effectuer une démarche » est \zh{办}, radical \zh{力} (la force, l'effort).\\
c) \zh{问间} doit s'écrire \zh{房间} (fángjiān, pièce). \zh{问} (wèn, demander) contient une bouche \zh{口} ; « pièce, espace » s'écrit \zh{间} avec le soleil \zh{日} à l'intérieur de la porte \zh{门}.
\end{exoresolu}

\begin{textebox}[Dictée guidée : à vous d'écrire]
Votre professeur dicte les cinq mots suivants. Écrivez-les en respectant l'ordre des traits, comptez les traits, puis auto-corrigez-vous avec le tableau récapitulatif de la section.\\
\zh{办理} (bànlǐ) — « effectuer une démarche »\\
\zh{服务} (fúwù) — « service »\\
\zh{证明} (zhèngmíng) — « certificat »\\
\zh{窗口} (chuāngkǒu) — « guichet »\\
\zh{材料} (cáiliào) — « documents »\\
Vérification : \zh{办} = 5 traits, \zh{务} = 5, \zh{证} = 7, \zh{窗} = 12, \zh{料} = 10. Avez-vous bien écrit \zh{穴} avant le bas de \zh{窗} ? Avez-vous mis \zh{米} et non \zh{禾} dans \zh{料} ? Avez-vous respecté l'ordre : crochet, points, puis radical \zh{力} pour \zh{办} ?
\end{textebox}

\begin{aretenir}[L'essentiel de la section]
\begin{itemize}
\item Quatre règles d'ordre des traits : haut avant bas, gauche avant droite, extérieur avant intérieur, fermeture en dernier.
\item Le radical \zh{讠} (parole) signale le sens : \zh{请}, \zh{证}, \zh{议}, \zh{话}.
\item La phonétique \zh{青} donne le son ; le radical choisit le sens : \zh{请} (parler), \zh{清} (eau), \zh{情} (cœur).
\item Ne jamais confondre \zh{办} (bàn, gérer, 5 traits, radical \zh{力}) et \zh{为} (wèi, pour), ni \zh{间} (\zh{门}+\zh{日}) et \zh{问} (\zh{门}+\zh{口}).
\end{itemize}
\end{aretenir}

\coursec{Structure et connecteurs de l'offre de services}

Dans la section précédente, nous avons appris à écrire correctement les caractères du thème administratif : \zh{服务} (fúwù, « service »), \zh{居民} (jūmín, « habitant »), \zh{申请} (shēnqǐng, « demande »), \zh{材料} (cáiliào, « pièces, documents »). Nous savons désormais tracer les traits et reconnaître les radicaux. Mais une offre de services ne se résume pas à des mots bien écrits : elle obéit à une \cle{architecture fixe} et à des \cle{connecteurs codifiés}. C'est précisément ce que nous allons étudier maintenant, afin de pouvoir rédiger un avis administratif complet, clair et poli, comme le ferait une mairie de Yaoundé ou de Douala.

\courssub{Observation : lire un avis de services avant de le construire}

Lisons d'abord un avis complet, tel qu'on pourrait l'afficher à l'entrée d'un centre de services municipal. Observez l'ordre des parties : nous le décortiquerons ensuite.

\begin{textebox}[Avis affiché au centre de services du quartier Nkolndongo, Yaoundé]
服务公告\\
Fúwù gōnggào\\
Avis de services\\[0.5em]
为了方便居民办理各种手续，本中心提供以下服务。\\
Wèile fāngbiàn jūmín bànlǐ gè zhǒng shǒuxù, běn zhōngxīn tígōng yǐxià fúwù.\\
Afin de faciliter aux habitants l'accomplissement de toutes sortes de démarches, notre centre offre les services suivants.\\[0.5em]
一、办理身份证和出生证明的申请。\\
Yī, bànlǐ shēnfènzhèng hé chūshēng zhèngmíng de shēnqǐng.\\
Premièrement, réception des demandes de cartes d'identité et d'actes de naissance.\\[0.5em]
二、提供户籍证明和各种表格。\\
Èr, tígōng hùjí zhèngmíng hé gè zhǒng biǎogé.\\
Deuxièmement, délivrance d'attestations de résidence et de divers formulaires.\\[0.5em]
三、回答居民的问题，提供信息咨询。\\
Sān, huídá jūmín de wèntí, tígōng xìnxī zīxún.\\
Troisièmement, réponses aux questions des habitants et information-conseil.\\[0.5em]
办公时间：上午八点到下午三点半，星期一到星期五。\\
Bàngōng shíjiān: shàngwǔ bā diǎn dào xiàwǔ sān diǎn bàn, xīngqīyī dào xīngqīwǔ.\\
Horaires d'ouverture : de 8 h à 15 h 30, du lundi au vendredi.\\[0.5em]
地址：雅温得市恩科隆东戈区。电话：二二二二二二二二。\\
Dìzhǐ: Yǎwēndé shì Ēnkēlóngdōnggē qū. Diànhuà: èr èr èr èr èr èr èr èr.\\
Adresse : quartier Nkolndongo, ville de Yaoundé. Téléphone : 22 22 22 22.\\[0.5em]
如有问题，请随时联系我们。谢谢合作！\\
Rú yǒu wèntí, qǐng suíshí liánxì wǒmen. Xièxie hézuò!\\
Si vous avez des questions, contactez-nous à tout moment. Merci de votre collaboration !
\end{textebox}

\textbf{Questions d'observation :}
\begin{enumerate}
\item Quel mot ouvre l'avis et annonce sa nature ? (\zh{公告})
\item Quelle expression introduit le but de l'avis ? (\zh{为了...})
\item Comment les services sont-ils numérotés ? (\zh{一、二、三})
\item Quelle formule termine l'avis ? (\zh{如有问题，请...；谢谢合作！})
\end{enumerate}

\courssub{La structure type en cinq blocs}

Tout avis administratif chinois d'offre de services suit un plan quasi immuable, en cinq blocs. Le mémoriser, c'est déjà savoir rédiger.

\begin{definition}[Les cinq blocs de l'avis de services]
\begin{enumerate}
\item \cle{Titre} : \zh{公告} (gōnggào, « avis public ») ou \zh{通知} (tōngzhī, « notification, note d'information »), souvent précédé d'un nom : \zh{服务公告} « avis de services ».
\item \cle{Objet} : introduit par \zh{为了} + but, suivi de \zh{本...提供以下服务} « notre... offre les services suivants ».
\item \cle{Énumération des services} : \zh{一、二、三} (style officiel) ou \zh{首先...其次...另外...最后}.
\item \cle{Informations pratiques} : \zh{办公时间} (horaires), \zh{地址} (adresse), \zh{电话} (téléphone).
\item \cle{Clôture} : \zh{如有问题，请...} puis \zh{谢谢合作！} « merci de votre collaboration ».
\end{enumerate}
\end{definition}

\begin{popfigure}
\begin{tikzpicture}[
bloc/.style={rectangle, rounded corners=3pt, draw=popblue, fill=popblueL, text width=9.5cm, align=center, inner sep=6pt, font=\small},
fleche/.style={-{Stealth[length=2.5mm]}, thick, poporange}
]
\node[bloc] (b1) at (0,0) {\textbf{1. Titre} : 公告 / 通知 (avis)};
\node[bloc, below=0.35cm of b1] (b2) {\textbf{2. Objet} : 为了 + but，本...提供以下服务。};
\node[bloc, below=0.35cm of b2] (b3) {\textbf{3. Services} : 一、... 二、... 三、...};
\node[bloc, below=0.35cm of b3] (b4) {\textbf{4. Infos pratiques} : 办公时间、地址、电话};
\node[bloc, below=0.35cm of b4] (b5) {\textbf{5. Clôture} : 如有问题，请... 谢谢合作！};
\draw[fleche] (b1) -- (b2);
\draw[fleche] (b2) -- (b3);
\draw[fleche] (b3) -- (b4);
\draw[fleche] (b4) -- (b5);
\end{tikzpicture}
\legende{Schéma de la structure type d'un avis de services en cinq blocs}
\end{popfigure}

\courssub{Le connecteur d'objet : 为了 + but}

Observons la deuxième phrase de notre avis : \zh{为了方便居民办理各种手续，本中心提供以下服务。} Le but (« faciliter les démarches des habitants ») est placé \emph{avant} l'annonce, introduit par \zh{为了}.

\begin{propriete}[La construction 为了 + but]
\zh{为了} (wèile, « afin de, pour ») introduit une proposition de but qui se place \cle{en tête de phrase}, suivie d'une virgule, puis de la proposition principale.\\
Structure : \textbf{为了 + verbe/nom de but ，sujet + verbe principal.}\\
Exemple : \zh{为了方便居民，我们提供以下服务。}\\
\emph{Wèile fāngbiàn jūmín, wǒmen tígōng yǐxià fúwù.}\\
« Pour faciliter la vie des habitants, nous offrons les services suivants. »\\
Autre exemple : \zh{为了提高服务质量，本中心每天开门。}\\
\emph{Wèile tígāo fúwù zhìliàng, běn zhōngxīn měi tiān kāi mén.}\\
« Afin d'améliorer la qualité du service, notre centre ouvre tous les jours. »
\end{propriete}

\begin{attention}[为了 ou 因为 ?]
Ne confondez pas \zh{为了} (wèile, « afin de » : le \emph{but}, orienté vers le futur) et \zh{因为} (yīnwèi, « parce que » : la \emph{cause}, orientée vers le passé). Les francophones traduisent souvent « pour » par \zh{因为} : erreur grave. Comparez : \zh{为了方便居民，我们延长了办公时间。} « Pour faciliter la vie des habitants, nous avons prolongé les horaires » (but) ; \zh{因为下大雨，中心今天不开门。} « À cause de la forte pluie, le centre n'ouvre pas aujourd'hui » (cause).
\end{attention}

\courssub{L'énumération administrative : 一、二、三 et les connecteurs}

Le chinois administratif dispose de deux systèmes pour énumérer.

\textbf{1. Le style officiel : les chiffres chinois suivis du signe 、}

\zh{一、} \zh{二、} \zh{三、} (yī, èr, sān) : c'est l'équivalent de nos « 1. 2. 3. » ou « Premièrement, deuxièmement... ». Le signe \zh{、} (dùnhào) est la virgule d'énumération chinoise, réservée aux listes.

\textbf{2. Les connecteurs de progression}

\begin{center}
\begin{tabularx}{\linewidth}{|X|X|X|X|}
\hline
\rowcolor{popblueL} Caractères & Pinyin & Nature & Traduction \\
\hline
首先 & shǒuxiān & connecteur & d'abord, en premier lieu \\
\hline
其次 & qícì & connecteur & ensuite, en second lieu \\
\hline
另外 & lìngwài & connecteur & de plus, par ailleurs \\
\hline
最后 & zuìhòu & connecteur & enfin, pour finir \\
\hline
以下 & yǐxià & marqueur & ci-dessous, ce qui suit \\
\hline
\end{tabularx}
\end{center}

\begin{exemplebox}[Énumération avec connecteurs]
\zh{首先，请填写申请表；其次，提交材料；最后，在窗口领取证明。}\\
\emph{Shǒuxiān, qǐng tiánxiě shēnqǐngbiǎo; qícì, tíjiāo cáiliào; zuìhòu, zài chuāngkǒu lǐngqǔ zhèngmíng.}\\
« D'abord, remplissez le formulaire de demande ; ensuite, déposez les pièces ; enfin, retirez l'attestation au guichet. »\\[0.5em]
Autre exemple : \zh{首先，本中心办理身份证；其次，我们提供咨询；另外，居民可以复印文件；最后，我们回答所有问题。}\\
\emph{Shǒuxiān, běn zhōngxīn bànlǐ shēnfènzhèng; qícì, wǒmen tígōng zīxún; lìngwài, jūmín kěyǐ fùyìn wénjiàn; zuìhòu, wǒmen huídá suǒyǒu wèntí.}\\
« D'abord, notre centre traite les cartes d'identité ; ensuite, nous offrons des conseils ; de plus, les habitants peuvent photocopier des documents ; enfin, nous répondons à toutes les questions. »
\end{exemplebox}

\begin{attention}[« Ci-dessous » ne se dit pas 下面 dans un avis]
Dans la langue courante, « en bas, dessous » se dit \zh{下面} (xiàmian). Mais dans un texte officiel, « ci-dessous, ce qui suit » exige la formule consacrée \zh{以下} (yǐxià) : \zh{提供以下服务} « offre les services ci-dessous ». Écrire \zh{提供下面的服务} dans un avis ferait très oral et peu professionnel. Retenez : \cle{avis officiel = 以下}.
\end{attention}

\begin{popfigure}
\begin{tikzpicture}[
etape/.style={circle, draw=poppurple, fill=poppurpleL, thick, minimum size=1.5cm, align=center, font=\small},
lbl/.style={below=0.15cm, font=\small\itshape, popdark},
fleche/.style={-{Stealth[length=3mm]}, very thick, poporange}
]
\node[etape] (e1) at (0,0) {首先};
\node[etape] (e2) at (3.2,0) {其次};
\node[etape] (e3) at (6.4,0) {另外};
\node[etape] (e4) at (9.6,0) {最后};
\node[lbl] at (e1.south) {d'abord};
\node[lbl] at (e2.south) {ensuite};
\node[lbl] at (e3.south) {de plus};
\node[lbl] at (e4.south) {enfin};
\draw[fleche] (e1) -- (e2);
\draw[fleche] (e2) -- (e3);
\draw[fleche] (e3) -- (e4);
\end{tikzpicture}
\legende{Frise des connecteurs d'énumération : la progression logique d'un avis}
\end{popfigure}

\courssub{La formule conditionnelle de service et la clôture}

La clôture d'un avis chinois est presque toujours une invitation polie à contacter l'administration, construite sur une condition.

\begin{propriete}[La structure 如（果）...，请...]
\zh{如果} (rúguǒ, « si ») introduit la condition ; dans le style écrit administratif, on utilise souvent la forme abrégée \zh{如}. La conséquence est un impératif poli introduit par \zh{请} (qǐng, « s'il vous plaît, veuillez »).\\
Structure : \textbf{如（果）+ condition ，请 + action.}\\
Modèle : \zh{如有问题，请随时联系我们。}\\
\emph{Rú yǒu wèntí, qǐng suíshí liánxì wǒmen.}\\
« Si vous avez des questions, contactez-nous à tout moment. »\\
Variantes : \zh{如需帮助，请到一号窗口。} « Si vous avez besoin d'aide, adressez-vous au guichet n\textsuperscript{o} 1. » ; \zh{如要申请，请带身份证。} « Si vous voulez faire une demande, apportez votre carte d'identité. »
\end{propriete}

La formule finale \zh{谢谢合作！} (Xièxie hézuò !, litt. « merci [de votre] coopération ») est la clôture standard des avis chinois ; elle correspond à notre « Merci de votre compréhension / collaboration ».

\courssub{Comparaison avec le français et registre de langue}

\begin{center}
\begin{tabularx}{\linewidth}{|X|X|X|}
\hline
\rowcolor{popblueL} Point de langue & En français & En chinois \\
\hline
Place du but & souvent après : « Nous offrons ces services afin de... » & avant : 为了...，本中心提供... \\
\hline
Articles et pluriel & « les services », « des formulaires » & aucun article, aucun pluriel : 服务, 表格 \\
\hline
Horaires & « de 8 h à 17 h » (prépositions) & sans préposition : 上午八点到下午五点 \\
\hline
« Notre » administration & « la mairie », « notre centre » & 本中心, 本校, 本市 (marqueur impersonnel) \\
\hline
Politesse & « veuillez », « nous vous remercions » & 请 + verbe ; 谢谢合作 \\
\hline
\end{tabularx}
\end{center}

\begin{attention}[Pas de 我/你 dans un avis officiel]
Un avis administratif chinois évite \zh{我} (« je ») et \zh{你} (« tu »). L'administration se désigne par \zh{本} + nom : \zh{本中心} « notre centre », \zh{本校} « notre école », \zh{本市} « notre ville ». L'usager est désigné par \zh{居民} (« les habitants ») ou interpellé indirectement par \zh{请}. Écrire \zh{我给你提供服务} dans un avis serait une faute de registre : trop direct, trop personnel.
\end{attention}

\begin{savaistu}[本 : le « notre » des administrations chinoises]
Le caractère \zh{本} (běn) signifie à l'origine « racine » (c'est l'arbre \zh{木} avec un trait à sa base). De « racine, origine », il est passé au sens de « celui-ci même, le présent » : \zh{本中心} « le présent centre », \zh{本市} « la présente ville », \zh{本人} « ma personne, moi-même » (sur les formulaires). C'est l'équivalent exact de nos formules « la présente administration », « le soussigné ». Au Cameroun, nos avis officiels disent « La Mairie informe les populations que... » ; le chinois dirait \zh{本市通知居民...} : même distance impersonnelle, même politesse institutionnelle.
\end{savaistu}

\courssub{Méthode de rédaction et entraînement}

\begin{methode}[Rédiger un avis de services en cinq étapes]
\begin{enumerate}
\item \textbf{Planifier en français} : remplir le plan en cinq blocs (titre, objet, trois services, infos pratiques, clôture) avec des mots simples.
\item \textbf{Traduire le titre} : choisir \zh{公告} ou \zh{通知}, éventuellement précédé d'un nom (\zh{服务公告}).
\item \textbf{Traduire l'objet} : construire \zh{为了} + but, puis \zh{本...提供以下服务。}
\item \textbf{Énumérer} : rédiger chaque service avec \zh{一、二、三} ou \zh{首先...其次...另外...最后}, une phrase courte par service.
\item \textbf{Clôturer} : ajouter horaires (\zh{办公时间：...}), adresse, téléphone, puis \zh{如有问题，请随时联系我们。谢谢合作！} Relire en vérifiant l'absence de \zh{我}/\zh{你} et de \zh{下面}.
\end{enumerate}
\end{methode}

\begin{exoresolu}[Transformer des notes en avis structuré]
Énoncé : voici des notes en désordre rédigées par le secrétaire du centre communautaire de Buea : « ouvert du lundi au vendredi, 8 h--15 h ; demandes de passeports ; conseils aux habitants ; photocopies ; pour aider les habitants ; si questions, téléphoner ». Rédigez le début de l'avis (blocs 1 à 3) en chinois.
\tcblower
Solution détaillée :\\
Bloc 1 (titre) : \zh{服务公告} (Fúwù gōnggào).\\
Bloc 2 (objet avec 为了) : \zh{为了帮助居民，本中心提供以下服务。} \emph{Wèile bāngzhù jūmín, běn zhōngxīn tígōng yǐxià fúwù.} « Pour aider les habitants, notre centre offre les services suivants. »\\
Bloc 3 (énumération) : \zh{一、办理护照申请。二、回答居民的问题。三、复印文件。} \emph{Yī, bànlǐ hùzhào shēnqǐng. Èr, huídá jūmín de wèntí. Sān, fùyìn wénjiàn.} « 1. Réception des demandes de passeport. 2. Réponses aux questions des habitants. 3. Photocopies de documents. »\\
Les horaires et la clôture (blocs 4 et 5) seront rédigés en production guidée : \zh{办公时间：上午八点到下午三点。如有问题，请打电话。谢谢合作！}
\end{exoresolu}

\begin{exercice}[À vous de structurer]
Remettez les phrases suivantes dans l'ordre des cinq blocs, puis traduisez-les en français : \zh{谢谢合作！} / \zh{一、办理出生证明。二、提供表格。} / \zh{服务通知} / \zh{办公时间：上午九点到下午四点。} / \zh{为了方便居民，本中心提供以下服务。}
\end{exercice}

\begin{corrige}[Exercice : À vous de structurer]
Ordre correct : 1) \zh{服务通知} (titre) ; 2) \zh{为了方便居民，本中心提供以下服务。} (objet) ; 3) \zh{一、办理出生证明。二、提供表格。} (services) ; 4) \zh{办公时间：上午九点到下午四点。} (infos pratiques) ; 5) \zh{谢谢合作！} (clôture).\\
Traduction : « Avis de services. — Pour faciliter la vie des habitants, notre centre offre les services suivants. — 1. Traitement des actes de naissance. 2. Fourniture de formulaires. — Horaires : de 9 h à 16 h. — Merci de votre collaboration ! »
\end{corrige}

\begin{aretenir}[L'essentiel de la section]
\begin{itemize}
\item Un avis de services se construit en \cle{cinq blocs} : titre (\zh{公告/通知}) → objet (\zh{为了...，本...提供以下服务}) → énumération (\zh{一、二、三} ou \zh{首先...其次...另外...最后}) → infos pratiques (\zh{办公时间、地址、电话}) → clôture (\zh{如有问题，请...；谢谢合作！}).
\item \zh{为了} + but se place en tête de phrase ; ne pas confondre avec \zh{因为} (cause).
\item « Ci-dessous » officiel = \zh{以下}, jamais \zh{下面} dans un avis.
\item Registre impersonnel : \zh{本中心}, \zh{请}, pas de \zh{我}/\zh{你} ; horaires sans préposition : \zh{上午八点到下午五点}.
\end{itemize}
\end{aretenir}

Nous disposons maintenant de tous les outils : le lexique, les caractères, la structure et les connecteurs. Dans la prochaine section, nous mettrons ces acquis à l'épreuve avec des exercices de traduction guidée, du français vers le chinois (thème) et du chinois vers le français (version).

\coursec{Thème et version guidés}

Dans la section précédente, nous avons appris à structurer une offre de services : annoncer le service, indiquer les horaires, préciser les pièces à fournir et orienter l'usager. Il est temps maintenant de mettre ces acquis à l'épreuve de la traduction, dans les deux sens : la \cle{version} (du chinois vers le français) et le \cle{thème} (du français vers le chinois). C'est l'exercice roi au baccalauréat : il vérifie à la fois votre compréhension, votre grammaire et votre écriture des caractères. Procédons pas à pas, avec méthode.

\courssub{La version : du chinois vers le français}

\textbf{Observer d'abord, traduire ensuite}\par

Lisez attentivement les deux phrases suivantes, extraites d'un avis affiché à l'entrée d'un bureau d'état civil :

\begin{exemplebox}[Deux phrases à traduire]
Phrase 1 : 本办公室为居民办理出生证明。办公时间：上午九点到下午四点。\\
\emph{Běn bàngōngshì wèi jūmín bànlǐ chūshēng zhèngmíng. Bàngōng shíjiān : shàngwǔ jiǔ diǎn dào xiàwǔ sì diǎn.}\\[0.3em]
Phrase 2 : 如有问题，请到三号窗口咨询。\\
\emph{Rú yǒu wèntí, qǐng dào sān hào chuāngkǒu zīxún.}
\end{exemplebox}

Avant de traduire, repérons les mots clés :

\begin{center}
\begin{tabularx}{\linewidth}{|X|X|X|X|}
\hline
\rowcolor{popblueL} Caractères & Pinyin & Nature & Traduction \\
\hline
本办公室 & běn bàngōngshì & nom & notre bureau, ce bureau \\
\hline
为 & wèi & préposition & pour, au bénéfice de \\
\hline
居民 & jūmín & nom & habitants, résidents \\
\hline
办理 & bànlǐ & verbe & traiter, effectuer (démarche administrative) \\
\hline
出生证明 & chūshēng zhèngmíng & nom & acte de naissance \\
\hline
办公时间 & bàngōng shíjiān & nom & horaires de bureau \\
\hline
到 & dào & verbe/prép. & aller à, jusqu'à \\
\hline
窗口 & chuāngkǒu & nom & guichet \\
\hline
咨询 & zīxún & verbe & demander des informations, consulter \\
\hline
\end{tabularx}
\end{center}

\begin{methode}[Méthode de version en 4 étapes]
\begin{enumerate}
\item \textbf{Traduire d'abord le sens global.} Lisez toute la phrase et demandez-vous : de quoi parle-t-on ? Qui fait quoi, pour qui ? Ici : un bureau, des habitants, des actes de naissance, des horaires.
\item \textbf{Identifier le sujet et le verbe principal.} Phrase 1 : sujet = 本办公室 (notre bureau), verbe = 办理 (délivrer/traiter).
\item \textbf{Vérifier chaque caractère clé.} Ne confondez pas les verbes administratifs : \zh{办理} (bànlǐ, « effectuer une démarche ») n'est ni \zh{服务} (fúwù, « servir ») ni \zh{提供} (tígōng, « fournir, offrir »). Chacun a son emploi propre.
\item \textbf{Rédiger en français naturel.} Le français exige des articles et des liaisons que le chinois n'a pas : 为居民 devient « aux habitants », pas « pour habitants ».
\end{enumerate}
\end{methode}

\begin{exoresolu}[Version 1 : l'avis du bureau d'état civil]
Énoncé. Traduisez en français : 本办公室为居民办理出生证明。办公时间：上午九点到下午四点。
\tcblower
Solution détaillée.
\begin{itemize}
\item 本办公室 = « notre bureau » (本 marque ici « notre / le présent », très fréquent dans les avis officiels).
\item 为居民 = « pour les habitants » → en français naturel : « aux habitants ».
\item 办理出生证明 = « traiter les actes de naissance » → on dira plus élégamment « délivrer les actes de naissance ».
\item 办公时间 = « horaires de bureau » ; 上午九点到下午四点 = « de 9 h du matin à 4 h de l'après-midi ».
\end{itemize}
Traduction finale : « Notre bureau délivre les actes de naissance aux habitants. Horaires : de 9 h à 16 h. »
\end{exoresolu}

\begin{exoresolu}[Version 2 : orientation de l'usager]
Énoncé. Traduisez en français : 如有问题，请到三号窗口咨询。
\tcblower
Solution détaillée.
\begin{itemize}
\item 如有问题 = « si (vous) avez des questions » ; 如 est la forme écrite et soutenue de 如果, typique des avis administratifs.
\item 请 = marque de politesse, « veuillez / s'il vous plaît ».
\item 到三号窗口 = « au guichet numéro 3 » ; la structure est 到 + lieu + verbe : « aller à un endroit pour faire quelque chose ».
\item 咨询 = « demander des informations ».
\end{itemize}
Traduction finale : « Pour toute question, adressez-vous au guichet n° 3. » (ou : « Si vous avez des questions, veuillez vous renseigner au guichet n° 3. »)
\end{exoresolu}

\begin{exemplebox}[La structure 到 + lieu + verbe]
\zh{请到二号窗口办理。}\\
\emph{Qǐng dào èr hào chuāngkǒu bànlǐ.}\\
« Veuillez vous adresser au guichet n° 2 (pour effectuer la démarche). »\\[0.4em]
Autres exemples :\\
\zh{请到一楼交表。} \emph{Qǐng dào yī lóu jiāo biǎo.} « Veuillez déposer le formulaire au rez-de-chaussée. »\\
\zh{居民到市政府办身份证。} \emph{Jūmín dào shì zhèngfǔ bàn shēnfènzhèng.} « Les habitants vont à la mairie pour faire leur carte d'identité. »
\end{exemplebox}

\courssub{Le thème : du français vers le chinois}

En thème, le danger est de traduire mot à mot en gardant l'ordre français. Or le chinois impose sa propre architecture. Observons trois phrases modèles.

\begin{exemplebox}[Trois phrases de thème corrigées]
1. « Notre mairie offre les services suivants aux habitants. »\\
\zh{我们市政府为居民提供以下服务。}\\
\emph{Wǒmen shì zhèngfǔ wèi jūmín tígōng yǐxià fúwù.}\\[0.4em]
2. « Vous devez apporter votre carte d'identité et le formulaire. »\\
\zh{您需要带身份证和申请表。}\\
\emph{Nín xūyào dài shēnfènzhèng hé shēnqǐng biǎo.}\\[0.4em]
3. « Le centre est ouvert du lundi au vendredi, de 8 h à 17 h. »\\
\zh{中心星期一至星期五开门，上午八点到下午五点。}\\
\emph{Zhōngxīn xīngqīyī zhì xīngqīwǔ kāimén, shàngwǔ bā diǎn dào xiàwǔ wǔ diǎn.}\\
Formule d'affiche : \zh{办公时间：星期一至星期五，上午八点到下午五点。}
\end{exemplebox}

\begin{propriete}[需要 + verbe : « devoir, il faut »]
\zh{需要} (xūyào) suivi d'un verbe exprime la nécessité : « il faut, il est nécessaire de ». C'est la tournure administrative standard, plus neutre que \zh{必须} (bìxū, « il faut absolument »).
\\[0.3em]
Structure : Sujet + 需要 + verbe (+ complément).\\
\zh{您需要提交材料。} \emph{Nín xūyào tíjiāo cáiliào.} « Vous devez déposer les pièces. »\\
\zh{办护照需要带照片。} \emph{Bàn hùzhào xūyào dài zhàopiàn.} « Pour le passeport, il faut apporter une photo. »
\end{propriete}

\begin{attention}[Le complément de temps se place AVANT le verbe]
En français on dit « Le centre ouvre \emph{à 8 h} » ; en chinois, l'indication de temps précède le verbe :\\
\zh{我们上午八点开门。} \emph{Wǒmen shàngwǔ bā diǎn kāimén.} « Nous ouvrons à 8 h du matin. »\\
Jamais : *我们开门上午八点。 (ordre français calqué, faute classique des francophones).\\
Règle générale : Sujet + (temps) + (lieu) + verbe + complément.\\
\zh{居民星期一上午到市政府办手续。} « Les habitants vont lundi matin à la mairie faire leurs démarches. »
\end{attention}

\begin{attention}[Ne confondez pas 为 et 办]
\zh{为} (wèi) est une préposition : « pour, au bénéfice de » → \zh{为居民服务} « servir les habitants ».\\
\zh{办} (bàn) est un verbe : « faire, traiter » → \zh{办身份证} « faire une carte d'identité ».\\
Piège : dans « la mairie délivre les actes \emph{pour} les habitants », le « pour » se rend par 为 : \zh{市政府为居民办理证明。} Ne dites pas *市政府办居民证明 (sens : « la mairie fait les habitants-certificats », incompréhensible).
\end{attention}

\begin{center}
\begin{tabularx}{\linewidth}{|X|X|X|}
\hline
\rowcolor{popblueL} Verbe administratif & Collocation typique & Traduction \\
\hline
办理 bànlǐ & 办理手续 / 办理证明 & effectuer une démarche / délivrer un acte \\
\hline
提供 tígōng & 提供服务 & fournir / offrir des services \\
\hline
提交 tíjiāo & 提交材料 / 提交申请表 & déposer les pièces / le formulaire \\
\hline
服务 fúwù & 为居民服务 & servir les habitants \\
\hline
咨询 zīxún & 到窗口咨询 & se renseigner au guichet \\
\hline
\end{tabularx}
\end{center}

\courssub{Stratégie de thème : la méthode en 4 étapes}

\begin{popfigure}
\begin{tikzpicture}[node distance=0.55cm, every node/.style={font=\small}]
\node[draw=popblue, fill=popblueL, rounded corners, align=center, text width=3.1cm, minimum height=1.1cm] (e1) {\textbf{1. Repérer}\\sujet et verbe};
\node[draw=poporange, fill=poporangeL, rounded corners, align=center, text width=3.1cm, minimum height=1.1cm, right=of e1] (e2) {\textbf{2. Placer}\\temps et lieu\\avant le verbe};
\node[draw=popgreen, fill=popgreenL, rounded corners, align=center, text width=3.1cm, minimum height=1.1cm, right=of e2] (e3) {\textbf{3. Choisir}\\la bonne collocation\\(办理 / 提供 / 提交)};
\node[draw=poppurple, fill=poppurpleL, rounded corners, align=center, text width=3.1cm, minimum height=1.1cm, right=of e3] (e4) {\textbf{4. Relire}\\tons et caractères};
\draw[-{Stealth}, thick, popdark] (e1) -- (e2);
\draw[-{Stealth}, thick, popdark] (e2) -- (e3);
\draw[-{Stealth}, thick, popdark] (e3) -- (e4);
\end{tikzpicture}
\legende{Méthode de thème en 4 étapes : de l'analyse française à la phrase chinoise correcte.}
\end{popfigure}

\begin{methode}[Appliquer la méthode : « Vous devez apporter votre carte d'identité et le formulaire. »]
\begin{enumerate}
\item \textbf{Sujet et verbe} : sujet = « vous » → 您 ; verbe d'action = « apporter » → 带 ; la nécessité « devoir » → 需要.
\item \textbf{Temps et lieu} : aucun ici ; on passe à l'étape suivante.
\item \textbf{Collocation} : « apporter une carte » → 带身份证 ; « le formulaire » administratif → 申请表 ; coordination → 和.
\item \textbf{Relire} : 您需要带身份证和申请表。 Vérifier : 需 (clé 雨 « pluie » en haut), 表 (pas de trait en trop), tons : nín xūyào dài shēnfènzhèng hé shēnqǐng biǎo.
\end{enumerate}
\end{methode}

\begin{exoresolu}[Thème guidé : les horaires du centre]
Énoncé. Traduisez en chinois : « Le centre est ouvert du lundi au vendredi, de 8 h à 17 h. »
\tcblower
Solution détaillée.
\begin{enumerate}
\item Sujet : « le centre » → 中心. Verbe : « être ouvert » → 开门 (litt. « ouvrir la porte »).
\item « du lundi au vendredi » : 星期一至星期五 (至 zhì = « jusqu'à », forme écrite ; à l'oral on dirait 到). Ce complément de temps se place avant le verbe : 中心星期一至星期五开门.
\item « de 8 h à 17 h » : 上午八点到下午五点, placé après en précision d'horaire, ou intégré dans la formule d'affiche 办公时间：…
\item Traduction finale : 中心星期一至星期五开门，上午八点到下午五点。 Variante d'affiche : 办公时间：星期一至星期五，上午八点到下午五点。
\end{enumerate}
Piège évité : ne pas écrire *中心开门星期一至星期五 (temps après le verbe = faute).
\end{exoresolu}

\courssub{Entraînement collectif : mini-dictée de thème}

\begin{textebox}[Mini-dictée de thème — 10 minutes, correction au tableau]
Traduisez les trois phrases suivantes en chinois, puis justifiez chaque choix (ordre des mots, verbe administratif, caractères).\\[0.4em]
1. « Notre bureau offre les services suivants aux habitants. »\\
2. « Vous devez apporter le formulaire et une photo. »\\
3. « Pour toute question, adressez-vous au guichet n° 2. »\\[0.4em]
\emph{Corrigé attendu :}\\
1. 我们办公室为居民提供以下服务。 \emph{Wǒmen bàngōngshì wèi jūmín tígōng yǐxià fúwù.}\\
2. 您需要带申请表和一张照片。 \emph{Nín xūyào dài shēnqǐng biǎo hé yī zhāng zhàopiàn.} (classificateur 张 pour les photos et les feuilles !)\\
3. 如有问题，请到二号窗口咨询。 \emph{Rú yǒu wèntí, qǐng dào èr hào chuāngkǒu zīxún.}
\end{textebox}

\begin{savaistu}[Administration et services au public : Cameroun - Chine]
Au Cameroun, la décentralisation a transféré de nombreuses compétences aux \textbf{communes} (mairies) : état civil, marchés, voirie, hygiène. L'usager se déplace souvent physiquement au \textbf{guichet unique} ou au service concerné. En Chine, l'administration locale (\textbf{街道办事处} jiēdào bànshìchù, « bureau des affaires de rue » au niveau quartier ; \textbf{镇政府} zhèn zhèngfǔ, « gouvernement de bourg » en zone rurale) gère l'état civil, la sécurité sociale, le registre de résidence (\textbf{户口} hùkǒu). La numérisation est massive : prise de rendez-vous via WeChat/Alipay (\textbf{预约} yùyuē), paiement en ligne, retrait de documents aux bornes automatiques (\textbf{自助终端} zìzhù zhōngduān). Le vocabulaire appris ici (窗口, 办理, 材料) reste la base, même dans le tout-numérique.
\end{savaistu}

\begin{experience}[Jeu de rôle : « Au guichet de la mairie »]
\textbf{Consigne :} Par binômes. L'un joue l'agent (坐在桌后), l'autre l'usager (arrive avec un sac).
\textbf{Scénario A (État civil) :} L'usager demande un acte de naissance. L'agent demande la carte d'identité, indique le formulaire à remplir, donne le délai (3 jours) et le guichet de retrait.
\textbf{Scénario B (Urbanisme) :} L'usager dépose un dossier de permis de construire. L'agent vérifie la liste des pièces (plan, titre foncier, photos), tamponne le récépissé, indique le numéro de suivi.
\textbf{Contraintes linguistiques :} Utiliser obligatoirement : 为您办理, 需要带, 请到...窗口, 提交材料, 办公时间. Noter les phrases clés sur une fiche avant de jouer.
\end{experience}

\begin{exercice}[À vous de jouer]
Traduisez en chinois :\\
1. « La mairie délivre les actes de mariage aux habitants. »\\
2. « Le bureau est ouvert du lundi au vendredi, de 9 h à 16 h. »\\
3. « Vous devez déposer les pièces au guichet n° 3. »\\
Traduisez en français :\\
4. 居民需要带身份证和出生证明。\\
5. 办公时间：上午八点到下午五点。
\end{exercice}

\begin{corrige}[Exercice « À vous de jouer »]
1. 市政府为居民办理结婚证明。 \emph{Shì zhèngfǔ wèi jūmín bànlǐ jiéhūn zhèngmíng.} (为 + bénéficiaire avant le verbe 办理.)\\
2. 办公室星期一至星期五开门，上午九点到下午四点。 \emph{Bàngōngshì xīngqīyī zhì xīngqīwǔ kāimén, shàngwǔ jiǔ diǎn dào xiàwǔ sì diǎn.} (temps avant le verbe 开门.)\\
3. 您需要到三号窗口提交材料。 \emph{Nín xūyào dào sān hào chuāngkǒu tíjiāo cáiliào.} (需要 + verbe ; 到 + lieu + verbe.)\\
4. « Les habitants doivent apporter leur carte d'identité et leur acte de naissance. »\\
5. « Horaires de bureau : de 8 h à 17 h. »
\end{corrige}

\begin{aretenir}[L'essentiel du thème et de la version]
\begin{itemize}
\item En version : sens global d'abord, puis vérification de chaque verbe clé (办理 ≠ 提供 ≠ 服务).
\item En thème : sujet + temps + lieu + verbe + complément ; le temps ne suit jamais le verbe.
\item « Devoir » administratif → 需要 + verbe ; politesse → 请 en tête de phrase ; orientation → 到 + lieu + verbe.
\item « Pour les habitants » → 为居民, jamais un calque avec 办.
\item Toujours relire : tons, caractères (需, 窗, 表), classificateurs (一张照片, 一份材料).
\end{itemize}
\end{aretenir}

\coursec{Production guidée et évaluation}

Dans la section précédente, vous avez appris à traduire des phrases simples sur les services municipaux, dans les deux sens (thème et version). Vous disposez maintenant du lexique (état civil, taxes, plaintes, informations), des structures (\zh{为了} pour l'objet, \zh{请} pour les instructions) et des connecteurs (\zh{首先}、\zh{其次}、\zh{最后}). Il est temps de mobiliser tout cela dans une \cle{production écrite complète} : la rédaction d'un avis d'offre de services aux usagers, exercice central du Module 3 et situation très proche de ce que l'on peut vous demander au baccalauréat.

\courssub{La tâche finale : rédiger un avis d'offre de services}

\begin{definition}[La tâche]
Rédiger, en chinois simplifié, un \cle{avis d'offre de services} pour la communauté urbaine de Douala (ou le council de Bamenda), destiné aux usagers. Longueur exigée : \textbf{80 à 120 caractères chinois} (la ponctuation ne compte pas). L'avis doit informer les habitants des services offerts, des horaires et des modalités de contact.
\end{definition}

Avant d'écrire, observons un modèle complet, légèrement plus long que la cible, pour bien identifier les cinq blocs obligatoires.

\begin{textebox}[Modèle commenté : avis de la communauté urbaine de Douala]
\zh{杜阿拉城市社区服务通知}\\
\emph{Dū'ālā chéngshì shèqū fúwù tōngzhī}\\
« Avis de services de la communauté urbaine de Douala »\\[4pt]
\zh{为了方便市民，本社区提供以下服务：}\\
\emph{Wèile fāngbiàn shìmín, běn shèqū tígōng yǐxià fúwù:}\\
« Pour faciliter la vie des habitants, notre communauté offre les services suivants : »\\[4pt]
\zh{一、办理出生证明和身份证；}\\
\emph{Yī, bànlǐ chūshēng zhèngmíng hé shēnfènzhèng;}\\
« 1. Établissement des actes de naissance et des cartes d'identité ; »\\[4pt]
\zh{二、收取地方税；}\\
\emph{Èr, shōuqǔ dìfāng shuì;}\\
« 2. Perception des taxes locales ; »\\[4pt]
\zh{三、接受市民的投诉和建议。}\\
\emph{Sān, jiēshòu shìmín de tóusù hé jiànyì.}\\
« 3. Réception des plaintes et des suggestions des habitants. »\\[4pt]
\zh{办公时间：周一至周五，上午八点到下午四点。}\\
\emph{Bàngōng shíjiān: zhōuyī zhì zhōuwǔ, shàngwǔ bā diǎn dào xiàwǔ sì diǎn.}\\
« Horaires : du lundi au vendredi, de 8 h à 16 h. »\\[4pt]
\zh{地址：杜阿拉市政府大楼二楼。电话：233-42-XX-XX。}\\
\emph{Dìzhǐ: Dū'ālā shì zhèngfǔ dàlóu èr lóu. Diànhuà: 233-42-XX-XX.}\\
« Adresse : 2\textsuperscript{e} étage de l'hôtel de ville de Douala. Téléphone : 233-42-XX-XX. »\\[4pt]
\zh{谢谢您的合作！}\\
\emph{Xièxie nín de hézuò!}\\
« Merci de votre coopération ! »
\end{textebox}

Observez : le titre annonce le thème ; \zh{为了} introduit l'objet ; les services sont numérotés avec \zh{一、二、三、} ; les deux-points \zh{：} introduisent la liste et les horaires ; la clôture remercie l'usager. C'est exactement le squelette que vous allez reproduire.

\courssub{Étape 1 : brainstormer les services et les traduire}

En classe entière ou par deux, listez d'abord en français les services qu'une collectivité locale peut offrir, puis traduisez-les à l'aide du glossaire. Ne retenez que trois services pour respecter la limite de caractères.

\begin{center}
\begin{tabularx}{\linewidth}{|X|X|X|X|}
\hline
\rowcolor{popblueL} Caractères & Pinyin & Nature & Traduction \\
\hline
办理出生证明 & bànlǐ chūshēng zhèngmíng & locution verbale & établir un acte de naissance \\
\hline
办理身份证 & bànlǐ shēnfènzhèng & locution verbale & établir une carte d'identité \\
\hline
收取地方税 & shōuqǔ dìfāng shuì & locution verbale & percevoir les taxes locales \\
\hline
接受投诉 & jiēshòu tóusù & locution verbale & recevoir les plaintes \\
\hline
提供信息 & tígōng xìnxī & locution verbale & fournir des informations \\
\hline
回答市民的问题 & huídá shìmín de wèntí & locution verbale & répondre aux questions des habitants \\
\hline
办公时间 & bàngōng shíjiān & nom & horaires d'ouverture \\
\hline
地址 & dìzhǐ & nom & adresse \\
\hline
电话 & diànhuà & nom & téléphone \\
\hline
\end{tabularx}
\end{center}

\begin{attention}[Collocations : ne traduisez pas mot à mot]
En chinois, chaque nom a son verbe attitré. On dit \zh{办理证明} (établir un document), jamais \zh{做证明} ; on dit \zh{收取税费} ou \zh{征收税费} (percevoir une taxe), jamais \zh{拿税}. Une mauvaise collocation coûte des points de lexique au Bac. Apprenez les couples verbe + nom comme des blocs : \zh{办理 + 证件}, \zh{收取/征收 + 费用/税}, \zh{接受 + 投诉/建议}, \zh{提供 + 服务/信息}.
\end{attention}

\courssub{Étape 2 : remplir le plan en cinq blocs}

Avant de rédiger, remplissez la fiche-plan suivante. Un plan rempli, c'est la moitié des points de structure déjà gagnés.

\begin{center}
\begin{tabularx}{\linewidth}{|X|X|X|}
\hline
\rowcolor{popblueL} Bloc & Contenu attendu & Outils linguistiques \\
\hline
1. Titre & 社区服务通知 (avis de services) & nom du service + 通知 \\
\hline
2. Objet & pourquoi cet avis & 为了 + verbe + ，本社区提供以下服务： \\
\hline
3. Services & 3 services numérotés & 一、…；二、…；三、…。 \\
\hline
4. Infos pratiques & horaires, adresse, téléphone & 办公时间：… 地址：… 电话：… \\
\hline
5. Clôture & formule de politesse & 谢谢您的合作！/ 欢迎大家提出意见。 \\
\hline
\end{tabularx}
\end{center}

\begin{exemplebox}[Deux modèles de clôture à réutiliser]
\zh{谢谢您的合作！} \emph{Xièxie nín de hézuò!} « Merci de votre coopération ! »\\
\zh{欢迎大家提出意见。} \emph{Huānyíng dàjiā tíchū yìjiàn.} « Nous accueillons vos remarques. » (litt. « Bienvenue à tous pour proposer des avis. »)
\end{exemplebox}

\begin{methode}[Compter ses caractères (80 à 120)]
\begin{enumerate}
\item Rédigez d'abord sans compter, en suivant le plan en cinq blocs.
\item Comptez ensuite \textbf{uniquement les caractères chinois}, sans la ponctuation ni les chiffres arabes.
\item Si vous êtes sous 80 : ajoutez un service ou une précision (un horaire, un guichet). Un avis trop court perd des points de \textbf{structure}.
\item Si vous dépassez 120 : supprimez les répétitions et les adjectifs inutiles. Un avis trop long \textbf{multiplie les occasions de fautes} de caractères.
\item Visez la zone idéale : 95 à 110 caractères.
\end{enumerate}
\end{methode}

\courssub{Étapes 3 et 4 : rédiger, puis relire en binôme}

À l'étape 3, rédigez au propre en appliquant l'ordre des traits (haut → bas, gauche → droite, horizontal avant vertical, extérieur avant intérieur pour les caractères fermés comme \zh{国}). Surveillez les radicaux du thème : la clé de la parole \zh{讠} dans \zh{诉} (plainte) et \zh{议} (avis), la clé du cœur \zh{心} dans \zh{意} (idée), la clé de la main \zh{扌} dans \zh{接} (recevoir) et \zh{投} (déposer).

À l'étape 4, échangez votre copie avec un voisin (\emph{pair review}) : chacun corrige l'avis de l'autre avec la grille ci-dessous, signe ses remarques au crayon, puis rend la copie à son auteur pour version finale.

\begin{attention}[Ponctuation chinoise obligatoire]
La ponctuation chinoise est \cle{pleine chasse} et fait partie de la note : \zh{。} (point), \zh{，} (virgule), \zh{、} (virgule d'énumération entre mots), \zh{；} (point-virgule entre services numérotés), \zh{：} (deux-points). Les deux-points \zh{：} introduisent obligatoirement les listes et les horaires : \zh{提供以下服务：} et \zh{办公时间：}. Écrire un point français « . » à la place de \zh{。} est une faute de présentation sanctionnée.
\end{attention}

\begin{popfigure}
\begin{tikzpicture}[node distance=0.9cm and 1.1cm,
  etape/.style={rectangle, rounded corners, draw=popblue, fill=popblueL, thick, align=center, minimum width=2.6cm, minimum height=1.1cm, font=\small},
  fleche/.style={-{Stealth[length=3mm]}, thick, popdark}]
\node[etape, align=center] (e1){\textbf{Étape 1}\\Brainstorming\\lexique + glossaire};
\node[etape, right=of e1, draw=poporange, fill=poporangeL, align=center] (e2){\textbf{Étape 2}\\Plan en\\5 blocs};
\node[etape, right=of e2, draw=popgreen, fill=popgreenL, align=center] (e3){\textbf{Étape 3}\\Rédaction\\80--120 caractères};
\node[etape, right=of e3, draw=poppurple, fill=poppurpleL, align=center] (e4){\textbf{Étape 4}\\Relecture croisée\\(pair review)};
\draw[fleche] (e1) -- (e2);
\draw[fleche] (e2) -- (e3);
\draw[fleche] (e3) -- (e4);
\end{tikzpicture}
\legende{Organigramme des quatre étapes de la production écrite.}
\end{popfigure}

\courssub{Grille d'évaluation sur 20 points}

\begin{propriete}[Ce que le correcteur regarde au Bac]
Au baccalauréat, une production écrite de ce type est notée sur trois piliers : l'\cle{exactitude des caractères} (traits et radicaux corrects, pas de caractère inventé), le \cle{respect de la structure} attendue (ici les cinq blocs de l'avis) et la \cle{richesse du lexique du thème} (collocations correctes : \zh{办理}、\zh{收取}、\zh{接受}、\zh{提供}). La cohérence (connecteurs, ordre logique) et la présentation (ponctuation chinoise, propreté) complètent la note. Entraînez-vous avec la grille officielle de la classe :
\end{propriete}

\begin{popfigure}
\begin{tikzpicture}
\matrix[column sep=0pt, row sep=0pt,
  nodes={draw=popline, align=left, text=popdark, font=\small, inner sep=5pt}] (m) {
  \node[fill=popblue, text=white, font=\small\bfseries, minimum width=7.2cm, minimum height=0.7cm] {Critère}; &
  \node[fill=popblue, text=white, font=\small\bfseries, minimum width=4.6cm, minimum height=0.7cm, align=center] {Barème}; \\
  \node[fill=popcream, minimum width=7.2cm, minimum height=0.8cm] {1. Structure complète (5 blocs)}; &
  \node[fill=popcream, minimum width=4.6cm, minimum height=0.8cm, align=center] {5 points}; \\
  \node[fill=popblueL, minimum width=7.2cm, minimum height=0.8cm] {2. Lexique et collocations corrects}; &
  \node[fill=popblueL, minimum width=4.6cm, minimum height=0.8cm, align=center] {5 points}; \\
  \node[fill=popcream, minimum width=7.2cm, minimum height=0.8cm] {3. Caractères exacts (traits, radicaux)}; &
  \node[fill=popcream, minimum width=4.6cm, minimum height=0.8cm, align=center] {5 points}; \\
  \node[fill=popblueL, minimum width=7.2cm, minimum height=0.8cm] {4. Connecteurs et cohérence}; &
  \node[fill=popblueL, minimum width=4.6cm, minimum height=0.8cm, align=center] {3 points}; \\
  \node[fill=popcream, minimum width=7.2cm, minimum height=0.8cm] {5. Présentation et ponctuation chinoise}; &
  \node[fill=popcream, minimum width=4.6cm, minimum height=0.8cm, align=center] {2 points}; \\
  \node[fill=popgold, minimum width=7.2cm, minimum height=0.8cm, font=\small\bfseries] {Total}; &
  \node[fill=popgold, minimum width=4.6cm, minimum height=0.8cm, align=center, font=\small\bfseries] {20 points}; \\
};
\end{tikzpicture}
\legende{Grille d'évaluation de l'avis d'offre de services (sur 20).}
\end{popfigure}

\begin{exoresolu}[Corriger un avis d'élève]
Énoncé : voici l'avis rédigé par un élève. Repérez les \textbf{quatre fautes}, corrigez-les et attribuez une note indicative sur 20.\\
\zh{为了方便市民，我们社区提供以下服务，一、做出生证明；二、收取地方税；三、接受投诉. 办公时间，周一到周五上午八点到下午四点。谢谢您的合作！}
\tcblower
Solution détaillée :
\begin{enumerate}
\item \zh{以下服务，} → \zh{以下服务：} : la liste numérotée doit être introduite par les deux-points \zh{：}, pas par une virgule.
\item \zh{做出生证明} → \zh{办理出生证明} : mauvaise collocation ; un document officiel se « traite » avec \zh{办理}, pas avec \zh{做}.
\item \zh{接受投诉.} → \zh{接受投诉。} : point français interdit, il faut le point chinois pleine chasse \zh{。}.
\item \zh{办公时间，} → \zh{办公时间：} : les horaires s'introduisent aussi avec \zh{：}.
\end{enumerate}
Note indicative : structure complète (5/5) ; lexique entaché d'une faute de collocation (3/5) ; caractères tous corrects (5/5) ; cohérence bonne (3/3) ; présentation pénalisée par deux fautes de ponctuation (1/2). Total : 17/20.
\end{exoresolu}

\courssub{Variante pour l'examen : la lettre au maire}

Au baccalauréat, le même lexique peut être réinvesti dans une \cle{lettre adressée au maire} pour demander un service. Trois changements seulement par rapport à l'avis : une formule d'appel (\zh{尊敬的市长先生：} « Monsieur le Maire, »), une demande polie avec \zh{我想申请…} « je souhaite demander… » ou \zh{请您帮我办理…} « veuillez m'aider à établir… », et une salutation finale (\zh{此致敬礼！} « Veuillez agréer mes salutations distinguées. », suivi du nom et de la date alignés à droite).

\begin{exemplebox}[Mini-lettre de demande de service]
\zh{尊敬的市长先生：您好！我想申请出生证明。请您告诉我需要什么材料。谢谢您的帮助！}\\
\emph{Zūnjìng de shìzhǎng xiānsheng: nín hǎo! Wǒ xiǎng shēnqǐng chūshēng zhèngmíng. Qǐng nín gàosu wǒ xūyào shénme cáiliào. Xièxie nín de bāngzhù!}\\
« Monsieur le Maire, bonjour ! Je souhaite demander un acte de naissance. Veuillez m'indiquer les pièces nécessaires. Merci de votre aide ! »
\end{exemplebox}

\begin{experience}[Atelier d'écriture chronométré (30 minutes)]
\begin{enumerate}
\item \textbf{5 min} — En binôme, brainstormez les services et remplissez le plan en cinq blocs sur une feuille de brouillon.
\item \textbf{15 min} — Rédigez individuellement votre avis pour la communauté urbaine de Douala (ou le council de Bamenda), entre 80 et 120 caractères.
\item \textbf{5 min} — Relecture croisée avec la grille d'évaluation ; notez la copie de votre voisin et justifiez chaque point retiré.
\item \textbf{5 min} — Version finale au propre ; les deux meilleures copies sont lues à voix haute et affichées au mur de la classe.
\end{enumerate}
\end{experience}

\begin{savaistu}[Les guichets uniques et la coopération sino-camerounaise]
Au Cameroun, les communautés urbaines de Douala et de Yaoundé ont mis en place des \cle{guichets uniques} où l'usager accomplit plusieurs démarches (état civil, taxes, permis) au même endroit — exactement le type de service décrit dans votre avis. Dans le cadre de la coopération sino-camerounaise (infrastructures, formation, échanges municipaux), savoir rédiger un avis administratif en chinois n'est pas un exercice artificiel : c'est une \cle{mise en situation professionnelle réaliste}, utile pour un futur emploi dans une administration, une entreprise chinoise implantée au Cameroun ou un organisme de coopération.
\end{savaistu}

\begin{exercice}[À faire à la maison]
\begin{enumerate}
\item Rédigez la version définitive de votre avis (80 à 120 caractères) en intégrant les corrections de la relecture croisée.
\item Transformez votre avis en lettre au maire demandant l'établissement d'une carte d'identité (60 à 90 caractères), avec appel, demande polie et salutation finale.
\item Version : traduisez en français — \zh{本社区办公时间是周一到周五，上午八点到下午四点。欢迎大家提出意见。}
\end{enumerate}
\end{exercice}

\begin{corrige}[Exercice « À faire à la maison »]
\begin{enumerate}
\item Production individuelle : vérifier les cinq blocs, la fourchette 80--120 caractères, les collocations (\zh{办理}、\zh{收取}、\zh{接受}) et la ponctuation chinoise.
\item Modèle possible : \zh{尊敬的市长先生：您好！我是本社区的市民。我想办理身份证，请您告诉我需要什么材料和费用。谢谢您的帮助！此致敬礼！} \emph{Zūnjìng de shìzhǎng xiānsheng: nín hǎo! Wǒ shì běn shèqū de shìmín. Wǒ xiǎng bànlǐ shēnfènzhèng, qǐng nín gàosu wǒ xūyào shénme cáiliào hé fèiyòng. Xièxie nín de bāngzhù! Cǐ zhì jìnglǐ!} « Monsieur le Maire, bonjour ! Je suis un habitant de votre commune. Je souhaite établir une carte d'identité ; veuillez m'indiquer les pièces et les frais nécessaires. Merci de votre aide ! Veuillez agréer mes salutations distinguées. »
\item « Les horaires d'ouverture de notre communauté sont du lundi au vendredi, de 8 h à 16 h. Nous accueillons vos remarques. »
\end{enumerate}
\end{corrige}

\newpage

\coursec{Activité d'intégration}

\coursec{Leçon 19 — Expression écrite : offre de services aux usagers des administrations (collectivités locales)}

\courssub{Objectifs de la leçon}

À l'issue de cette leçon, tu seras capable de :
\begin{itemize}
\item produire un \cle{avis officiel d'offre de services} (\zh{公告}) complet en chinois simplifié (80 à 120 caractères) pour une collectivité locale (mairie, \emph{council}), en respectant la structure administrative codifiée ;
\item écrire correctement les \cle{caractères du thème} (traits, radicaux, différenciation) : \zh{服务、办理、证明、材料、窗口、综合、设立、地址、电话} ;
\item mobiliser les \cle{connecteurs et formules figées} de l'écrit administratif : \zh{为了}、\zh{第一/第二/第三}、\zh{至}、\zh{欢迎……前来办理}、\zh{特此公告} ;
\item réaliser un \cle{thème} (français $\to$ chinois) et une \cle{version} (chinois $\to$ français) de phrases et textes administratifs simples ;
\item comprendre et poser une \cle{question d'usager} à l'oral (\zh{请问，……需要什么材料？}) et y répondre.
\end{itemize}

\courssub{Découverte : un avis de services à la mairie}

\courssub{Situation de communication}

La mairie de ta ville (ou le \emph{council} de ta région anglophone) ouvre un \cle{guichet unique} (\zh{综合服务窗口}) pour rapprocher l'administration des habitants. Le maire souhaite informer la communauté chinoise locale — commerçants, ingénieurs, étudiants — ainsi que les partenaires chinois de la commune. Tu dois rédiger l'\cle{avis officiel en chinois} affiché à l'entrée du guichet et au marché central.

\begin{textebox}[Note de service du maire — Support authentique bilingue]
市长办公室通知。\\
Shìzhǎng bàngōngshì tōngzhī.\\
« Note du bureau du maire. »\\[4pt]
为了方便市民，本市政府设立综合服务窗口，提供以下服务：\\
Wèile fāngbiàn shìmín, běn shì zhèngfǔ shèlì zònghé fúwù chuāngkǒu, tígōng yǐxià fúwù:\\
« Pour faciliter la vie des habitants, la mairie crée un guichet unique offrant les services suivants : »\\[4pt]
一、办理出生证明；二、办理居住证；三、咨询税务问题。\\
Yī, bànlǐ chūshēng zhèngmíng; èr, bànlǐ jūzhùzhèng; sān, zīxún shuìwù wèntí.\\
« 1. Délivrance des actes de naissance ; 2. Délivrance des cartes de séjour ; 3. Renseignements sur les impôts. »\\[4pt]
办公时间：星期一至星期五，上午八点到下午四点。地址：市政大楼一层。电话：二三三 六五四 三二一。\\
Bàngōng shíjiān: xīngqīyī zhì xīngqīwǔ, shàngwǔ bā diǎn dào xiàwǔ sì diǎn. Dìzhǐ: shìzhèng dàlóu yī céng. Diànhuà: èr sān sān, liù wǔ sì, sān èr yī.\\
« Horaires : lundi au vendredi, 8 h à 16 h. Adresse : rez-de-chaussée de l'hôtel de ville. Tél. : 233 654 321. »
\end{textebox}

\begin{exercice}[Repérage guidé (10 min)]
Sur ce texte, identifie et surligne :
\begin{enumerate}
\item Le mot qui introduit l'objectif (\zh{为了}).
\item Les trois services (verbes + objets).
\item L'expression des horaires (\zh{星期一至星期五}, \zh{上午八点到下午四点}).
\item L'adresse (\zh{地址：……}) et le téléphone (\zh{电话：……}).
\end{enumerate}
Vérifie en binôme que chaque information est comprise.
\end{exercice}

\begin{corrige}[Exercice 1]
\begin{itemize}
\item Objectif : \zh{为了方便市民} (Pour faciliter les habitants).
\item Services : \zh{办理出生证明} (traiter acte naissance), \zh{办理居住证} (traiter carte séjour), \zh{咨询税务问题} (consulter questions fiscales).
\item Horaires : \zh{办公时间：星期一至星期五，上午八点到下午四点}.
\item Lieu/Contact : \zh{地址：市政大楼一层。电话：二三三 六五四 三二一}.
\end{itemize}
\end{corrige}

\courssub{Lexique des services aux usagers}

\begin{center}
\begin{tabularx}{\linewidth}{|X|X|X|X|}
\hline
\rowcolor{popblueL} Caractères & Pinyin & Nature & Traduction \\
\hline
公告 & gōnggào & n. & avis officiel, annonce administrative \\
\hline
服务 & fúwù & n./v. & service ; servir, assurer un service \\
\hline
窗口 & chuāngkǒu & n. & guichet, fenêtre (administratif) \\
\hline
综合服务窗口 & zònghé fúwù chuāngkǒu & n. & guichet unique / guichet de services intégrés \\
\hline
办理 & bànlǐ & v. & traiter, effectuer (une démarche administrative) \\
\hline
设立 & shèlì & v. & créer, établir (une structure) \\
\hline
提供 & tígōng & v. & fournir, offrir (un service) \\
\hline
出生证明 & chūshēng zhèngmíng & n. & acte de naissance \\
\hline
居住证 & jūzhùzhèng & n. & carte de séjour, certificat de résidence \\
\hline
建筑许可证 & jiànzhú xǔkězhèng & n. & permis de construire \\
\hline
市场税 & shìchǎng shuì & n. & taxes de marché \\
\hline
咨询 & zīxún & v. & consulter, demander des renseignements \\
\hline
材料 & cáiliào & n. & documents, pièces justificatives, dossier \\
\hline
身份证 & shēnfènzhèng & n. & carte d'identité \\
\hline
照片 & zhàopiàn & n. & photo d'identité \\
\hline
市民 & shìmín & n. & habitants, citoyens (ville) \\
\hline
广大市民 & guǎngdà shìmín & n. & l'ensemble des habitants (formule admin.) \\
\hline
办公时间 & bàngōng shíjiān & n. & horaires d'ouverture (bureau) \\
\hline
至 & zhì & prep. & jusqu'à (plus formel que 到) \\
\hline
地址 & dìzhǐ & n. & adresse \\
\hline
市政大楼 & shìzhèng dàlóu & n. & hôtel de ville, bâtiment administratif \\
\hline
一层 & yī céng & n. & rez-de-chaussée (1er étage chinois = RDC) \\
\hline
电话 & diànhuà & n. & téléphone \\
\hline
欢迎 & huānyíng & v. & accueillir, souhaiter la bienvenue \\
\hline
前来 & qiánlái & v. & venir (vers le locuteur, formel) \\
\hline
特此公告 & tècǐ gōnggào & expr. & « À cet effet, le présent avis » (clôture admin.) \\
\hline
\end{tabularx}
\end{center}

\begin{exercice}[Entraînement lexical]
Traduis en chinois (caractères + pinyin) :
\begin{enumerate}
\item Guichet unique
\item Délivrer un acte de naissance
\item Pièces justificatives
\item Du lundi au vendredi
\item Rez-de-chaussée de l'hôtel de ville
\item Bienvenue pour venir effectuer les démarches
\end{enumerate}
\end{exercice}

\begin{corrige}[Exercice 2]
\begin{enumerate}
\item \zh{综合服务窗口} (zònghé fúwù chuāngkǒu)
\item \zh{办理出生证明} (bànlǐ chūshēng zhèngmíng)
\item \zh{材料} (cáiliào) ou \zh{所需材料} (suǒxū cáiliào)
\item \zh{星期一至星期五} (xīngqīyī zhì xīngqīwǔ)
\item \zh{市政大楼一层} (shìzhèng dàlóu yī céng)
\item \zh{欢迎前来办理} (huānyíng qiánlái bànlǐ)
\end{enumerate}
\end{corrige}

\courssub{Écriture des caractères du thème}

\begin{methode}[Écrire correctement les caractères administratifs clés]
\begin{enumerate}
\item \textbf{Analyser le radical (clé) et la structure} : identifier la clé sémantique (souvent à gauche ou en haut) et la partie phonétique.
\item \textbf{Respecter l'ordre des traits} : horizontal avant vertical ; gauche avant droite ; haut avant bas ; extérieur avant intérieur ; trait central avant les ailes ; point final en dernier.
\item \textbf{Vérifier la proportion} : le caractère doit tenir dans un carré imaginaire, parties équilibrées.
\item \textbf{S'exercer} : tracer 3 fois chaque caractère sur papier quadrillé (田字格), en prononçant le pinyin et le sens.
\end{enumerate}
\end{methode}

\begin{propriete}[Analyse des caractères difficiles du thème]
\begin{center}
\begin{tabularx}{\linewidth}{|X|X|X|X|X|}
\hline
\rowcolor{popblueL} Caractère & Pinyin & Radical (Clé) & Nb traits & Ordre des traits principal / Pièges \\
\hline
服 & fú & \zh{月} (viande/chair) & 8 & Haut \zh{月} (4 traits) + bas \zh{彐} + \zh{ㄗ} (4 traits). Ne pas confondre avec \zh{福} (fú, bonheur, radical \zh{礻}). \\
\hline
办 & bàn & \zh{力} (force) & 8 & Structure gauche-droite. \zh{力} à droite (2 traits). Partie gauche \zh{卩} + \zh{日} + \zh{丨}. Attention : \zh{办} (simplifié) $\neq$ \zh{辦} (traditionnel). \\
\hline
证 & zhèng & \zh{讠} (parole) & 6 & Gauche \zh{讠} (2 traits : point, trait horizontal coudé). Droite \zh{正} (5 traits : horizontal, vertical, horizontal, vertical, horizontal). Ne pas écrire \zh{正明}. \\
\hline
材 & cái & \zh{木} (bois) & 7 & Haut \zh{木} (4 traits). Bas \zh{才} (3 traits : horizontal, oblique, point). Sens : matière première, bois $\to$ documents. \\
\hline
窗 & chuāng & \zh{穴} (caverne/toit) & 12 & Haut \zh{穴} (5 traits). Bas \zh{囱} (7 traits). Ne pas oublier le point en haut de \zh{穴}. \\
\hline
综 & zòng & \zh{纟} (soie) & 14 & Gauche \zh{纟} (3 traits). Droite \zh{东} (5 traits) + \zh{心} (4 traits) simplifié. Complexe : écrire grand. \\
\hline
设 & shè & \zh{讠} (parole) & 11 & Gauche \zh{讠}. Droite \zh{殳} (arme) + \zh{口}. Attention à l'ordre dans \zh{殳}. \\
\hline
地 & dì & \zh{土} (terre) & 6 & Gauche \zh{土} (3 traits). Droite \zh{也} (3 traits). Différencier \zh{地} (lieu, adv -de) et \zh{的} (possessif). \\
\hline
话 & huà & \zh{讠} (parole) & 13 & Gauche \zh{讠}. Droite \zh{舌} (6 traits) + \zh{刂} (2 traits). \zh{电话} (téléphone = parole électrique). \\
\hline
\end{tabularx}
\end{center}
\end{propriete}

\begin{attention}[Caractères proches / Pièges visuels]
\begin{itemize}
\item \zh{证} (zhèng, certificat, radical \zh{讠}) $\neq$ \zh{正} (zhèng, correct, radical \zh{止}) $\neq$ \zh{政} (zhèng, politique, radical \zh{攵}).
\item \zh{材} (cái, matériaux) $\neq$ \zh{才} (cái, talent, seulement).
\item \zh{窗} (chuāng, fenêtre/guichet) $\neq$ \zh{床} (chuáng, lit, radical \zh{广}).
\item \zh{层} (céng, étage/couche) $\neq$ \zh{曾} (céng, autrefois, radical \zh{曰}).
\item \zh{特} (tè, spécial) $\neq$ \zh{持} (chí, tenir, radical \zh{手}) $\neq$ \zh{时} (shí, temps, radical \zh{日}).
\end{itemize}
\end{attention}

\begin{exercice}[Entraînement à l'écriture]
Sur ton cahier d'écriture (田字格), écris 5 fois chaque caractère ci-dessous en respectant l'ordre des traits. À côté, note le pinyin, le radical et le sens.
\zh{服、办、证、材、窗、综、设、地、话、层}
\end{exercice}

\courssub{Structure et connecteurs de l'offre de services}

\begin{aretenir}[Structure type de l'avis officiel \zh{公告} (Modèle canonique)]
\zh{【标题】} \zh{XXXX公告} (Nom administration + 公告)\\
\zh{【缘由】} \zh{为了 + Objectif，+ Sujet + 设立/提供……} \\
\zh{【内容】} \zh{第一，……；第二，……；第三，……。} (Verbes d'action : \zh{办理}、\zh{咨询}、\zh{受理})\\
\zh{【时限】} \zh{办公时间：时间范围 (从……到…… / ……至……)} \\
\zh{【地点】} \zh{地址：具体地点} \\
\zh{【联系】} \zh{电话：号码 (chiffre par chiffre, 1 = yāo)} \\
\zh{【结语】} \zh{欢迎广大市民前来办理。特此公告。} \\
\zh{【落款】} \zh{发布单位，日期 (chiffres chinois : 二〇二五年五月)}
\end{aretenir}

\begin{propriete}[Points grammaticaux essentiels pour l'avis administratif]
\begin{enumerate}
\item \textbf{\cle{为了} + Objectif (Finalité) en tête de phrase} \\
Structure : \zh{为了 + Verbe/Objectif，+ Sujet + Verbe principal}. \\
Ex. \zh{为了方便市民，本市政府设立了窗口。} (Pour faciliter les habitants, la mairie a créé un guichet.) \\
\emph{Remarque :} Pas de \zh{是} (shì) devant le verbe principal. \textbf{Erreur :} \zh{我们是提供服务} $\to$ \textbf{Correct :} \zh{我们提供服务}.
\item \textbf{Énumération formelle : \zh{第一}、\zh{第二}、\zh{第三}} \\
Suivis d'une proposition verbe + objet. Séparés par point-virgule \zh{；}, fin par point \zh{。}. \\
Ex. \zh{第一，办理出生证明；第二，办理居住证；第三，咨询税务问题。}
\item \textbf{Ordre des mots : Temps -- Lieu -- Verbe (T-L-V)} \\
Sujet + Temps + Lieu + Verbe (+ Objet). \\
Ex. \zh{我们 星期一 上午八点 在市政大楼 开门。} (Nous lundi 8h à l'hôtel de ville ouvrons.)
\item \textbf{Expression de la durée / intervalle : \zh{从……到……} ou \zh{……至……}} \\
\zh{从上午八点到下午四点} (oral/standard) ; \zh{上午八点至下午四点} (écrit/admin, plus concis). \\
\zh{星期一至星期五} (lundi à vendredi).
\item \textbf{Formule de clôture figée : \zh{特此公告}} \\
\zh{特} (spécialement) \zh{此} (ceci) \zh{公告} (annoncer). Ne se traduit pas mot à mot. Équivalent : « Le présent avis est publié à cet effet. » / « Fait pour être affiché et publié. »
\item \textbf{Date administrative : chiffres chinois formels} \\
\zh{零 (líng)}, \zh{一 (yī)}, \zh{二 (èr)}, \zh{三 (sān)}, \zh{四 (sì)}, \zh{五 (wǔ)}, \zh{六 (liù)}, \zh{七 (qī)}, \zh{八 (bā)}, \zh{九 (jiǔ)}, \zh{十 (shí)}. \\
Année 2025 : \zh{二〇二五年} (èr líng èr wǔ nián). Mois : \zh{五月} (wǔ yuè).
\end{enumerate}
\end{propriete}

\begin{savaistu}[Pratiques administratives : Cameroun vs Chine]
\begin{itemize}
\item \textbf{Cameroun :} Les avis de mairie (communes, \emph{councils}) sont souvent en français/anglais, affichés sur panneaux extérieurs. Structure : Titre, Objet, Services, Horaires, Signature du Maire/Secretary General. Moins de formules figées codifiées.
\item \textbf{Chine :} L'avis (\zh{公告} / \zh{通知}) suit un formalisme strict hérité de la tradition impériale et du style partisan. La formule \zh{特此公告} (ou \zh{特此通知}) est \textbf{obligatoire} en bas à droite avant la signature/date. L'usage de \zh{至} (jusqu'à) et \zh{前来} (venir) marque le registre soutenu (\zh{书面语}).
\item \textbf{Point commun :} L'affichage public vise l'accessibilité de l'information (horaires, lieu, pièces à fournir). Au Cameroun comme en Chine, la numérisation progresse (sites web, applications \zh{政务服务网}), mais l'affichage physique reste légal.
\end{itemize}
\end{savaistu}

\courssub{Thème et version guidés (Entraînement individuel)}

\begin{methode}[Méthode du Thème (Français $\to$ Chinois) : Phrase administrative]
\begin{enumerate}
\item \textbf{Analyser} : Identifier le sujet, le temps, le lieu, l'action, le but.
\item \textbf{Structurer} : Appliquer l'ordre Sujet -- Temps -- Lieu -- Verbe -- Objet. Placer \zh{为了……} au début si but.
\item \textbf{Lexicaliser} : Choisir le verbe administratif précis (\zh{办理} pour démarche, \zh{咨询} pour renseignement, \zh{申请} pour demande, \zh{缴纳} pour paiement).
\item \textbf{Vérifier} : Caractères (radicaux), pinyin (tons), ponctuation chinoise (， 。 ；), connecteurs (\zh{第一}、\zh{至}、\zh{特此公告}).
\end{enumerate}
\end{methode}

\begin{methode}[Méthode de la Version (Chinois $\to$ Français) : Style administratif]
\begin{enumerate}
\item \textbf{Segmenter} : Repérer les blocs (Titre, 为了, 第一/二/三, 时间, 地址, 电话, 特此公告, 落款).
\item \textbf{Traduire le sens, pas le mot} : \zh{特此公告} $\to$ « Le présent avis est publié à cet effet » (pas « Spécialement ceci annoncer »).
\item \textbf{Adapter le registre} : \zh{办理} $\to$ « délivrer / traiter / effectuer » ; \zh{市民} $\to$ « les habitants / les usagers / la population » ; \zh{前来} $\to$ « se présenter / se rendre ».
\item \textbf{Soigner la mise en page} : Respecter la typographie française (majuscules, deux-points, point-virgules).
\end{enumerate}
\end{methode}

\begin{exoresolu}[Thème guidé : Phrases isolées]
Traduis en chinois (caractères + pinyin).
\begin{enumerate}
\item Pour faciliter les usagers, la mairie de Bafoussam a créé un guichet unique.
\item 1. Traiter les cartes de séjour ; 2. Renseigner sur les taxes de marché.
\item Horaires : du mardi au samedi, de 7 h 30 à 15 h 30.
\item Adresse : 2e étage de l'hôtel de ville. Tél. : 233 6 99 88 77.
\item Bienvenue aux habitants pour venir faire les démarches. Le présent avis est publié à cet effet.
\end{enumerate}
\tcblower
\textbf{Corrigé commenté :}
\begin{enumerate}
\item \zh{为了方便市民，巴富萨市政府设立了综合服务窗口。} (Wèile fāngbiàn shìmín, Bāfùsà shì zhèngfǔ shèlìle zònghé fúwù chuāngkǒu.) \emph{Le 了 marque l'accomplissement de la création.}
\item \zh{第一，办理居住证；第二，咨询市场税问题。} (Dìyī, bànlǐ jūzhùzhèng; dì'èr, zīxún shìchǎng shuì wèntí.)
\item \zh{办公时间：星期二至星期六，上午七点半至下午三点半。} (Bàngōng shíjiān: xīngqī'èr zhì xīngqīliù, shàngwǔ qī diǎn bàn zhì xiàwǔ sān diǎn bàn.) \emph{半 (bàn) pour "demi".}
\item \zh{地址：市政大楼二层。电话：二三三 六 九九 八八 七七。} (Dìzhǐ: shìzhèng dàlóu èr céng. Diànhuà: èr sān sān liù jiǔ jiǔ bā bā qī qī.) \emph{Chiffres lus un à un.}
\item \zh{欢迎广大市民前来办理。特此公告。} (Huānyíng guǎngdà shìmín qiánlái bànlǐ. Tècǐ gōnggào.)
\end{enumerate}
\end{exoresolu}

\begin{exoresolu}[Version guidée : Court texte]
Traduis en français (style administratif) :
\zh{雅温得市政府公告。为了方便广大市民，本市政府设立综合服务窗口，提供以下服务：第一，办理出生证明；第二，申请建筑许可证；第三，缴纳市场税。办公时间：星期一至星期五，上午八点至下午四点。地址：市政大楼一层。电话：二三三 六五五 四四三三。欢迎广大市民前来办理。特此公告。雅温得市政府，二〇二五年六月。}
\tcblower
\textbf{Traduction modèle :}
\textbf{Avis de la Mairie de Yaoundé.}
Pour faciliter la vie de l'ensemble des habitants, la mairie met en place un guichet unique offrant les services suivants :
1. Délivrance des actes de naissance ;
2. Demande de permis de construire ;
3. Paiement des taxes de marché.
Horaires d'ouverture : du lundi au vendredi, de 8 h à 16 h.
Adresse : Rez-de-chaussée de l'Hôtel de ville.
Téléphone : 233 6 55 44 33.
Tous les habitants sont invités à se présenter pour effectuer leurs démarches.
Le présent avis est publié à cet effet.
\textbf{Mairie de Yaoundé, juin 2025.}
\end{exoresolu}

\courssub{Production guidée et évaluation : Activité d'intégration (2 h)}

\courssub{Consignes de l'activité finale}

Travail en groupes de 3 ou 4 élèves. Durée totale : 2 heures.

\begin{enumerate}
\item \textbf{Choix de la collectivité et adaptation des services (10 min).} Choisissez une collectivité réelle : Mairie de Yaoundé, Douala, Bafoussam, ou \emph{Council} de Buea, Bamenda. Adaptez les 3 services (ex. : \zh{建筑许可证} permis de construire, \zh{市场税} taxes marché, \zh{营业执照} licence commerce).
\item \textbf{Rédaction au brouillon (25 min).} Rédigez l'avis (80--120 caractères) en respectant \textbf{strictement} la structure de la boîte \textbf{À retenir} (p. 3) : Titre \zh{公告} $\to$ \zh{为了……} $\to$ \zh{第一/二/三} $\to$ \zh{办公时间} $\to$ \zh{地址} $\to$ \zh{电话} $\to$ \zh{欢迎……前来办理。特此公告。} $\to$ Nom + Date.
\item \textbf{Relecture linguistique croisée (15 min).} Échangez les brouillons avec un autre groupe. Vérifiez :
\begin{itemize}
\item Ordre des mots (Temps/Lieu avant Verbe).
\item Caractères difficiles : \zh{服} (radical \zh{月}), \zh{证} (radical \zh{讠}), \zh{窗} (radical \zh{穴}), \zh{综} (radical \zh{纟}), \zh{设} (radical \zh{讠}), \zh{层} (ne pas confondre avec \zh{曾}).
\item Connecteurs : \zh{为了}、\zh{第一/二/三}、\zh{至}、\zh{特此公告}.
\item Absence de calque français (\zh{是} inutile, \zh{的} abusif).
\end{itemize}
\item \textbf{Production au propre et affichage (20 min).} Écrivez la version finale soignée (grand format A3) au feutre noir. Affichez au mur de la classe.
\item \textbf{Jeu de rôle : Usagers vs Administration (25 min).} Les groupes circulent. Les « usagers » lisent l'avis et posent \textbf{une question orale} en chinois au groupe « administration » (auteurs de l'avis).
\begin{itemize}
\item Question type : \zh{请问，办理……需要什么材料？} (Qǐngwèn, bànlǐ…… xūyào shénme cáiliào ?)
\item Réponse type : \zh{需要身份证、照片和……} (Xūyào shēnfènzhèng, zhàopiàn hé……)
\end{itemize}
\item \textbf{Évaluation par les pairs (15 min).} Remplissez la grille pour chaque avis affiché (sur 10 pts) :
\begin{center}
\begin{tabularx}{\linewidth}{|X|X|X|}
\hline
\rowcolor{popblueL} Critère & Points \\
\hline
Structure complète (7 étapes obligatoires) & /2 \\
\hline
Caractères corrects (traits, radicaux, pas de confusion) & /2 \\
\hline
Grammaire & connecteurs (为了, 第一, 至, 特此公告) & /2 \\
\hline
Clarté pour l'usager (infos lisibles, cohérentes) & /2 \\
\hline
Réponses orales aux questions (pertinence, prononciation) & /2 \\
\hline
\textbf{Total} & \textbf{/10} \\
\hline
\end{tabularx}
\end{center}
La classe désigne l'avis « \cle{Meilleur avis officiel} » (note max + meilleure clarté).
\item \textbf{Prolongement individuel : Version (15 min).} Sur ta copie, traduis en français \textbf{l'avis gagnant} (version complète, style administratif soigné). Remets à l'enseignant pour notation.
\end{enumerate}

\courssub{Proposition de production modèle et corrigé commenté}

\begin{exoresolu}[Avis modèle — Mairie de Douala (environ 110 caractères)]
\textbf{Production finale :}\\[4pt]
\zh{杜阿拉市政府公告}\\
\emph{Dù'ālā shì zhèngfǔ gōnggào.}\\
« Avis de la mairie de Douala. »\\[4pt]
\zh{为了方便广大市民，本市政府设立综合服务窗口，提供以下服务：}\\
\emph{Wèile fāngbiàn guǎngdà shìmín, běn shì zhèngfǔ shèlì zònghé fúwù chuāngkǒu, tígōng yǐxià fúwù:}\\
« Pour faciliter la vie de l'ensemble des habitants, la mairie crée un guichet unique offrant les services suivants : »\\[4pt]
\zh{第一，办理出生证明；第二，办理居住证；第三，咨询市场和税务问题。}\\
\emph{Dìyī, bànlǐ chūshēng zhèngmíng; dì'èr, bànlǐ jūzhùzhèng; dìsān, zīxún shìchǎng hé shuìwù wèntí.}\\
« 1. Délivrance des actes de naissance ; 2. Délivrance des cartes de séjour ; 3. Renseignements sur les taxes de marché et les impôts. »\\[4pt]
\zh{办公时间：星期一至星期五，上午八点至下午四点。}\\
\emph{Bàngōng shíjiān: xīngqīyī zhì xīngqīwǔ, shàngwǔ bā diǎn zhì xiàwǔ sì diǎn.}\\
« Horaires : du lundi au vendredi, de 8 h à 16 h. »\\[4pt]
\zh{地址：市政大楼一层。电话：二三三 四二 零九 八七。}\\
\emph{Dìzhǐ: shìzhèng dàlóu yī céng. Diànhuà: èr sān sān, sì èr, líng jiǔ, bā qī.}\\
« Adresse : rez-de-chaussée de l'hôtel de ville. Tél. : 233 42 09 87. »\\[4pt]
\zh{欢迎广大市民前来办理。特此公告。}\\
\emph{Huānyíng guǎngdà shìmín qiánlái bànlǐ. Tècǐ gōnggào.}\\
« Tous les habitants sont les bienvenus pour effectuer les démarches. Le présent avis est publié à cet effet. »\\[4pt]
\zh{杜阿拉市政府，二〇二五年五月}\\
\emph{Dù'ālā shì zhèngfǔ, èr líng èr wǔ nián wǔ yuè.}\\
« Mairie de Douala, mai 2025. »\\[6pt]
\textbf{Commentaire linguistique et stylistique :}
\begin{itemize}
\item \textbf{Titre} : \zh{杜阿拉市政府公告} (Nom administration + \zh{公告}), sans verbe, norme chinoise.
\item \textbf{Objet} : \zh{为了方便广大市民} (formule \zh{广大市民} = « l'ensemble des usagers »). Verbes \zh{设立} (créer structure) + \zh{提供} (offrir service).
\item \textbf{Services} : Variété verbale : \zh{办理} (démarche papier) pour 1 et 2 ; \zh{咨询} (renseignement) pour 3. Ponctuation \zh{；} entre items, \zh{。} final.
\item \textbf{Horaires} : \zh{至} (zhì) plus formel que \zh{到} (dào). \zh{上午八点至下午四点} concis.
\item \textbf{Adresse} : \zh{一层} (yī céng) = Rez-de-chaussée (Chine : 1er étage = RDC ; France : 1er étage = 1er au-dessus). Préciser si besoin \zh{底层} (dǐcéng) pour RDC, mais \zh{一层} standard admin.
\item \textbf{Téléphone} : Groupes de chiffres séparés par espaces en pinyin. \zh{一} (yī) se prononce souvent \emph{yāo} dans les numéros (évite confusion avec \zh{七} qī), mais s'écrit \zh{一}.
\item \textbf{Clôture} : \zh{欢迎广大市民前来办理} (invitation polie) + \zh{特此公告} (formule figée incontournable).
\item \textbf{Date} : Chiffres chinois formels \zh{二〇二五年五月} (zéro = \zh{〇} ou \zh{零}).
\end{itemize}
\end{exoresolu}

\begin{popfigure}
\begin{tikzpicture}[
  node distance=1.2cm and 1.5cm,
  box/.style={rectangle, draw=popblue, thick, fill=popblueL, rounded corners, align=center, minimum width=3.5cm, minimum height=1.2cm, font=\small},
  arrow/.style={-Stealth, thick, popdark},
  title/.style={font=\bfseries\large, text=popdark, align=center}
]
\node[title] (t) at (0,0) {Structure de l'Avis Officiel \zh{公告}};
\node[box, below=1.5cm of t] (titre) {Titre : \zh{XX公告}};
\node[box, below=0.8cm of titre, align=center] (objet){Objet : \zh{为了 + But，}\\\zh{Sujet + 设立/提供}};
\node[box, below=0.8cm of objet] (services) {Services : \zh{第一……} \zh{；} \zh{第二……} \zh{；} \zh{第三……}};
\node[box, below=0.8cm of services, align=center] (horaires){Horaires : \zh{办公时间：}\\\zh{从/至}};
\node[box, below=0.8cm of horaires] (lieu) {Lieu : \zh{地址：……}};
\node[box, below=0.8cm of lieu] (tel) {Contact : \zh{电话：……}};
\node[box, below=0.8cm of tel, align=center] (cloture){Clôture : \zh{欢迎前来办理。}\\\zh{特此公告。}};
\node[box, below=0.8cm of cloture] (signature) {Signature : \zh{单位，日期}};

\draw[arrow] (titre) -- (objet);
\draw[arrow] (objet) -- (services);
\draw[arrow] (services) -- (horaires);
\draw[arrow] (horaires) -- (lieu);
\draw[arrow] (lieu) -- (tel);
\draw[arrow] (tel) -- (cloture);
\draw[arrow] (cloture) -- (signature);
\end{tikzpicture}
\legende{Schéma de la structure obligatoire de l'avis administratif (flux vertical).}
\end{popfigure}

\courssub{Bilan de la leçon}

À la fin de cette leçon, vérifie que tu maîtrises les savoir-faire suivants, réinvestissables en expression écrite (épreuves du Baccalauréat) :

\begin{itemize}
\item \textbf{Rédiger un \zh{公告} complet} : Titre, \zh{为了} + but, 3 services (\zh{第一/二/三}), horaires (\zh{至}), adresse, téléphone, \zh{欢迎……前来办理。特此公告。}, signature/date.
\item \textbf{Écrire sans faute les 10 caractères cibles} : \zh{服、务、办、理、证、材、窗、综、设、层} (radicaux, traits, homophones).
\item \textbf{Manipuler la syntaxe admin.} : \zh{为了} en tête ; ordre T-L-V ; \zh{第一/二/三} + \zh{；} ; \zh{特此公告} figé.
\item \textbf{Traduire (Thème/Version)} : Passer du français au chinois (registre soutenu, vocabulaire précis \zh{办理/咨询/申请/缴纳}) et du chinois au français (style administratif, pas de mot-à-mot).
\item \textbf{Interagir à l'oral} : Poser \zh{请问，……需要什么材料？} et répondre \zh{需要……} avec vocabulaire adéquat (\zh{身份证、照片、表格}).
\item \textbf{Comparer cultures admin.} : Cameroun (affichage bilingue, signature hiérarchique) vs Chine (codification stricte \zh{特此公告}, date en chiffres chinois, \zh{前来}).
\end{itemize}

\begin{attention}[Derniers rappels pour l'examen]
\begin{itemize}
\item \textbf{Caractères} : \zh{证明} (pas \zh{正明}), \zh{窗口} (pas \zh{床口}), \zh{层} (pas \zh{曾}), \zh{特} (pas \zh{持}).
\item \textbf{Grammaire} : Pas de \zh{是} devant verbe d'action (\zh{我们提供} $\neq$ \zh{我们是提供}). Complément de temps/lieu \textbf{avant} le verbe.
\item \textbf{Ponctuation} : Utilise \zh{，} \zh{。} \zh{；} \zh{：} \zh{、} (pleine chasse). Pas d'espace après.
\item \textbf{Pinyin} : Espace entre mots (\zh{Duōshǎo qián}), majuscule nom propre (\zh{Dù'ālā}), \zh{ü} = \zh{ü} (lǜ, nǚ).
\end{itemize}
\end{attention}

\newpage

\coursec{Méthodes et exercices résolus}

\courssub{Méthodes et exercices résolus}

\begin{methode}[Rédiger un avis administratif d'offre de services (结构与连接词)]
\textbf{Objectif :} Produire un texte court, clair et formel informant les usagers des services disponibles à la mairie.
\textbf{Étapes :}
\begin{enumerate}
    \item \textbf{Analyser la consigne :} Identifier l'émetteur (mairie/collectivité), le destinataire (usagers/habitants), l'objet (services offerts) et le ton (formel, informatif, \cle{通知} ou \cle{公告}).
    \item \textbf{Structurer le texte selon le modèle standard :}
    \begin{itemize}
        \item \textbf{Titre :} Centré, police grande. Ex: \zh{关于提供便民服务的通知} (Guānyú tígōng biànmín fúwù de tōngzhī).
        \item \textbf{Formule d'appel :} \zh{尊敬的市民朋友们：} (Zūnjìng de shìmín péngyǒumen :) ou \zh{广大市民：}.
        \item \textbf{Corps du texte (3 parties) :}
        \begin{itemize}
            \item \textit{Contexte/But :} \zh{为进一步优化营商环境，方便市民办事……} (Wèi jìn yībù yōuhuà yíngshāng huánjìng, fāngbiàn shìmín bànshì…).
            \item \textit{Liste des services (avec connecteurs) :} \zh{现推出以下服务：} (Xiàn tuīchū yǐxià fúwù:) + énumération \zh{一、… 二、… 三、…} ou \zh{首先、其次、最后}.
            \item \textit{Modalités pratiques :} Lieu, horaires, contact, plateforme en ligne. \zh{办理地点：… 办理时间：… 咨询电话：…}.
        \end{itemize}
        \item \textbf{Formule de politesse finale :} \zh{特此通知。} (Tè cǐ tōngzhī.) + \textbf{Signature/Date} : \zh{XX市人民政府} (XX Shì Rénmín Zhèngfǔ) + \zh{202X年X月X日}.
    \end{itemize}
    \item \textbf{Mobiliser les connecteurs logiques (逻辑连接词) :}
    \begin{center}
    \begin{tabularx}{\linewidth}{|X|X|X|}\hline
    \rowcolor{popblueL} Fonction & Connecteur chinois (Pinyin) & Exemple d'usage \\ \hline
    But / Cause & \zh{为/为了} (wèi/wèile), \zh{由于} (yóuyú) & \zh{为了方便市民，我们延长办公时间。} \\ \hline
    Énumération & \zh{首先/第一} (shǒuxiān/dì yī), \zh{其次/第二} (qícì/dì èr), \zh{最后} (zuìhòu) & \zh{首先，可网上预约；其次，增设窗口。} \\ \hline
    Progressif & \zh{进而} (jìn'ér), \zh{进一步} (jìn yībù) & \zh{简化流程，进而提高效率。} \\ \hline
    Condition & \zh{如有需要} (rú yǒu xūyào), \zh{若} (ruò) & \zh{如有需要，请拨打热线。} \\ \hline
    Conclusion & \zh{特此} (tè cǐ), \zh{总之} (zǒngzhī) & \zh{特此通知。} \\ \hline
    \end{tabularx}
    \end{center}
    \item \textbf{Relire et vérifier :} Caractères (简体字), tons, ordre SUJET + TEMPS + LIEU + VERBE + OBJET, classificateurs (\cle{项} pour services, \cle{个} pour fenêtres), ponctuation pleine chasse (，。：).
\end{enumerate}
\end{methode}

\begin{exoresolu}[Rédaction d'un avis : « Services en ligne à la mairie de Douala »]
\textbf{Énoncé :} Rédigez un avis officiel (\cle{通知}) de 80 à 100 caractères chinois pour la Mairie de Douala (杜拉市政府 Dùlā Shì Zhèngfǔ) annonçant le lancement de trois nouveaux services en ligne : 1. Prise de rendez-vous (预约 yùyuē), 2. Téléchargement d'actes d'état civil (下载证件 xiàzài zhèngjiàn), 3. Paiement des taxes locales (缴纳税费 jiǎonà shuìfèi). Indiquez le site web (www.douala.gov.cm) et le numéro vert (112). Date : 15 octobre 2025.
\tcblower
\textbf{Solution détaillée :}
\begin{enumerate}
    \item \textbf{Choix du format :} \cle{通知} (Tōngzhī) car c'est une information descendante administration $\to$ usagers.
    \item \textbf{Construction du titre :} \zh{关于开通网上办事服务的通知} (Guānyú kāitōng wǎngshàng bànshì fúwù de tōngzhī) — « Avis sur l'ouverture des services de démarches en ligne ». Structure \zh{关于 + Sujet + 的通知}.
    \item \textbf{Corps du texte — Application de la méthode :}
    \begin{itemize}
        \item \textit{Appel :} \zh{尊敬的市民朋友们：}
        \item \textit{Contexte :} \zh{为进一步方便市民办事，提高办事效率，杜拉市政府决定开通网上办事大厅。} (But : \zh{为/为了} + double objectif).
        \item \textit{Énumération (connecteurs \zh{一}、\zh{二}、\zh{三}) :}
        \zh{现提供以下三项服务：} \\
        \zh{一、网上预约 (wǎngshàng yùyuē)；} \\
        \zh{二、证件在线下载 (zhèngjiàn zàixiàn xiàzài)；} \\
        \zh{三、税费网上缴纳 (shuìfèi wǎngshàng jiǎonà)。}
        \item \textit{Modalités :} \zh{请市民登录网站 www.douala.gov.cm 或拨打服务热线 112 咨询办理。} (Impératif poli \zh{请} + alternatives \zh{或}).
    \end{itemize}
    \item \textbf{Conclusion :} \zh{特此通知。}
    \item \textbf{Signature/Date :} \zh{杜拉市政府} \\ \zh{2025年10月15日}
\end{enumerate}
\textbf{Texte final (caractères + pinyin) :}
\begin{textebox}[Avis officiel — Mairie de Douala]
关于开通网上办事服务的通知

尊敬的市民朋友们：
为进一步方便市民办事，提高办事效率，杜拉市政府决定开通网上办事大厅。现提供以下三项服务：
一、网上预约 (wǎngshàng yùyuē)；
二、证件在线下载 (zhèngjiàn zàixiàn xiàzài)；
三、税费网上缴纳 (shuìfèi wǎngshàng jiǎonà)。
请市民登录网站 www.douala.gov.cm 或拨打服务热线 112 咨询办理。
特此通知。

杜拉市政府
2025年10月15日
\end{textebox}
\textbf{Conclusion :} Le texte respecte la structure administrative, utilise le vocabulaire ciblé (\zh{预约, 下载, 缴纳, 热线}) et les connecteurs (\zh{为/为了, 一/二/三, 或, 特此}).
\end{exoresolu}

\begin{methode}[Écrire correctement les caractères du thème : Analyse composant par composant (字形结构分析)]
\textbf{Objectif :} Maîtriser l'écriture standard (规范字) des mots-clés du thème « services administratifs » pour éviter les confusions visuelles et les erreurs de traits.
\textbf{Étapes :}
\begin{enumerate}
    \item \textbf{Identifier le radical (部首 bùshǒu) et la structure (结构 jiégòu) :} Gauche-Droite ([G-D] ), Haut-Bas ([H-B] ), Enfermement ([cadre] /[cadre haut] /[cadre bas] /[cadre gauche] ), Caractère unique.
    \item \textbf{Respecter l'ordre des traits (笔顺 bǐshùn) universel :}
    \begin{itemize}
        \item Horizontal avant Vertical (横竖) : \zh{十} (héng zhù).
        \item Gauche avant Droite (撇捺) : \zh{人} (piě nà).
        \item Haut avant Bas (上下) : \zh{三} (shàng xià).
        \item Extérieur avant Intérieur (外里) : \zh{同} (wài lǐ).
        \item Fermer en dernier (封口最后) : \zh{口} $\to$ \zh{日} $\to$ \zh{田}.
        \item Trait central avant les ailes (中间两边) : \zh{小} (zhōngjiān liǎngbiān).
    \end{itemize}
    \item \textbf{Repérer les « caractères proches » (易混淆字) et les différencier par le radical ou un trait :}
    \begin{center}
    \begin{tabularx}{\linewidth}{|X|X|X|X|}\hline
    \rowcolor{popblueL} Caractère cible & Radical / Clé & Caractère proche (Piège) & Différence clé \\ \hline
    \zh{务} (wù) & \zh{力} (force) & \zh{务} vs \zh{務} (traditionnel) & Simplifié : \zh{力} en bas. Ne pas ajouter de trait au milieu. \\ \hline
    \zh{办} (bàn) & \zh{力} (force) & \zh{办} vs \zh{务} & \zh{办} = \zh{力} + \zh{卜} (divination) à gauche. \zh{务} = \zh{力} + \zh{矛} (lance) simplifié. \\ \hline
    \zh{税} (shuì) & \zh{禾} (riz) & \zh{税} vs \zh{稅} (traditionnel) & Simplifié : \zh{兑} (échange) à droite. Traditionnel : \zh{說} (parler) à droite. \\ \hline
    \zh{证} (zhèng) & \zh{讠} (parole) & \zh{证} vs \zh{證} (traditionnel) & Simplifié : \zh{正} (correct) à droite. Traditionnel : \zh{登} (monter) au-dessus de \zh{言}. \\ \hline
    \zh{窗} (chuāng) & \zh{穴} (caverne/toit) & \zh{窗} vs \zh{囱} (cheminée) & \zh{窗} = \zh{穴} + \zh{囱}. \zh{囱} = \zh{囗} + \zh{工}. Ne pas confondre le toit \zh{穴} et l'enfermement \zh{囗}. \\ \hline
    \zh{咨} (zī) & \zh{讠} (parole) & \zh{咨} vs \zh{资} (zī, ressource) & \zh{咨} (consulter) = \zh{讠} + \zh{咨} (phonétique). \zh{资} (capital) = \zh{贝} (coquillage/monnaie) + \zh{次}. \\ \hline
    \end{tabularx}
    \end{center}
    \item \textbf{Entraînement ciblé :} Écrire 5 fois chaque caractère en dictant l'ordre des traits à voix haute. Utiliser du papier quadrillé (田字格) pour respecter les proportions (centre de gravité).
\end{enumerate}
\end{methode}

\begin{exoresolu}[Dictée commentée et discrimination de caractères]
\textbf{Énoncé :}
\begin{enumerate}
    \item Écrivez en caractères (avec ordre des traits décrit) : \zh{服务} (fúwù), \zh{办理} (bànlǐ), \zh{缴费} (jiǎofèi), \zh{咨询} (zīxún).
    \item Choisissez le bon caractère entre parenthèses :
    \begin{itemize}
        \item Veuillez vous rendre au guichet (窗/囱) 口 3.
        \item Il faut présenter une pièce d'identité (证/识) 件。
        \item Cette démarche (务/办) 理 取消了。
        \item Nous offrons un service de (咨/资) 询 免费。
    \end{itemize}
\end{enumerate}
\tcblower
\textbf{Solution détaillée :}
\textbf{1. Écriture commentée (Ordre des traits) :}
\begin{itemize}
    \item \zh{服} (fú) [8 traits, [G-D]  \zh{月} + \zh{ㄗ}] : \zh{月} (丿, (trait brisé), 一, 一) $\to$ partie droite ((trait brisé), 丨, 丿, 丶). \textit{Rappel :} \zh{月} s'écrit : piě, héng zhé gōu, héng, héng. Ne pas confondre avec \zh{用} ou \zh{肉}.
    \item \zh{务} (wù) [7 traits, [H-B]  \zh{矛} + \zh{力}] : \zh{矛} (一, 丨, 丿, 丶, 一) $\to$ \zh{力} ((trait brisé), 丿). \textbf{\cle{Attention}} : Le 5e trait de \zh{矛} est une horizontale qui traverse, pas un crochet.
    \item \zh{办} (bàn) [8 traits, [G-D]  \zh{卜} + \zh{力}] : \zh{卜} (丶, (trait brisé)) $\to$ \zh{力} ((trait brisé), 丿). \textit{Différence avec 务 :} Le composant gauche n'est pas \zh{矛} mais \zh{卜} (un point, un trait zigzag horizontal-pliant-gauche).
    \item \zh{理} (lǐ) [11 traits, [G-D]  \zh{王} + \zh{里}] : \zh{王} (一, 一, 丨, 一) $\to$ \zh{里} (田: (trait brisé), 丨, 一, 一, 一 + 土: 一, 丨). \textit{Note :} \zh{里} = \zh{田} + \zh{土}.
    \item \zh{缴} (jiǎo) [9 traits, [G-D]  \zh{纟} + \zh{交}] : \zh{纟} (丿, (trait brisé), ㇀) $\to$ \zh{交} (丨, (trait brisé), 丿, ㇀, 丿, 丶). \textit{Simplifié :} \zh{纟} (3 traits) au lieu de \zh{糸} (6 traits).
    \item \zh{费} (fèi) [9 traits, [G-D]  \zh{弗} + \zh{贝}] : \zh{弗} (一, 丨, (trait brisé), 乙, 乀) $\to$ \zh{贝} ((trait brisé), 丨, (trait brisé), 丶, 丶). \textit{Radical \zh{贝} :} Argent/Échange.
    \item \zh{咨} (zī) [9 traits, [G-D]  \zh{讠} + \zh{咨}] : \zh{讠} (丶, (trait brisé)) $\to$ \zh{咨} (phonétique, 7 traits). \zh{资} (zī) [7 traits, [G-D]  \zh{贝} + \zh{次}] : \zh{贝} $\to$ \zh{次} (欠 + 一). \textbf{Clé :} \zh{讠} = langage (consulter) vs \zh{贝} = bien/argent (ressource).
\end{itemize}
\textbf{2. Discrimination (QCM) :}
\begin{enumerate}
    \item \zh{窗口} (chuāngkǒu) — \textbf{\zh{窗}} (radical \zh{穴} toit/caverne = ouverture dans un mur). \zh{囱} (cōng) = cheminée (radical \zh{囗} enfermement).
    \item \zh{证件} (zhèngjiàn) — \textbf{\zh{证}} (radical \zh{讠} parole/preuve = document officiel). \zh{识} (shí) = connaître/reconnaître.
    \item \zh{办理} (bànlǐ) — \textbf{\zh{办}} (radical \zh{力} force/action = traiter une affaire). \zh{务} (wù) = affaire/tâche (nom), ne s'utilise pas seul comme verbe « traiter ».
    \item \zh{咨询} (zīxún) — \textbf{\zh{咨}} (radical \zh{讠} = demander avis). \zh{资} (zī) = fonds/capital (radical \zh{贝}).
\end{enumerate}
\textbf{Conclusion :} La connaissance des radicaux (\cle{部首}) est le levier principal pour ne pas confondre les homophones ou quasi-homophones très fréquents dans l'administration (\zh{咨/资, 证/识, 务/办, 窗/囱}).
\end{exoresolu}

\begin{methode}[Traduire du français vers le chinois (Thème) : Stratégie de transposition syntaxique]
\textbf{Objectif :} Passer de la structure française (SVO mais avec prépositions complexes, relatives, passif) à la structure chinoise (SVO stricte, thème-rhème, pas de conjugaison, compléments antéposés).
\textbf{Étapes :}
\begin{enumerate}
    \item \textbf{Segmenter la phrase française en unités de sens (groupes nominaux, verbaux, prépositionnels).}
    \item \textbf{Identifier le « Sujet chinois » (主语) :} Souvent le même, mais parfois le complément d'agent (passif) ou le thème (sujet zéro). \textit{Ex:} « Le service est offert par la mairie » $\to$ Sujet chinois = \zh{市政府} (Mairie) ou \zh{这项服务} (Ce service) + \zh{由/被}.
    \item \textbf{Placer les compléments selon l'ordre canonique chinois :} \cle{Sujet + (Adverbial: Temps) + (Adverbial: Lieu/Manière/Moyen) + Verbe + (Complément de résultat/degré) + Objet}.
    \begin{center}
    \begin{tabularx}{\linewidth}{|X|X|}\hline
    \rowcolor{popblueL} Français (ordre variable) & Chinois (ordre fixe) \\ \hline
    « Je vais \textbf{demain} \textbf{à la mairie} \textbf{faire} une carte. » & \zh{我 明天 去市政府 办 卡。} (Temps $\to$ Lieu $\to$ Verbe) \\ \hline
    « Il paie \textbf{par internet} \textbf{les taxes}. » & \zh{他 在网上 缴纳 税费。} (Lieu/Moyen \zh{在网上} avant Verbe) \\ \hline
    « Ce formulaire \textbf{rempli par vous} est valide. » & \zh{您 填写 的 这份 表格 有效。} (Relatif \zh{您填写的} antéposé au Nom) \\ \hline
    \end{tabularx}
    \end{center}
    \item \textbf{Traiter les mots-outils critiques :}
    \begin{itemize}
        \item \cle{的} (de) : Après adjectif/relatif modifiant un nom. \zh{免费的 服务}.
        \item \cle{地} (de) : Après adjectif modifiant un verbe. \zh{免费 地 提供} (souvent omis si monosyllabique : \zh{免费提供}).
        \item \cle{得} (de) : Après verbe, avant complément de degré. \zh{办 得 很快}.
        \item \cle{了} (le) : Aspect accompli (verbe + \zh{了} + quantifié) ou changement d'état (phrase + \zh{了}). \zh{我 已经 办 了 证件 了}.
        \item \cle{过} (guò) : Expérience passée. \zh{我 去 过 市政府}.
    \end{itemize}
    \item \textbf{Convertir les connecteurs français en connecteurs chinois (voir Méthode 1).}
    \item \textbf{Vérification finale :} Accord tons (3e ton + 3e ton $\to$ 2e ton + 3e ton), classificateurs (\cle{一份} 表格, \cle{一项} 服务, \cle{一个} 窗口), pas d'espace entre caractères.
\end{enumerate}
\end{methode}

\begin{exoresolu}[Thème : Phrases administratives types]
\textbf{Énoncé :} Traduisez en chinois (caractères + pinyin).
\begin{enumerate}
    \item Pour faciliter les démarches des citoyens, la mairie a ouvert un guichet unique.
    \item Les usagers peuvent désormais payer les taxes en ligne sans se déplacer.
    \item Le formulaire que vous avez rempli hier est incomplet ; veuillez le refaire.
    \item Avez-vous déjà consulté le site web de la mairie ?
\end{enumerate}
\tcblower
\textbf{Solution détaillée :}
\textbf{Phrase 1 :} « Pour faciliter les démarches des citoyens, la mairie a ouvert un guichet unique. »
\begin{itemize}
    \item \textit{Analyse :} But (Pour faciliter...) + Sujet (la mairie) + Verbe (a ouvert) + Objet (un guichet unique).
    \item \textit{Transposition :} But antéposé (\zh{为了/为} + V + \zh{,}) + Sujet + Temps (implicite passé/accompli) + Verbe + Objet.
    \item \textit{Vocabulaire :} Faciliter \zh{方便} (fāngbiàn), démarches \zh{办事} (bànshì) / \zh{办理手续} (bànlǐ shǒuxù), citoyens \zh{市民} (shìmín) / \zh{群众} (qúnzhòng), mairie \zh{市政府} (shì zhèngfǔ) / \zh{区政府} (qū zhèngfǔ), ouvrir \zh{设立} (shèlì) / \zh{开设} (kāishè), guichet unique \zh{综合窗口} (zōnghé chuāngkǒu) / \zh{一站式服务窗口} (yī zhàn shì fúwù chuāngkǒu).
    \item \textit{Proposition :} \zh{为方便市民办事，市政府设立了综合窗口。} (Wèi fāngbiàn shìmín bànshì, shì zhèngfǔ shèlì le zōnghé chuāngkǒu.)
    \item \textit{Justification :} \zh{为} (but) en tête. \zh{了} après \zh{设立} car action accomplie dans le passé avec résultat actuel. \zh{综合窗口} terme technique standard.
\end{itemize}

\textbf{Phrase 2 :} « Les usagers peuvent désormais payer les taxes en ligne sans se déplacer. »
\begin{itemize}
    \item \textit{Analyse :} Sujet (Usagers) + Modal (peuvent) + Temps (désormais) + Verbe (payer) + Objet (taxes) + Lieu/Moyen (en ligne) + Condition/Négation (sans se déplacer).
    \item \textit{Ordre chinois :} Sujet + Temps (\zh{现在/如今}) + Modal (\zh{可以/能}) + Lieu/Moyen (\zh{在网上/通过网络}) + Verbe + Objet + Complément négatif (\zh{不用/无需} + Verbe).
    \item \textit{Vocabulaire :} Usagers \zh{办事群众} (bànshì qúnzhòng) / \zh{市民} (shìmín), désormais \zh{现在} (xiànzài) / \zh{如今} (rújīn), payer \zh{缴纳} (jiǎonà) / \zh{交} (jiāo), taxes \zh{税费} (shuìfèi), en ligne \zh{在线} (zàixiàn) / \zh{网上} (wǎngshàng), se déplacer \zh{跑腿} (pǎotuǐ - familier adm) / \zh{亲自前往} (qīnzì qiánwǎng).
    \item \textit{Proposition :} \zh{市民现在可以在网上缴纳税费，不用亲自跑腿了。} (Shìmín xiànzài kěyǐ zài wǎngshàng jiǎonà shuìfèi, bùyòng qīnzì pǎotuǐ le.)
    \item \textit{Justification :} \zh{在网上} (locatif) placé AVANT le verbe \zh{缴纳}. \zh{不用} (pas besoin) + Verbe pour « sans + gérondif ». \zh{了} final : changement d'état (nouvelle situation).
\end{itemize}

\textbf{Phrase 3 :} « Le formulaire que vous avez rempli hier est incomplet ; veuillez le refaire. »
\begin{itemize}
    \item \textit{Analyse :} Sujet complexe (Formulaire [que vous avez rempli hier]) + Verbe être + Adj (incomplet) ; Impératif (Veuillez) + Objet (le) + Verbe (refaire).
    \item \textit{Transposition :} Relatif chinois antéposé : \zh{您 昨天 填写 的} (Vous hier rempli \cle{的}) + Nom \zh{表格} (Formulaire) + \zh{不完整} (incomplet) + \zh{请 重新 填写} (Veuillez re-remplir).
    \item \textit{Vocabulaire :} Formulaire \zh{表格} (biǎogé) / \zh{申请表} (shēnqǐngbiǎo), remplir \zh{填写} (tiánxiě), hier \zh{昨天} (zuótiān), incomplet \zh{不完整} (bù wánzhěng) / \zh{填得不全} (tián de bù quán), veuillez \zh{请} (qǐng) / \zh{麻烦} (máfan), refaire \zh{重新填写} (chóngxīn tiánxiě) / \zh{补齐} (bǔqí - compléter).
    \item \textit{Proposition :} \zh{您昨天填写的表格不完整，请重新填写。} (Nín zuótiān tiánxiě de biǎogé bù wánzhěng, qǐng chóngxīn tiánxiě.)
    \item \textit{Justification :} Structure \zh{的} pour la relative. \zh{不完整} (adj predicatif). Phrase impérative sans sujet (sujet zéro = vous).
\end{itemize}

\textbf{Phrase 4 :} « Avez-vous déjà consulté le site web de la mairie ? »
\begin{itemize}
    \item \textit{Analyse :} Question expérience passée (déjà) + Verbe (consulter) + Objet (site web).
    \item \textit{Grammaire :} \cle{过} (guò) pour l'expérience. Structure : Sujet + \zh{去过/看过/访问过} + Objet + \zh{吗} ? Ou \zh{有没有} + V + \zh{过} + Objet ?
    \item \textit{Vocabulaire :} Déjà \zh{已经} (yǐjīng) / \zh{曾经} (céngjīng), consulter/visiter \zh{访问} (fǎngwèn) / \zh{浏览} (liúlǎn) / \zh{查看} (chákàn), site web \zh{网站} (wǎngzhàn) / \zh{官网} (guānwǎng - site officiel).
    \item \textit{Proposition :} \zh{您访问过市政府官网吗？} (Nín fǎngwèn guò shì zhèngfǔ guānwǎng ma?) ou \zh{您有没有去过市政府的网站？} (Nín yǒu méi yǒu qù guò shì zhèngfǔ de wǎngzhàn?)
    \item \textit{Justification :} \zh{过} placé immédiatement après le verbe \zh{访问}. \zh{官网} terme précis pour « site web de l'administration ». \zh{吗} particule interrogative finale.
\end{itemize}
\end{exoresolu}

\begin{methode}[Traduire du chinois vers le français (Version) : Déverbalisation et reformulation]
\textbf{Objectif :} Produire un français naturel, administratif ou courant, en évitant le calque littéral (mot à mot) qui donne un français « chinois ».
\textbf{Étapes :}
\begin{enumerate}
    \item \textbf{Segmenter le chinois en blocs syntaxiques :} Identifier le Sujet (souvent zéro), le Prédicat (Verbe + compléments), les modificateurs (\zh{的}, \zh{地}, \zh{得}).
    \item \textbf{Restaurer la grammaire française implicite :}
    \begin{itemize}
        \item \textbf{Temps/Aspect :} \zh{了} (accompli/passé récent), \zh{过} (expérience/vécu), \zh{在/正在} (en train de), \zh{会/要/将} (futur/projet). \textit{Ex:} \zh{已经办好了} $\to$ « a déjà fini de traiter » / « est déjà traité ».
        \item \textbf{Passif / Voice moyenne :} \zh{被} (bèi), \zh{由} (yóu), \zh{让} (ràng), \zh{给} (gěi) $\to$ Passif français (« est traité par »). Sujet zéro + Verbe transitif $\to$ Passif réfléchi (« cela se fait ») ou Impersonnel (« il faut »).
        \item \textbf{Relatives (\zh{的}) :} \zh{您办理的业务} $\to$ « L'affaire \textbf{que vous traitez} » / « Votre démarche ».
        \item \textbf{Compléments de lieu/manière antéposés :} \zh{在网上办理} $\to$ « Traiter \textbf{en ligne} » (pas « Sur internet traiter »).
    \end{itemize}
    \item \textbf{Choisir le registre lexicologique français approprié :}
    \begin{center}
    \begin{tabularx}{\linewidth}{|X|X|X|}\hline
    \rowcolor{popblueL} Chinois (Registre admin) & Calque littéral (À éviter) & Français standard (Cible) \\ \hline
    \zh{办理手续} (bànlǐ shǒuxù) & « Traiter des procédures » & \textbf{Effectuer des démarches / Faire les formalités} \\ \hline
    \zh{跑腿} (pǎotuǐ) & « Courir les jambes » & \textbf{Se déplacer / Faire la queue / Faire le déplacement} \\ \hline
    \zh{一站式服务} (yī zhàn shì fúwù) & « Service de type une gare » & \textbf{Guichet unique / Service tout-en-un} \\ \hline
    \zh{优化营商环境} (yōuhuà yíngshāng huánjìng) & « Optimiser l'environnement commercial » & \textbf{Améliorer le climat des affaires} \\ \hline
    \zh{便民利民} (biànmín lìmín) & « Commodité peuple profit peuple » & \textbf{Simplifier la vie des usagers / Services de proximité} \\ \hline
    \zh{特此通知} (tè cǐ tōngzhī) & « Spécialement par ceci notifier » & \textbf{Le présent avis / Par la présente} \\ \hline
    \end{tabularx}
    \end{center}
    \item \textbf{Soigner la ponctuation et la typographie françaises :} Majuscules, accents, espaces insécables avant : ; ? !, guillemets « ».
\end{enumerate}
\end{methode}

\begin{exoresolu}[Version : Extrait d'avis officiel de la mairie de Yaoundé]
\textbf{Énoncé :} Traduisez en français naturel le texte suivant :
\begin{textebox}[Source : Communiqué mairie de Yaoundé (约恩德市政府通告)]
约恩德市政府为优化政务服务，提高办事效率，自2025年11月1日起，在市民服务中心设立“综合窗口”，提供“最多跑一次”服务。市民办理身份证、营业执照、不动产登记等高频事项，只需在一个窗口提交材料，无需跑多个科室。如有疑问，请拨打12345热线咨询。特此通告。
\end{textebox}
\tcblower
\textbf{Solution détaillée :}
\textbf{1. Segmentation et analyse structurelle :}
\begin{itemize}
    \item \zh{约恩德市政府} (Sujet : Mairie de Yaoundé)
    \item \zh{为 优化 政务 服务，提高 办事 效率} (But : \zh{为} + VO, VO)
    \item \zh{自 2025年11月1日 起} (Temps : À partir du 1er nov 2025)
    \item \zh{在 市民 服务 中心} (Lieu : Au centre de service aux citoyens)
    \item \zh{设立 “综合 窗口”} (Verbe + Objet : Établir « guichet unique »)
    \item \zh{提供 “最多 跑 一次” 服务} (Verbe + Objet : Offrir service « au plus une fois »)
    \item \zh{市民 办理 身份证、营业 执照、不动产 登记 等 高频 事项} (Nouveau Sujet/Thème : Citoyens traitant CI, licence, foncier... fréquents)
    \item \zh{只需 在 一个 窗口 提交 材料} (Verbe modal + Lieu + Verbe + Obj : N'ont qu'à déposer dossiers à un guichet)
    \item \zh{无需 跑 多个 科室} (Négation + Verbe + Obj : Pas besoin de courir plusieurs services)
    \item \zh{如有 疑问，请 拨打 12345 热线 咨询} (Condition + Impératif : Si questions, appeler hotline)
    \item \zh{特此 通告} (Formule close : Par la présente avis).
\end{itemize}

\textbf{2. Reformulation progressive (Déverbalisation) :}
\begin{itemize}
    \item \textit{Brouillon littéral :} « Yaoundé Mairie pour optimiser affaires-gouvernement service, augmenter traiter-affaires efficacité, depuis 2025-11-1 commencer, à citoyen service centre établir "complet guichet", fournir "plusieurs fois courir une fois" service. Citoyens traiter carte-identité, commerce licence, immobilier enregistrement etc haut-fréquence affaires, seulement besoin à un guichet soumettre matériaux, pas besoin courir plusieurs départements. Si avoir doute, s'il-vous-plaît composer 12345 hotline consulter. Spécialement ceci communiquer. »
    \item \textit{Correction syntaxique :} « La Mairie de Yaoundé, \textbf{afin d'optimiser les services administratifs et d'améliorer l'efficacité du traitement des dossiers}, \textbf{à compter du 1er novembre 2025}, \textbf{met en place} au Centre de services aux citoyens un \textbf{guichet unique} offrant un service \textbf{"au plus un déplacement"}.
    \item \textit{Traitement de la 2e phrase :} « Pour les démarches à forte fréquentation telles que \textbf{la carte d'identité, la licence commerciale, l'enregistrement foncier}, les usagers \textbf{n'ont plus qu'à déposer leur dossier à un seul guichet}, \textbf{sans avoir à se rendre dans plusieurs services}. »
    \item \textit{Fin :} « Pour tout renseignement, \textbf{veuillez appeler le 12345}. \textbf{Le présent avis est publié pour information.} »
\end{itemize}

\textbf{3. Version finale (Français administratif standard) :}
\begin{quote}
\textbf{COMMUNIQUÉ DE LA MAIRIE DE YAOUNDÉ}

\textit{Dans le but d'optimiser les services administratifs et d'améliorer l'efficacité du traitement des dossiers, la Mairie de Yaoundé met en place, à compter du 1er novembre 2025, un guichet unique au Centre de services aux citoyens, offrant la garantie du « déplacement unique » (service « au plus une fois »).}

Pour les démarches à forte fréquentation — carte d'identité nationale, licence commerciale, immatriculation foncière —, les usagers n'ont désormais plus qu'à déposer leur dossier à un seul guichet, sans avoir à se rendre dans plusieurs services.

Pour tout renseignement, merci de contacter le 12345 (numéro vert).

\textbf{Le présent communiqué est publié pour information.}
\end{quote}
\textbf{Points de vigilance (Pièges évités) :}
\begin{itemize}
    \item \zh{政务服务} $\neq$ « services politiques » mais \textbf{services administratifs / publics}.
    \item \zh{最多跑一次} $\neq$ « courir au plus une fois » mais terme technique \textbf{« au plus un déplacement / guichet unique »}.
    \item \zh{高频事项} $\neq$ « affaires haute fréquence » mais \textbf{« démarches les plus courantes / à forte demande »}.
    \item \zh{科室} $\neq$ « salles de classe » mais \textbf{« services / bureaux / départements »}.
    \item \zh{特此通告} $\to$ Formule figée \textbf{« Le présent communiqué / avis »}.
\end{itemize}
\end{exoresolu}

\begin{methode}[Production orale guidée : Jeu de rôle « Guichet unique / Service en ligne » (情景对话)]
\textbf{Objectif :} Simuler une interaction usager / agent d'accueil pour s'informer, prendre rendez-vous ou déposer un dossier. Compétence : Interaction orale (HSK 3-4).
\textbf{Étapes :}
\begin{enumerate}
    \item \textbf{Assigner les rôles :} A = \zh{市民} (Usager) / B = \zh{大厅工作人员 / 志愿者} (Agent d'accueil / Bénévole).
    \item \textbf{Préparer les "cartes scénarios" (carte A / carte B) avec objectifs distincts :}
    \begin{itemize}
        \item \textit{Scénario 1 (Information) :} A veut savoir comment faire une carte d'identité (\zh{身份证}) pour son enfant mineur. B guide : pièces (\zh{户口本, 出生证, 照片}), prise RDV (\zh{预约}), délai (\zh{7个工作日}).
        \item \textit{Scénario 2 (Problème technique) :} A a essayé de payer une taxe (\zh{缴纳税费}) sur le site \zh{官网} mais paiement échoué (\zh{支付失败}). B vérifie : numéro commande (\zh{订单号}), propose solution (\zh{重新支付 / 线下缴纳}).
        \item \textit{Scénario 3 (Réclamation) :} A s'est déplacé (\zh{跑腿}) pour rien car guichet fermé (\zh{窗口关闭}) ou dossier incomplet non signalé. B s'excuse (\zh{抱歉/对不起}), vérifie dossier (\zh{查一下}), reprogramme (\zh{重新安排}).
    \end{itemize}
    \item \textbf{Structurer l'échange (3 temps) :}
    \begin{enumerate}
        \item \textbf{Accueil / Identification du besoin :} \zh{您好，请问有什么可以帮您？} / \zh{我想咨询一下……} / \zh{我来办理……}.
        \item \textbf{Échange d'informations (Q/R) :} Utiliser \zh{请问} (politesse), \zh{需要带什么材料？} (Quels docs ?), \zh{在哪里/几点办？} (Où/Quand ?), \zh{可以网上办吗？} (En ligne possible ?).
        \item \textbf{Clôture / Confirmation :} \zh{好的，我明白了。谢谢。} / \zh{不客气，祝您办事顺利。} (Formule admin standard).
    \end{enumerate}
    \item \textbf{Contraintes linguistiques :} Utiliser au moins 3 connecteurs (\zh{首先/然后/最后}, \zh{如果...就...}, \zh{虽然...但是...}) et 5 mots du vocabulaire thème (\zh{预约, 材料, 窗口, 热线, 办理, 证件, 网上, 缴纳, 取号, 叫号}).
    \item \textbf{Auto-évaluation / Pair-évaluation :} Grille : Prononciation (tons), Fluidité, Vocabulaire précis, Politesse (\zh{您/请/谢谢}), Résolution du problème.
\end{enumerate}
\end{methode}

\begin{exoresolu}[Simulation orale : Scénario « Paiement taxe en ligne échoué »]
\textbf{Énoncé :} Complétez le dialogue ci-dessous (rôle B : Agent). Puis jouez-le à deux.
\begin{textebox}[Dialogue à trous — Guichet Mairie Douala]
\textbf{A (Usager) :} \zh{您好，我想咨询一下网上缴税的问题。}
\textbf{B (Agent) :} \zh{\_\_\_\_\_\_\_\_\_\_\_\_（欢迎/询问细节）}
\textbf{A :} \zh{我刚才在官网缴纳房产税，支付显示失败，但银行卡扣钱了。}
\textbf{B :} \zh{\_\_\_\_\_\_\_\_\_\_\_\_（安抚/索要单号）}
\textbf{A :} \zh{订单号是 202511010088。}
\textbf{B :} \zh{\_\_\_\_\_\_\_\_\_\_\_\_（查系统/给方案：退款或重缴）}
\textbf{A :} \zh{那我现在可以直接在这儿重新缴纳吗？}
\textbf{B :} \zh{\_\_\_\_\_\_\_\_\_\_\_\_（确认窗口功能/指引）}
\textbf{A :} \zh{好的，谢谢您。}
\textbf{B :} \zh{\_\_\_\_\_\_\_\_\_\_\_\_（结束语）}
\end{textebox}
\tcblower
\textbf{Proposition de complétion (Modèle corrigé) :}
\begin{textebox}[Dialogue complet — Corrigé type]
\textbf{A :} \zh{您好，我想咨询一下网上缴税的问题。} (Nín hǎo, wǒ xiǎng zīxún yíxià wǎngshàng jiǎoshuì de wèntí.)
\textbf{B :} \zh{您好！请问您遇到什么具体问题了？请说一下情况。} (Nín hǎo! Qǐngwèn nín yùdào shénme jùtǐ wèntí le? Qǐng shuō yíxià qíngkuàng.) \textit{[Accueil + Question ouverte]}
\textbf{A :} \zh{我刚才在官网缴纳房产税，支付显示失败，但银行卡扣钱了。} (Wǒ gāngcái zài guānwǎng jiǎonà fángchǎnshuì, zhīfù xiǎnshì shībài, dàn yínhángkǎ kòu qián le.)
\textbf{B :} \zh{别着急，这种情况偶尔会发生。请提供一下订单号，我帮您查询一下系统记录。} (Bié zhāojí, zhè zhǒng qíngkuàng ǒu'ěr huì fāshēng. Qǐng tígōng yíxià dìngdānhào, wǒ bāng nín cháxún yíxià xìtǒng jìlù.) \textit{[Empathie + Action concrète]}
\textbf{A :} \zh{订单号是 202511010088。} (Dìngdānhào shì 202511010088.)
\textbf{B :} \zh{稍等，我查一下……系统显示订单状态是“支付失败”，款项会在3-5个工作日原路退回。您可以现在直接在咱们这儿的“综合窗口”刷卡重新缴纳，也可以回家等退款后再网上操作。} (Shāo děng, wǒ chá yíxià… xìtǒng xiǎnshì dìngdān zhuàngtài shì “zhīfù shībài”, kuǎnxiàng huì zài 3-5 gè gōngzuòrì yuánlù tuìhuí. Nín kěyǐ xiànzài zhíjiē zài zánmen zhèr de “zōnghé chuāngkǒu” shuākǎ chóngxīn jiǎonà, yě kěyǐ huíjiā děng tuìkuǎn hòu zài wǎngshàng cāozuò.) \textit{[Diagnostic + Deux solutions alternatives]}
\textbf{A :} \zh{那我现在可以直接在这儿重新缴纳吗？} (Nà wǒ xiànzài kěyǐ zhíjiē zài zhèr chóngxīn jiǎonà ma?)
\textbf{B :} \zh{可以的。请到3号窗口，刷身份证取号，业务员会协助您完成缴纳。} (Kěyǐ de. Qǐng dào 3 hào chuāngkǒu, shuā shēnfènzhèng qǔ hào, yèwùyuán huì xiézhù nín wánchéng jiǎonà.) \textit{[Instruction précise : Lieu + Action + Soutien]}
\textbf{A :} \zh{好的，谢谢您。} (Hǎo de, xièxie nín.)
\textbf{B :} \zh{不客气，祝您办事顺利，再见。} (Bú kèqi, zhù nín bànshì shùnlì, zàijiàn.) \textit{[Formule de clôture admin standard]}
\end{textebox}
\textbf{Analyse didactique :}
\begin{itemize}
    \item \textbf{Vocabulaire ciblé utilisé :} \zh{咨询, 缴纳, 房产税, 支付, 失败, 扣钱, 订单号, 查询, 系统, 退款, 综合窗口, 刷卡, 取号, 业务员, 协助}.
    \item \textbf{Structures :} \zh{别着急} (Impératif négatif rassurant), \zh{会在...原路退回} (Passif \zh{被} implicite / Verbe \zh{会} futur), \zh{可以...也可以...} (Alternative), \zh{请到...刷...取号...} (Série d'actions V1 V2 V3).
    \item \textbf{Registre :} \zh{您} (politesse), \zh{咱们} (inclusif « notre guichet »), \zh{祝您办事顺利} (formule consacrée).
\end{itemize}
\end{exoresolu}

\begin{methode}[Réviser et consolider : Fiche de synthèse « Services aux usagers » (知识清单)]
\textbf{Objectif :} Disposer d'une fiche de révision unique (mindmap/liste) pour le Bac : vocabulaire, structures, caractères, culture.
\textbf{Contenu de la fiche (à photocopier/compléter) :}
\begin{enumerate}
    \item \textbf{Vocabulaire cœur (20 mots) — Catégorisé :}
    \begin{center}
    \begin{tabularx}{\linewidth}{|X|X|}\hline
    \rowcolor{popblueL} \textbf{Acteurs / Lieux} & \textbf{Actions / Documents} \\ \hline
    \zh{市民} (shìmín) usager / \zh{工作人员} (gōngzuò rényuán) agent & \zh{办理} (bànlǐ) traiter / \zh{申请} (shēnqǐng) demander \\
    \zh{市政府/区政府} (shì/qū zhèngfǔ) mairie & \zh{提交/交} (tíjiāo/jiāo) déposer \\
    \zh{服务大厅/中心} (fúwù dàtīng/zhōngxīn) hall/centre & \zh{审核} (shěnhé) examiner / \zh{批准} (pīzhǔn) approuver \\
    \zh{综合窗口} (zōnghé chuāngkǒu) guichet unique & \zh{领取/取} (lǐngqǔ/qǔ) retirer \\
    \zh{科室/部门} (kēshì/bùmén) service/bureau & \zh{预约} (yùyuē) RDV / \zh{取号} (qǔhào) prendre ticket \\
    \zh{热线/电话} (rèxiàn/diànhuà) hotline & \zh{缴费/缴纳} (jiǎofèi/jiǎonà) payer \\
    \zh{官网/网站} (guānwǎng/wǎngzhàn) site officiel & \zh{材料/证件} (cáiliào/zhèngjiàn) pièces justificatives \\
    \hline
    \end{tabularx}
    \end{center}
    \item \textbf{Structures grammaticales « Bac » (3 indispensables) :}
    \begin{enumerate}
        \item \textbf{把 字句 (bǎ zì jù) — Action sur objet défini / Résultat :} \zh{请把材料交给窗口。} (Qǐng bǎ cáiliào jiāo gěi chuāngkǒu.) $\to$ « Veuillez remettre les documents au guichet. » \textit{Condition :} Objet connu/spécifique + Verbe + Complément de résultat/direction.
        \item \textbf{被 字句 (bèi zì jù) — Passif (souvent contrainte) :} \zh{我的申请被退回了。} (Wǒ de shēnqǐng bèi tuíhuí le.) $\to$ « Ma demande a été rejetée/retournée. » \textit{Marque l'adversité.}
        \item \textbf{连词 复句 — Logique condition/hypothèse :} \zh{如果材料齐全，当场就能办结。} (Rúguǒ cáiliào qíquán, dāngchǎng jiù néng bànjié.) $\to$ « Si le dossier est complet, on peut boucler sur place. »
    \end{enumerate}
    \item \textbf{Caractères « Pièges » à écrire parfaitement (Dictée Bac) :}
    \zh{务 办 税 证 窗 咨 料 领 执 登} (wù, bàn, shuì, zhèng, chuāng, zī, liào, lǐng, zhí, dēng).
    \item \textbf{Repères socioculturels Cameroun / Chine :}
    \begin{itemize}
        \item \textbf{Cameroun :} Décentralisation (Communes, Régions). Mairies = \zh{市政府} (ville) / \zh{区政府} (arrondissement). Services : État civil, Urbanisme, Marchés, Impôts locaux (\zh{房产税} rare pour particuliers, \zh{营业税} commerçants). Numérique : \zh{“E-Gouv”} (portail), paiement Mobile Money (\zh{手机支付}) en essor.
        \item \textbf{Chine :} \zh{“放管服”改革} (Fàng guǎn fú gǎigé) — « Délégation, Régulation, Service » (Réforme admin 2013+). Objectif : \zh{“最多跑一次”} (Zuìduō pǎo yīcì) — « Au plus un déplacement ». \zh{“一网通办”} (Yī wǎng tōng bàn) — « Tout sur un seul réseau ». \zh{政务服务中心} (Zhèngwù fúwù zhōngxīn) — Centres uniques physiques + numériques. \zh{电子营业执照} (Diànzǐ yíngyè zhízhào) — Licence dématérialisée.
    \end{itemize}
\end{enumerate}
\end{methode}

\begin{exoresolu}[Bac Blanc : Sujet complet Type (Compréhension + Expression + Traduction)]
\textbf{Énoncé :} Durée : 1h. Documents autorisés : aucun.
\textbf{Partie A : Compréhension écrite (10 pts)} — Lisez le texte ci-dessous (Avis mairie de Buea) et répondez en français.
\begin{textebox}[Avis de la Mairie de Buea (布埃亚市政府通知) — 2025]
关于推行“掌上办事”服务的通知

尊敬的市民朋友们：
为深化“放管服”改革，让数据多跑路，让群众少跑腿，布埃亚市政府决定自2025年12月1日起，正式上线“布埃亚掌上办”微信小程序。市民可通过手机实名注册后，在线办理户籍证明、经营许可证补办、垃圾费缴纳等20项高频事项。
办理流程：打开微信 $\to$ 搜索“布埃亚掌上办” $\to$ 实名认证 $\to$ 选择事项 $\to$ 上传材料照片 $\to$ 提交申请 $\to$ 等待审核 $\to$ 电子证照发放至手机。
如需线下领取纸质证件，可预约最近的社区服务站自助取证。咨询热线：02-33-XX-XX-XX。
特此通知。

布埃亚市政府
2025年11月20日
\end{textebox}
\textbf{Questions :}
\begin{enumerate}
    \item Quel est l'objectif principal de ce nouveau service ? (2 pts)
    \item Citez 3 exemples de démarches concernées. (3 pts)
    \item Expliquez la procédure en 4 étapes clés (verbes chinois + français). (3 pts)
    \item Comment récupérer un document papier si on le souhaite ? (2 pts)
\end{enumerate}

\textbf{Partie B : Expression écrite (10 pts)} — Vous êtes le Secrétaire Général de la Mairie de Dschang (江城). Rédigez un \cle{通知} (80-100 caractères) pour annoncer l'extension des horaires d'ouverture du guichet unique : désormais ouvert le samedi matin (8h-12h) pour les actifs. Indiquez la date d'effet (1er mars 2026) et le numéro d'info.

\textbf{Partie C : Traduction (10 pts)}
\begin{enumerate}
    \item \textbf{Thème (Fr $\to$ Zh) :} « Afin de faciliter la vie des habitants, la mairie a mis en place un service de prise de rendez-vous en ligne. » (5 pts)
    \item \textbf{Version (Zh $\to$ Fr) :} \zh{市民只需在手机上点一点，就能把证件办好，真正实现了“数据多跑路，群众少跑腿”。} (5 pts)
\end{enumerate}
\tcblower
\textbf{Corrigé détaillé et barème :}

\textbf{Partie A : Compréhension}
\begin{enumerate}
    \item \textbf{Objectif :} Approfondir la réforme « 放管服 » (Délégation/Régulation/Service) $\to$ \textbf{Faire circuler les données plutôt que les gens} (\zh{让数据多跑路，让群众少跑腿}). But : Simplifier, réduire déplacements. (2 pts : 1 pt « réforme », 1 pt « moins se déplacer »).
    \item \textbf{Démarches (3 sur 5) :} \zh{户籍证明} (Attestation de résidence/domicile), \zh{经营许可证补办} (Renouvellement licence commerciale), \zh{垃圾费缴纳} (Paiement taxe ordures), \zh{...等20项} (etc. 20 items). (1 pt/item).
    \item \textbf{Procédure (Verbes chinois) :}
        1. \zh{搜索} (Sǒusuǒ) — Rechercher / Chercher.
        2. \zh{实名认证} (Shímíng rènzhèng) — Authentification réelle / Vérifier identité.
        3. \zh{上传} (Shàngchuán) — Télécharger / Uploader (photos).
        4. \zh{提交} (Tíjiāo) — Soumettre / Déposer (la demande).
        5. \zh{审核} (Shěnhé) — Examiner / Instruire (par l'admin).
        6. \zh{发放} (Fāfàng) — Délivrer / Envoyer (certificat électronique).
        (0.5 pt/verbe correct + traduction, max 3 pts).
    \item \textbf{Récupération papier :} \zh{预约最近的社区服务站自助取证} (Yùyuē zuìjìn de shèqū fúwùzhàn zìzhù qǔzhèng) $\to$ \textbf{Prendre RDV à la station de service communautaire la plus proche pour retrait en libre-service (borne automatique).} (1 pt « prévoir/réserver », 1 pt « borne/autoservice »).
\end{enumerate}

\textbf{Partie B : Expression écrite (Modèle de réponse)}
\begin{textebox}[Proposition de rédaction — Mairie de Dschang]
关于延长综合窗口周六服务时间的通知

尊敬的市民朋友们：
为方便上班族办事，江城市政府决定自2026年3月1日起，综合窗口增设周六上午服务（上午8:00-12:00），正常办理身份证、营业执照等高频业务。请市民根据需要预约办理。咨询热线：02-33-YY-YY-YY。
特此通知。

江城市政府
2026年2月20日
\end{textebox}
\textbf{Barème (10 pts) :} Structure titre/appel/corps/clôture/signature/date (3 pts) ; Vocabulaire précis (\zh{上班族, 增设, 高频业务, 预约}) (3 pts) ; Grammaire/Ordre mots/Connecteurs (\zh{为...决定...起, 正常办理...请...}) (2 pts) ; Caractères/Pinyin/Ponctuation corrects (2 pts).

\textbf{Partie C : Traduction}
\begin{enumerate}
    \item \textbf{Thème :} \zh{为方便市民生活，市政府推出了网上预约服务。} (Wèi fāngbiàn shìmín shēnghuó, shì zhèngfǔ tuīchū le wǎngshàng yùyuē fúwù.)
    \textit{Points :} \zh{为} but (1), \zh{市政府} sujet (1), \zh{推出} verbe « lancer/mettre en place » (1), \zh{网上预约服务} objet composé correct (1), \zh{了} aspect accompli (1). \textit{Variante acceptée :} \zh{设立了...} / \zh{提供了...}.
    \item \textbf{Version :} \textbf{« Les citoyens n'ont qu'à appuyer sur leur téléphone (faire un clic) pour obtenir leurs documents, réalisant véritablement le "parcours des données multiple, parcours des masses réduit". »}
    \textit{Analyse :}
    \begin{itemize}
        \item \zh{市民只需} (Les citoyens n'ont qu'à / Il suffit aux citoyens de).
        \item \zh{在手机上点一点} (Litt. « sur téléphone taper un coup » $\to$ \textbf{faire une manipulation simple / quelques clics}).
        \item \zh{就能把证件办好} (Structure \zh{把} : Objet \zh{证件} antéposé + \zh{办好} résultat « bien traité/obtenu » $\to$ \textbf{obtenir leurs papiers / voir leur dossier traité}).
        \item \zh{真正实现了} (A véritablement réalisé/accompli).
        \item \zh{“数据多跑路，群众少跑腿”} (Slogan officiel : \textbf{« Que les données fassent le chemin, que les gens fassent moins de déplacements »} ou \textbf{« Zéro déplacement pour l'usager, traitement numérique des données »}).
    \end{itemize}
    \textit{Barème :} Sens global fidélité (2), Vocabulaire précis \zh{把/办好/实现} (2), Français naturel/admin (1).
\end{enumerate}
\end{exoresolu}

\begin{attention}[Les 5 erreurs qui coûtent des points au Bac — Leçon 19]
\begin{enumerate}
    \item \textbf{Confusion \zh{务} (wù, tâche/affaire) vs \zh{办} (bàn, traiter/verbe) vs \zh{业} (yè, métier/activité).}
    \textit{Piège :} Écrire \zh{*办务} au lieu de \zh{办事} (traiter affaires) ou \zh{服务} (service). \zh{办} est verbe, \zh{务} est nom. \zh{营业执照} (Licence) s'écrit avec \zh{业}, pas \zh{务}.
    \item \textbf{Mauvaise place du complément de lieu/moyen (\zh{在网上}, \zh{在窗口}) :} \textit{Erreur :} \zh{*我缴纳税费在网上。} \textbf{Correct :} \zh{我在网上缴纳税费。} (Sujet + \textbf{Lieu/Moyen} + Verbe + Objet). Règle d'or : \cle{Le lieu/le moyen précède toujours le verbe en chinois}.
    \item \textbf{Oubli ou mauvais usage de \zh{的} (de) dans les relatifs :} \textit{Erreur :} \zh{*我填写表格不完整。} (Phrase mal segmentée). \textbf{Correct :} \zh{我填写\textbf{的}表格不完整。} (Le formulaire \textbf{que} j'ai rempli). \textit{Règle :} Verbe/Adjectif + \zh{的} + Nom. Sans \zh{的}, c'est une phrase principale « Je remplis le formulaire [qui est] incomplet ».
    \item \textbf{Caractères « miroir » mal écrits : \zh{窗} (chuāng, guichet) vs \zh{囱} (cōng, cheminée) ; \zh{咨} (zī, consulter) vs \zh{资} (zī, ressources).} \textit{Perte :} 1 pt par caractère fautif en dictée/rédaction. \textbf{Astuce :} \zh{穴} (toit/caverne) = ouverture dans mur (\zh{窗, 窝, 空}) ; \zh{囗} (enfermement) = volume clos (\zh{囱, 园, 囚}). \zh{讠} (parole) = langage (\zh{咨, 询, 讯}) ; \zh{贝} (coquillage/monnaie) = valeur (\zh{资, 费, 贵}).
    \item \textbf{Calque du français « Pour + Infinitif » $\to$ *\zh{为 + Verbe} sans sujet logique :} \textit{Erreur :} \zh{*为了方便市民，网上预约系统开通了。} (Sujet « système » ne fait pas l'action « faciliter »). \textbf{Correct :} \zh{为方便市民，\textbf{市政府}开通了网上预约系统。} ou \zh{为了方便市民，\textbf{我们}开通了……}. \textit{Règle :} Le sujet de la principale doit être l'agent du but exprimé par \zh{为/为了}.
\end{enumerate}
\end{attention}

\newpage

\coursec{Exercices}

\courssub{Je vérifie mes connaissances}

\begin{exercice}[QCM : le bon choix]
Pour chaque question, entoure la bonne réponse.

1. Pour dire « délivrer une carte d'identité », on utilise le verbe :
\begin{itemize}
\item A. \zh{发} (fā)
\item B. \zh{办理} (bànlǐ)
\item C. \zh{发身份证} est impossible
\end{itemize}

2. « 如有问题，请联系我们 » signifie :
\begin{itemize}
\item A. « Si vous avez des questions, appelez la police. »
\item B. « Pour toute question, contactez-nous. »
\item C. « Nous avons un problème, merci de patienter. »
\end{itemize}

3. Le guichet numéro 2 se dit :
\begin{itemize}
\item A. \zh{二号窗口} (èr hào chuāngkǒu)
\item B. \zh{窗口二号} (chuāngkǒu èr hào)
\item C. \zh{两个窗口} (liǎng ge chuāngkǒu)
\end{itemize}

4. « Nous sommes ouverts du lundi au vendredi » se traduit par :
\begin{itemize}
\item A. \zh{我们星期一到星期五开门。}
\item B. \zh{我们从星期一到星期五开门。}
\item C. \zh{我们开门星期一到星期五。}
\end{itemize}
\end{exercice}

\begin{exercice}[Vrai ou faux ? Justifie ta réponse]
Dis si chaque affirmation est vraie ou fausse, puis justifie en une phrase en français.

1. \zh{服务中心} (fúwù zhōngxīn) signifie « centre de services ».
2. En chinois, la formule de politesse \zh{请} (qǐng) se place à la fin de la phrase.
3. \zh{办公时间} (bàngōng shíjiān) désigne les horaires d'ouverture au public.
4. On peut écrire \zh{请把材料交在二号窗口} sans aucune erreur de structure.
5. \zh{谢谢您的合作。} (Xièxie nín de hézuò.) est une formule de clôture adaptée à un avis administratif.
\end{exercice}

\begin{exercice}[Vocabulaire : complète le tableau]
Recopie et complète le tableau suivant avec la nature grammaticale et la traduction française.

\begin{center}
\begin{tabularx}{\linewidth}{|X|X|X|X|}
\hline
\rowcolor{popblueL} Caractères & Pinyin & Nature & Traduction \\
\hline
用户 & yònghù & \ldots & \ldots \\
\hline
材料 & cáiliào & \ldots & \ldots \\
\hline
申请 & shēnqǐng & \ldots & \ldots \\
\hline
免费 & miǎnfèi & \ldots & \ldots \\
\hline
复印 & fùyìn & \ldots & \ldots \\
\hline
工作人员 & gōngzuò rényuán & \ldots & \ldots \\
\hline
\end{tabularx}
\end{center}
\end{exercice}

\begin{exercice}[Associe caractères et pinyin]
Relie chaque caractère ou mot à son pinyin.

\begin{center}
\begin{tabularx}{\linewidth}{|X|X|}
\hline
\rowcolor{popblueL} Caractères & Pinyin (dans le désordre) \\
\hline
1. 窗口 & a. zhèngjiàn \\
\hline
2. 证件 & b. děnghòu \\
\hline
3. 排队 & c. chuāngkǒu \\
\hline
4. 等候 & d. liúyán \\
\hline
5. 留言 & e. páiduì \\
\hline
\end{tabularx}
\end{center}
\end{exercice}

\courssub{Je m'entraîne}

\begin{exercice}[Thème : traduis en chinois]
Traduis les cinq phrases suivantes en chinois (caractères et pinyin).

1. Notre centre délivre les cartes d'identité.
2. Vous devez déposer les documents au guichet numéro 2.
3. Nous sommes ouverts du lundi au vendredi.
4. Pour toute question, contactez-nous.
5. Merci de votre coopération.
\end{exercice}

\begin{exercice}[Texte à trous : les mots-outils]
Complète l'avis suivant avec les mots de la liste : \zh{请} — \zh{为} — \zh{从} — \zh{把} — \zh{的} — \zh{了}.

\begin{textebox}[Avis à compléter]
市民服务中心 \ldots\ldots 大家提供多种服务。\\
我们 \ldots\ldots 星期一 \ldots\ldots 星期五开门。\\
\ldots\ldots 您 \ldots\ldots 材料交到三号窗口。\\
办理完 \ldots\ldots ，您可以在等候区休息。\\
谢谢您的合作！
\end{textebox}
\end{exercice}

\begin{exercice}[Remets l'avis en ordre]
Les cinq blocs de cet avis de la mairie ont été mélangés. Remets-les dans le bon ordre (titre, objet, services, informations pratiques, clôture) et justifie ton choix en citant les connecteurs ou formules repérés.

\begin{itemize}
\item \textbf{A.} \zh{谢谢您的理解与配合！} (Xièxie nín de lǐjiě yǔ pèihé !)
\item \textbf{B.} \zh{雅温得市政府通知} (Yǎwēndé shì zhèngfǔ tōngzhī)
\item \textbf{C.} \zh{办公时间：星期一到星期五，上午八点到下午四点。} (Bàngōng shíjiān : xīngqīyī dào xīngqīwǔ, shàngwǔ bā diǎn dào xiàwǔ sì diǎn.)
\item \textbf{D.} \zh{为了更好地为市民服务，本中心从下周一起调整服务安排。} (Wèile gèng hǎo de wèi shìmín fúwù, běn zhōngxīn cóng xià zhōu yī qǐ tiáozhěng fúwù ānpái.)
\item \textbf{E.} \zh{我们提供以下服务：办理身份证、出生证明和居住证。} (Wǒmen tígōng yǐxià fúwù : bànlǐ shēnfènzhèng, chūshēng zhèngmíng hé jūzhùzhèng.)
\end{itemize}
\end{exercice}

\begin{exercice}[Corrige les erreurs]
Cet avis rédigé par un élève contient \textbf{cinq fautes} (ordre des mots, emploi de \zh{为}/\zh{办}, ponctuation, oubli de \zh{请}, caractère mal formé). Repère-les, souligne-les et écris la version corrigée.

\begin{textebox}[Avis de l'élève]
服务中心开门从星期一到星期五。\\
我们办大家提供很多服务。\\
您交材料到一号窗口。\\
如果有问题，联系我们打电话。\\
谢谢你的合作，
\end{textebox}
\end{exercice}

\courssub{Je me prépare au Bac}

\begin{exercice}[Compréhension de l'écrit : un avis du centre de services]
Lis attentivement l'avis suivant, puis réponds aux questions en français.

\begin{textebox}[杜阿拉市市民服务中心通知]
杜阿拉市市民服务中心通知\\
Dūālā shì shìmín fúwù zhōngxīn tōngzhī\\[4pt]
亲爱的市民朋友们：\\
Qīn'ài de shìmín péngyoumen :\\[4pt]
为了更好地为大家服务，本中心提供以下服务：办理身份证、出生证明、居住证，以及文件复印和翻译。所有表格都可以免费索取。\\
Wèile gèng hǎo de wèi dàjiā fúwù, běn zhōngxīn tígōng yǐxià fúwù : bànlǐ shēnfènzhèng, chūshēng zhèngmíng, jūzhùzhèng, yǐjí wénjiàn fùyìn hé fānyì. Suǒyǒu biǎogé dōu kěyǐ miǎnfèi suǒqǔ.\\[4pt]
办公时间：星期一到星期五，上午七点半到下午三点半。星期六上午只办理紧急业务。\\
Bàngōng shíjiān : xīngqīyī dào xīngqīwǔ, shàngwǔ qī diǎn bàn dào xiàwǔ sān diǎn bàn. Xīngqīliù shàngwǔ zhǐ bànlǐ jǐnjí yèwù.\\[4pt]
请把材料交到二号窗口，然后在等候区等候。\\
Qǐng bǎ cáiliào jiāo dào èr hào chuāngkǒu, ránhòu zài děnghòuqū děnghòu.\\[4pt]
如有问题，请拨打我们的电话，或在服务台留言。\\
Rú yǒu wèntí, qǐng bōdǎ wǒmen de diànhuà, huò zài fúwùtái liúyán.\\[4pt]
谢谢您的合作！\\
Xièxie nín de hézuò !
\end{textebox}

\textbf{Questions de compréhension (6 points)}
\begin{enumerate}
\item Quels sont les quatre services principaux mentionnés dans cet avis ?
\item Que peut-on obtenir gratuitement dans ce centre ?
\item Quels sont les horaires d'ouverture en semaine ? Et le samedi ?
\item Que doit faire l'usager après avoir déposé ses documents ?
\item Quelles sont les deux façons de poser une question au centre ?
\end{enumerate}

\textbf{Questions de langue (4 points)}
\begin{enumerate}
\item Relevez dans le texte la structure avec \zh{把} et recopiez-la avec son pinyin.
\item Expliquez en français le rôle de \zh{为了} en début de phrase.
\item Traduisez en français : \zh{如有问题，请拨打我们的电话。}
\item Trouvez dans le texte un synonyme de \zh{市民} au sens de « public, usagers ».
\end{enumerate}
\end{exercice}

\begin{exercice}[Thème et version]
\textbf{A. Thème (français → chinois) — 5 points.} Traduis en chinois (caractères et pinyin) :

1. Notre centre délivre les cartes d'identité et les actes de naissance.
2. Vous devez déposer les documents au guichet numéro 2.
3. Nous sommes ouverts du lundi au vendredi, de huit heures à seize heures.
4. Pour toute question, contactez-nous par téléphone.
5. Merci de votre coopération.

\textbf{B. Version (chinois → français) — 5 points.} Traduis en français le texte suivant :

\begin{textebox}[Version]
各位用户：本办公室为大家办理居住证和出生证明。\\
Gè wèi yònghù : běn bàngōngshì wèi dàjiā bànlǐ jūzhùzhèng hé chūshēng zhèngmíng.\\[4pt]
请先到服务台领取表格，填好后交到一号窗口。\\
Qǐng xiān dào fúwùtái lǐngqǔ biǎogé, tián hǎo hòu jiāo dào yī hào chuāngkǒu.\\[4pt]
办公时间：星期一到星期五，上午八点到下午三点。\\
Bàngōng shíjiān : xīngqīyī dào xīngqīwǔ, shàngwǔ bā diǎn dào xiàwǔ sān diǎn.\\[4pt]
如有问题，请询问工作人员。感谢您的配合！\\
Rú yǒu wèntí, qǐng xúnwèn gōngzuò rényuán. Gǎnxiè nín de pèihé !
\end{textebox}
\end{exercice}

\begin{exercice}[Expression écrite : rédige un avis administratif]
\textbf{Sujet (10 points).} Tu es employé(e) à la mairie de ta ville (Yaoundé, Douala, Bafoussam, Buea\ldots). Le maire te demande de rédiger en chinois un \textbf{avis aux usagers} affiché à l'entrée du centre de services.

\textbf{Consignes :}
\begin{itemize}
\item Longueur : \textbf{80 à 120 caractères chinois}.
\item Ton avis doit comporter les cinq parties étudiées : titre, formule d'appel, liste des services offerts, informations pratiques (horaires, guichet), formule de clôture polie.
\item Utilise au moins \textbf{deux connecteurs} parmi : \zh{为了}， \zh{首先}， \zh{然后}， \zh{以及}， \zh{如果}.
\item Emploie au moins une structure avec \zh{请} et une structure avec \zh{把}.
\item Écris les caractères avec soin : attention aux familles \zh{请}/\zh{清}/\zh{情} et \zh{问}/\zh{间}.
\end{itemize}

\textbf{Éléments d'évaluation :} respect des consignes et du nombre de caractères (2 pts) ; correction des structures grammaticales (3 pts) ; richesse et exactitude du vocabulaire administratif (3 pts) ; qualité de l'écriture des caractères (2 pts).
\end{exercice}

\newpage

\coursec{Corrigés des exercices}

\begin{corrige}[Exercice 1 — QCM : le bon choix]

\textbf{1. Réponse B.} Le verbe \zh{办理} (bànlǐ) signifie « effectuer une démarche, traiter (un dossier) » : c'est le verbe administratif par excellence. \zh{发} (fā) signifie « émettre, distribuer » ; on peut dire \zh{发身份证} (« délivrer une carte d'identité »), donc la réponse C est fausse, mais dans le registre des services administratifs, \zh{办理} est la réponse attendue.

\textbf{2. Réponse B.} \zh{如} = si (registre écrit), \zh{有问题} = avoir des questions, \zh{请} = s'il vous plaît, \zh{联系我们} = nous contacter. La réponse A est un faux sens (\zh{联系} ne veut pas dire « appeler la police »), la C inverse le sujet.

\textbf{3. Réponse A.} En chinois, le numéro précède le nom : \zh{二号窗口} (èr hào chuāngkǒu), littéralement « numéro 2 guichet ». La réponse B inverse l'ordre ; la C signifie « deux guichets » (avec le classificateur \zh{个}).

\textbf{4. Réponse B.} La structure correcte est \zh{从}A\zh{到}B (cóng A dào B) : \zh{我们从星期一到星期五开门。} (Wǒmen cóng xīngqīyī dào xīngqīwǔ kāimén.) En A, il manque \zh{从} ; en C, le complément de temps est mal placé (il doit se situer avant le verbe).
\end{corrige}

\begin{corrige}[Exercice 2 — Vrai ou faux ? Justifie ta réponse]

\textbf{1. Vrai.} \zh{服务} (fúwù) = service ; \zh{中心} (zhōngxīn) = centre. \zh{服务中心} signifie bien « centre de services ».

\textbf{2. Faux.} \zh{请} (qǐng) se place \emph{avant} le verbe, en début de phrase ou de proposition : \zh{请把材料交到二号窗口。} (Qǐng bǎ cáiliào jiāo dào èr hào chuāngkǒu.) « Veuillez déposer les documents au guichet n° 2. » Le placer en fin de phrase est une erreur typique des francophones qui calquent « s'il vous plaît ».

\textbf{3. Vrai.} \zh{办公} (bàngōng) = travailler (dans un bureau) ; \zh{时间} (shíjiān) = temps, horaires. \zh{办公时间} désigne les horaires d'ouverture au public d'une administration.

\textbf{4. Vrai.} La phrase \zh{请把材料交在二号窗口} est correcte : structure \zh{把} + objet + verbe \zh{交} + complément de lieu introduit par \zh{在}. (Variante tout aussi correcte : \zh{交到二号窗口} avec \zh{到} « vers ».)

\textbf{5. Vrai.} \zh{谢谢您的合作。} (Xièxie nín de hézuò.) « Merci de votre coopération » est une formule de clôture standard des avis administratifs ; \zh{您} (nín), forme polie de « vous », convient au registre officiel.
\end{corrige}

\begin{corrige}[Exercice 3 — Vocabulaire : complète le tableau]

\begin{center}
\begin{tabularx}{\linewidth}{|X|X|X|X|}
\hline
\rowcolor{popblueL} Caractères & Pinyin & Nature & Traduction \\
\hline
用户 & yònghù & nom & usager, utilisateur \\
\hline
材料 & cáiliào & nom & documents, pièces (d'un dossier) \\
\hline
申请 & shēnqǐng & verbe / nom & demander, faire une demande ; demande \\
\hline
免费 & miǎnfèi & adjectif / adverbe & gratuit ; gratuitement \\
\hline
复印 & fùyìn & verbe / nom & photocopier ; photocopie \\
\hline
工作人员 & gōngzuò rényuán & nom (locution) & personnel, agent(s) \\
\hline
\end{tabularx}
\end{center}

\textbf{Points de vigilance :} \zh{免费} peut fonctionner comme adverbe : \zh{免费索取} (miǎnfèi suǒqǔ) « obtenir gratuitement ». \zh{申请} est à la fois verbe (\zh{申请居住证} « demander un titre de séjour ») et nom (\zh{申请表} « formulaire de demande »).
\end{corrige}

\begin{corrige}[Exercice 4 — Associe caractères et pinyin]

\begin{center}
\begin{tabularx}{\linewidth}{|X|X|X|}
\hline
\rowcolor{popblueL} Caractères & Pinyin & Traduction \\
\hline
1. 窗口 & c. chuāngkǒu & guichet, fenêtre \\
\hline
2. 证件 & a. zhèngjiàn & document officiel, pièce \\
\hline
3. 排队 & e. páiduì & faire la queue \\
\hline
4. 等候 & b. děnghòu & attendre \\
\hline
5. 留言 & d. liúyán & laisser un message \\
\hline
\end{tabularx}
\end{center}

\textbf{Aide-mémoire :} \zh{窗} porte la clé du « trou » (\zh{穴}) avec \zh{囱} ; \zh{证} porte la clé de la parole (\zh{讠}) — un document « atteste » par la parole écrite ; \zh{排} porte la clé de la main (\zh{扌}) — on se « range » dans la file.
\end{corrige}

\begin{corrige}[Exercice 5 — Thème : traduis en chinois]

\textbf{1.} \zh{我们中心办理身份证。}\\
\emph{Wǒmen zhōngxīn bànlǐ shēnfènzhèng.}\\
(Variante acceptée : \zh{我们中心发放身份证。} avec \zh{发放} « délivrer ».)

\textbf{2.} \zh{您要把材料交到二号窗口。}\\
\emph{Nín yào bǎ cáiliào jiāo dào èr hào chuāngkǒu.}\\
Structure \zh{把} : \zh{把} + objet (\zh{材料}) + verbe (\zh{交}) + complément de direction (\zh{到} + lieu). Le numéro précède le nom : \zh{二号窗口}.

\textbf{3.} \zh{我们从星期一到星期五开门。}\\
\emph{Wǒmen cóng xīngqīyī dào xīngqīwǔ kāimén.}\\
Ne pas oublier \zh{从} ; le complément de temps se place avant le verbe \zh{开门}.

\textbf{4.} \zh{如有问题，请联系我们。}\\
\emph{Rú yǒu wèntí, qǐng liánxì wǒmen.}\\
Formule figée du registre administratif ; \zh{如} est la forme écrite de \zh{如果}.

\textbf{5.} \zh{谢谢您的合作。}\\
\emph{Xièxie nín de hézuò.}\\
Registre poli : \zh{您} et non \zh{你}.
\end{corrige}

\begin{corrige}[Exercice 6 — Texte à trous : les mots-outils]

Texte complété :

\begin{textebox}[Avis complété — corrigé]
市民服务中心\textbf{为}大家提供多种服务。\\
Shìmín fúwù zhōngxīn wèi dàjiā tígōng duō zhǒng fúwù.\\
我们\textbf{从}星期一\textbf{到}星期五开门。\\
Wǒmen cóng xīngqīyī dào xīngqīwǔ kāimén.\\
\textbf{请}您\textbf{把}材料交到三号窗口。\\
Qǐng nín bǎ cáiliào jiāo dào sān hào chuāngkǒu.\\
办理完\textbf{了}，您可以在等候区休息。\\
Bànlǐ wán le, nín kěyǐ zài děnghòuqū xiūxi.\\
谢谢您的合作！\\
Xièxie nín de hézuò !
\end{textebox}

Ordre des mots à insérer : \zh{为} — \zh{从} — \zh{到} — \zh{请} — \zh{把} — \zh{了}.

\textbf{Justifications :}
\begin{itemize}
\item \zh{为}大家提供 : structure \zh{为} + bénéficiaire + verbe (« fournir \emph{à} tout le monde »).
\item \zh{从}\ldots\zh{到}\ldots : cadre temporel « de\ldots à\ldots » (deux trous consécutifs).
\item \zh{请} en tête de la consigne ; \zh{把} devant l'objet \zh{材料} (l'objet est mis en avant avant le verbe \zh{交}).
\item \zh{办理完了} : \zh{了} marque l'accomplissement — « une fois la démarche terminée ».
\end{itemize}

\textbf{Remarque :} le mot \zh{的} de la liste n'est pas utilisé dans ce texte (distracteur) ; il apparaît dans \zh{谢谢您的合作} déjà imprimé.
\end{corrige}

\begin{corrige}[Exercice 7 — Remets l'avis en ordre]

\textbf{Ordre correct : B — D — E — C — A.}

\begin{enumerate}
\item \textbf{B.} \zh{雅温得市政府通知} : le \emph{titre} vient toujours en premier ; il annonce l'autorité émettrice (mairie de Yaoundé).
\item \textbf{D.} \zh{为了更好地为市民服务，本中心从下周一起调整服务安排。} : l'\emph{objet} de l'avis, introduit par le connecteur de but \zh{为了} (« afin de »), placé juste après le titre.
\item \textbf{E.} \zh{我们提供以下服务：办理身份证、出生证明和居住证。} : la \emph{liste des services}, signalée par la formule \zh{以下服务} (« les services suivants ») suivie des deux-points chinois.
\item \textbf{C.} \zh{办公时间：星期一到星期五，上午八点到下午四点。} : les \emph{informations pratiques} (horaires), introduites par l'étiquette \zh{办公时间}.
\item \textbf{A.} \zh{谢谢您的理解与配合！} : la \emph{formule de clôture} polie, toujours en dernière position.
\end{enumerate}

\textbf{Justification par les marqueurs :} \zh{为了} ouvre logiquement l'annonce ; \zh{以下} annonce une énumération ; \zh{谢谢您的\ldots} ne peut que clore le texte. Cet ordre (titre → objet → services → infos pratiques → clôture) est le plan type de tout avis administratif, en Chine comme au Cameroun.
\end{corrige}

\begin{corrige}[Exercice 8 — Corrige les erreurs]

\textbf{Les cinq fautes repérées :}

\begin{enumerate}
\item \zh{服务中心开门从星期一到星期五。} : \textbf{ordre des mots} fautif. Le cadre \zh{从}\ldots\zh{到}\ldots se place \emph{avant} le verbe.
\item \zh{我们办大家提供很多服务。} : \textbf{confusion} \zh{办}/\zh{为}. « Fournir \emph{à} quelqu'un » exige la préposition \zh{为} (wèi), non le verbe \zh{办} (bàn).
\item \zh{您交材料到一号窗口。} : \textbf{oubli de} \zh{请} (et idéalement de la structure \zh{把}) dans une consigne adressée à l'usager ; sans \zh{请}, le ton est sec, inadapté à un avis.
\item \zh{如果有问题，联系我们打电话。} : \textbf{structure double} fautive — deux verbes se télescopent. On dit \zh{请联系我们} \emph{ou} \zh{请给我们打电话}, pas les deux à la fois.
\item \zh{谢谢你的合作，} : \textbf{ponctuation} fautive (virgule au lieu du point final) et registre : \zh{您} est préférable à \zh{你} dans un avis officiel.
\end{enumerate}

\textbf{Version corrigée :}

\begin{textebox}[Avis corrigé]
服务中心从星期一到星期五开门。\\
Fúwù zhōngxīn cóng xīngqīyī dào xīngqīwǔ kāimén.\\
我们为大家提供很多服务。\\
Wǒmen wèi dàjiā tígōng hěn duō fúwù.\\
请您把材料交到一号窗口。\\
Qǐng nín bǎ cáiliào jiāo dào yī hào chuāngkǒu.\\
如果有问题，请联系我们。\\
Rúguǒ yǒu wèntí, qǐng liánxì wǒmen.\\
谢谢您的合作。\\
Xièxie nín de hézuò.
\end{textebox}
\end{corrige}

\begin{corrige}[Exercice 9 — Compréhension de l'écrit : un avis du centre de services]

\textbf{Barème indicatif :} compréhension 6 pts (1 + 1 + 2 + 1 + 1) ; langue 4 pts (1 + 1 + 1 + 1).

\textbf{Questions de compréhension}

\textbf{1.} Les quatre services principaux sont : la délivrance des cartes d'identité (\zh{办理身份证}), des actes de naissance (\zh{出生证明}), des titres de séjour (\zh{居住证}), ainsi que la photocopie et la traduction de documents (\zh{文件复印和翻译}). \hfill (1 pt : 0,25 par service)

\textbf{2.} On peut obtenir gratuitement tous les formulaires : \zh{所有表格都可以免费索取} (« tous les formulaires peuvent être demandés gratuitement »). \hfill (1 pt)

\textbf{3.} En semaine, le centre est ouvert du lundi au vendredi, de 7 h 30 à 15 h 30. Le samedi matin, il ne traite que les affaires urgentes (\zh{只办理紧急业务}). \hfill (2 pts)

\textbf{4.} Après avoir déposé ses documents au guichet n° 2, l'usager doit attendre dans la salle (la zone) d'attente : \zh{然后在等候区等候}. \hfill (1 pt)

\textbf{5.} On peut poser une question de deux façons : en téléphonant au centre (\zh{拨打我们的电话}) ou en laissant un message au bureau d'accueil (\zh{在服务台留言}). \hfill (1 pt)

\textbf{Questions de langue}

\textbf{1.} Structure avec \zh{把} : \zh{请把材料交到二号窗口} (Qǐng bǎ cáiliào jiāo dào èr hào chuāngkǒu). \hfill (1 pt)

\textbf{2.} \zh{为了} (wèile) introduit le \emph{but} de l'action : « afin de, pour ». Placé en tête de phrase, il annonce l'objectif poursuivi par l'administration (ici : mieux servir le public) avant d'énoncer la mesure prise. \hfill (1 pt)

\textbf{3.} \zh{如有问题，请拨打我们的电话。} : « Si vous avez des questions, veuillez composer notre numéro (nous appeler). » \hfill (1 pt)

\textbf{4.} Un synonyme de \zh{市民} au sens de « public, usagers » : \zh{大家} (dàjiā, « tout le monde ») dans \zh{为大家服务} ; on accepte aussi \zh{市民朋友们} (« chers concitoyens »). \hfill (1 pt)

\textbf{Points de vigilance :} ne pas traduire \zh{办公时间} par « temps de travail des employés » mais « horaires d'ouverture » ; \zh{紧急业务} = « affaires urgentes », non « business d'urgence » ; \zh{七点半} = 7 h 30 (litt. « sept heures et demie »).
\end{corrige}

\begin{corrige}[Exercice 10 — Thème et version]

\textbf{A. Thème (5 points — 1 pt par phrase)}

\textbf{1.} \zh{我们中心办理身份证和出生证明。}\\
\emph{Wǒmen zhōngxīn bànlǐ shēnfènzhèng hé chūshēng zhèngmíng.}

\textbf{2.} \zh{您要把材料交到二号窗口。}\\
\emph{Nín yào bǎ cáiliào jiāo dào èr hào chuāngkǒu.}

\textbf{3.} \zh{我们从星期一到星期五开门，从上午八点到下午四点。}\\
\emph{Wǒmen cóng xīngqīyī dào xīngqīwǔ kāimén, cóng shàngwǔ bā diǎn dào xiàwǔ sì diǎn.}\\
(Forme condensée acceptée : \zh{办公时间是星期一到星期五，上午八点到下午四点。})

\textbf{4.} \zh{如有问题，请打电话联系我们。}\\
\emph{Rú yǒu wèntí, qǐng dǎ diànhuà liánxì wǒmen.}\\
(Variante : \zh{请通过电话联系我们。})

\textbf{5.} \zh{谢谢您的合作。}\\
\emph{Xièxie nín de hézuò.}

\textbf{B. Version (5 points)}

« Chers usagers,\\
notre bureau traite pour vous les demandes de titres de séjour et d'actes de naissance.\\
Veuillez d'abord retirer un formulaire au bureau d'accueil, puis, une fois rempli, le déposer au guichet n° 1.\\
Horaires d'ouverture : du lundi au vendredi, de 8 heures à 15 heures.\\
Pour toute question, veuillez vous adresser au personnel. Merci de votre coopération ! »

\textbf{Barème indicatif (version) :} 1 pt par segment correctement traduit (appel, services, consigne en deux étapes, horaires, clôture).

\textbf{Points de vigilance :} \zh{各位用户} = « chers usagers » (\zh{位} est le classificateur poli des personnes) ; \zh{领取} = « retirer, recevoir (un document) » ; \zh{填好后} = « une fois rempli » (\zh{好} marque l'achèvement, \zh{后} = « après ») ; \zh{询问} = « s'enquérir auprès de », plus soutenu que \zh{问}.
\end{corrige}

\begin{corrige}[Exercice 11 — Expression écrite : rédige un avis administratif]

\textbf{Barème indicatif (10 pts) :} respect des consignes et du nombre de caractères (2 pts) ; correction grammaticale (3 pts) ; vocabulaire administratif riche et exact (3 pts) ; qualité de l'écriture des caractères (2 pts).

\textbf{Proposition de rédaction modèle :}

\begin{textebox}[Modèle d'avis — mairie de Douala]
杜阿拉市市民服务中心通知\\
Dūālā shì shìmín fúwù zhōngxīn tōngzhī\\[4pt]
亲爱的市民朋友们：\\
Qīn'ài de shìmín péngyoumen :\\[4pt]
为了更好地为大家服务，本中心提供以下服务：办理身份证、出生证明，以及文件复印。\\
Wèile gèng hǎo de wèi dàjiā fúwù, běn zhōngxīn tígōng yǐxià fúwù : bànlǐ shēnfènzhèng, chūshēng zhèngmíng, yǐjí wénjiàn fùyìn.\\[4pt]
请先到服务台领取表格，然后把材料交到二号窗口。\\
Qǐng xiān dào fúwùtái lǐngqǔ biǎogé, ránhòu bǎ cáiliào jiāo dào èr hào chuāngkǒu.\\[4pt]
办公时间：星期一到星期五，上午八点到下午四点。如果有问题，请询问工作人员。\\
Bàngōng shíjiān : xīngqīyī dào xīngqīwǔ, shàngwǔ bā diǎn dào xiàwǔ sì diǎn. Rúguǒ yǒu wèntí, qǐng xúnwèn gōngzuò rényuán.\\[4pt]
谢谢您的合作！\\
Xièxie nín de hézuò !
\end{textebox}

(Environ 110 caractères hors titre et pinyin : dans la fourchette demandée.)

\textbf{Vérification des consignes :}
\begin{itemize}
\item Cinq parties présentes : titre (\zh{\ldots 通知}), appel (\zh{亲爱的市民朋友们}), services (\zh{提供以下服务：\ldots}), informations pratiques (guichet, horaires), clôture (\zh{谢谢您的合作！}).
\item Connecteurs utilisés : \zh{为了}, \zh{以及}, \zh{然后}, \zh{如果} (quatre, minimum deux exigés).
\item Structure avec \zh{请} : \zh{请先到服务台领取表格} ; structure avec \zh{把} : \zh{把材料交到二号窗口}.
\end{itemize}

\textbf{Points de vigilance pour l'écriture des caractères :}
\begin{itemize}
\item \zh{请} (clé de la parole \zh{讠} + \zh{青}) : ne pas confondre avec \zh{清} (clé de l'eau \zh{氵}, « clair ») ni \zh{情} (clé du cœur \zh{忄}, « sentiment »). Même prononciation proche (qǐng / qīng / qíng), sens très différents.
\item \zh{问} (wèn, « demander » : la porte \zh{门} avec la bouche \zh{口} à l'intérieur) ne se confond pas avec \zh{间} (jiān, « entre, pièce » : la porte \zh{门} avec le soleil \zh{日}). Dans \zh{时间}, c'est bien \zh{间}.
\item Ordre des traits : pour \zh{窗}, la clé \zh{穴} (toit-trou) s'écrit d'abord, de haut en bas ; pour \zh{证} et \zh{请}, la clé \zh{讠} s'écrit avant la partie droite (gauche → droite).
\item Ponctuation pleine chasse : deux-points \zh{：} après la formule d'appel, point \zh{。} en fin de phrase, énumération avec \zh{、} entre les services.
\end{itemize}
\end{corrige}

\newpage

\coursec{Fiche bilan}

\begin{popfigure}
\begin{center}\tcbox[colback=popblue,colframe=popblue,coltext=white,arc=9pt,boxsep=4pt,left=6pt,right=6pt]{\bfseries\small Offre de services aux usagers (mairie)}\end{center}
\vspace{-2pt}
\begin{tcbraster}[raster columns=3,raster equal height=rows,raster column skip=6pt,raster row skip=6pt,enhanced,blankest]
\begin{tcolorbox}[enhanced,colback=white,colframe=poporange,boxrule=0.9pt,arc=3pt,title={Lexique des services},fonttitle=\bfseries\scriptsize,coltitle=white,colbacktitle=poporange,left=4pt,right=4pt,top=2pt,bottom=2pt,boxsep=1pt,toptitle=1.5pt,bottomtitle=1.5pt,halign=left]
\begin{itemize}[nosep,leftmargin=*,itemsep=1pt,font=\scriptsize]
\item 服务 fúwù, 窗口 chuāngkǒu, 证件 zhèngjiàn
\item 办理 bànlǐ, 申请 shēnqǐng, 免费 miǎnfèi
\end{itemize}
\end{tcolorbox}
\begin{tcolorbox}[enhanced,colback=white,colframe=popgreen,boxrule=0.9pt,arc=3pt,title={Structures clés},fonttitle=\bfseries\scriptsize,coltitle=white,colbacktitle=popgreen,left=4pt,right=4pt,top=2pt,bottom=2pt,boxsep=1pt,toptitle=1.5pt,bottomtitle=1.5pt,halign=left]
\begin{itemize}[nosep,leftmargin=*,itemsep=1pt,font=\scriptsize]
\item 我们为您提供…
\item 请带…到…
\end{itemize}
\end{tcolorbox}
\begin{tcolorbox}[enhanced,colback=white,colframe=poppurple,boxrule=0.9pt,arc=3pt,title={Connecteurs},fonttitle=\bfseries\scriptsize,coltitle=white,colbacktitle=poppurple,left=4pt,right=4pt,top=2pt,bottom=2pt,boxsep=1pt,toptitle=1.5pt,bottomtitle=1.5pt,halign=left]
\begin{itemize}[nosep,leftmargin=*,itemsep=1pt,font=\scriptsize]
\item 首先…其次…最后
\item 而且, 如果…就…
\end{itemize}
\end{tcolorbox}
\begin{tcolorbox}[enhanced,colback=white,colframe=poppink,boxrule=0.9pt,arc=3pt,title={Caractères du thème},fonttitle=\bfseries\scriptsize,coltitle=white,colbacktitle=poppink,left=4pt,right=4pt,top=2pt,bottom=2pt,boxsep=1pt,toptitle=1.5pt,bottomtitle=1.5pt,halign=left]
\begin{itemize}[nosep,leftmargin=*,itemsep=1pt,font=\scriptsize]
\item 办 ≠ 为 ; 证 ≠ 正
\item Radicaux 讠 辶 阝
\end{itemize}
\end{tcolorbox}
\begin{tcolorbox}[enhanced,colback=white,colframe=popgold,boxrule=0.9pt,arc=3pt,title={Traduction},fonttitle=\bfseries\scriptsize,coltitle=white,colbacktitle=popgold,left=4pt,right=4pt,top=2pt,bottom=2pt,boxsep=1pt,toptitle=1.5pt,bottomtitle=1.5pt,halign=left]
\begin{itemize}[nosep,leftmargin=*,itemsep=1pt,font=\scriptsize]
\item Thème : ordre S + V + C
\item Version : repérer le sujet
\end{itemize}
\end{tcolorbox}
\begin{tcolorbox}[enhanced,colback=white,colframe=popblue!60,boxrule=0.9pt,arc=3pt,title={Modèle d'avis},fonttitle=\bfseries\scriptsize,coltitle=white,colbacktitle=popblue!60,left=4pt,right=4pt,top=2pt,bottom=2pt,boxsep=1pt,toptitle=1.5pt,bottomtitle=1.5pt,halign=left]
\begin{itemize}[nosep,leftmargin=*,itemsep=1pt,font=\scriptsize]
\item Titre + services + horaires + contact
\end{itemize}
\end{tcolorbox}
\begin{tcolorbox}[enhanced,colback=white,colframe=popmuted,boxrule=0.9pt,arc=3pt,title={Méthode de rédaction},fonttitle=\bfseries\scriptsize,coltitle=white,colbacktitle=popmuted,left=4pt,right=4pt,top=2pt,bottom=2pt,boxsep=1pt,toptitle=1.5pt,bottomtitle=1.5pt,halign=left]
\begin{itemize}[nosep,leftmargin=*,itemsep=1pt,font=\scriptsize]
\item Saluer → offrir → expliquer → remercier
\end{itemize}
\end{tcolorbox}
\end{tcbraster}

\legende{Carte mentale de la leçon 19 : l'offre de services aux usagers des collectivités locales.}
\end{popfigure}

\begin{bilan}
\textbf{1. Lexique clé à connaître par cœur}
\begin{itemize}
  \item \zh{服务} fúwù : service ; \zh{为市民服务} wèi shìmín fúwù : servir les citoyens.
  \item \zh{办理} bànlǐ : effectuer (une démarche) ; \zh{申请} shēnqǐng : demander, déposer une demande.
  \item \zh{证件} zhèngjiàn : pièce, document officiel ; \zh{身份证} shēnfènzhèng : carte d'identité.
  \item \zh{窗口} chuāngkǒu : guichet ; \zh{办公室} bàngōngshì : bureau.
  \item \zh{免费} miǎnfèi : gratuit ; \zh{复印} fùyìn : photocopier ; \zh{表格} biǎogé : formulaire.
  \item \zh{工作时间} gōngzuò shíjiān : horaires d'ouverture ; \zh{市政府} shì zhèngfǔ : mairie.
\end{itemize}

\textbf{2. Règles de grammaire à connaître par cœur}
\begin{center}
\begin{tabularx}{\linewidth}{|X|X|X|}
\hline
\rowcolor{popblueL} Structure & Exemple & Traduction \\
\hline
Sujet + 为 + personne + 提供 + service & \zh{我们为您提供免费服务。} Wǒmen wèi nín tígōng miǎnfèi fúwù. & Nous vous offrons un service gratuit. \\
\hline
请 + verbe + complément & \zh{请带身份证到二号窗口。} Qǐng dài shēnfènzhèng dào èr hào chuāngkǒu. & Veuillez apporter votre carte d'identité au guichet n° 2. \\
\hline
如果 … 就 … & \zh{如果有问题，就请打电话。} Rúguǒ yǒu wèntí, jiù qǐng dǎ diànhuà. & Si vous avez un problème, appelez-nous. \\
\hline
可以 / 需要 + verbe & \zh{您可以在这里办理证件。} Nín kěyǐ zài zhèlǐ bànlǐ zhèngjiàn. & Vous pouvez faire vos documents ici. \\
\hline
先 … 然后 … 最后 … & \zh{先填表，然后交费，最后取证。} Xiān tián biǎo, ránhòu jiāo fèi, zuìhòu qǔ zhèng. & D'abord remplir le formulaire, puis payer, enfin retirer le document. \\
\hline
\end{tabularx}
\end{center}

\textbf{3. Écriture des caractères du thème}
\begin{itemize}
  \item \zh{办} (bàn) : 4 traits, ne pas confondre avec \zh{为} (wèi) ; \zh{证} (zhèng) : clé de la parole 讠 + \zh{正}.
  \item \zh{窗} (chuāng) : clé du toit 穴 en haut ; \zh{费} (fèi) : clé du coquillage 贝 (argent) en bas.
  \item Ordre des traits : toujours de haut en bas, de gauche à droite, l'horizontal avant le vertical.
\end{itemize}

\textbf{4. Méthodes express}
\begin{itemize}
  \item \textbf{Rédiger un avis de services :} titre → formule d'accueil (\zh{亲爱的市民朋友们}) → liste des services avec 提供/可以 → horaires et lieu → contact → remerciement (\zh{谢谢大家！}).
  \item \textbf{Thème (français → chinois) :} repérer sujet et verbe, placer les compléments de lieu et de temps AVANT le verbe, vérifier les tons et les classificateurs.
  \item \textbf{Version (chinois → français) :} identifier le sujet, traduire bloc par bloc, restituer les connecteurs (因为, 所以, 而且).
\end{itemize}
\end{bilan}

\begin{aretenir}[Checklist avant le Bac]
\begin{itemize}
  \item $\square$ Je connais le lexique des services : 服务, 办理, 申请, 证件, 窗口, 免费, 表格.
  \item $\square$ Je sais utiliser la structure \zh{为 + personne + 提供 + service}.
  \item $\square$ Je maîtrise les impératifs polis avec \zh{请} + verbe.
  \item $\square$ Je sais enchaîner les étapes avec \zh{首先 / 先 … 然后 … 最后 …}.
  \item $\square$ Je place correctement les compléments de temps et de lieu avant le verbe.
  \item $\square$ Je distingue les caractères proches : 办/为, 证/正, 费/废.
  \item $\square$ J'écris le pinyin avec les bons tons sur les voyelles.
  \item $\square$ Je connais le plan d'un avis : titre, accueil, services, horaires, contact, remerciement.
  \item $\square$ Je sais traduire une phrase simple dans les deux sens (thème et version).
  \item $\square$ J'utilise les connecteurs 而且, 如果…就…, 因为…所以… pour lier mes phrases.
  \item $\square$ Je relis ma production : tons, ordre des mots, ponctuation chinoise ，。
  \item $\square$ Je sais présenter les services d'une mairie camerounaise en chinois simple.
\end{itemize}
\end{aretenir}

\findecours

\end{document}
